% **************************************************************************************************************
% A Classic Thesis Style
% An Homage to The Elements of Typographic Style
%
% Copyright (C) 2015 André Miede http://www.miede.de
%
% If you like the style then I would appreciate a postcard. My address
% can be found in the file ClassicThesis.pdf. A collection of the
% postcards I received so far is available online at
% http://postcards.miede.de
%
% License:
% This program is free software; you can redistribute it and/or modify
% it under the terms of the GNU General Public License as published by
% the Free Software Foundation; either version 2 of the License, or
% (at your option) any later version.
%
% This program is distributed in the hope that it will be useful,
% but WITHOUT ANY WARRANTY; without even the implied warranty of
% MERCHANTABILITY or FITNESS FOR A PARTICULAR PURPOSE.  See the
% GNU General Public License for more details.
%
% You should have received a copy of the GNU General Public License
% along with this program; see the file COPYING.  If not, write to
% the Free Software Foundation, Inc., 59 Temple Place - Suite 330,
% Boston, MA 02111-1307, USA.
%
% **************************************************************************************************************
\RequirePackage{fix-cm} % fix some latex issues see: http://texdoc.net/texmf-dist/doc/latex/base/fixltx2e.pdf
\documentclass[ titlepage,numbers=noenddot,headinclude,oneside, %1headlines,% letterpaper a4paper
                footinclude=true,cleardoublepage=empty,abstractoff, % <--- obsolete, remove (todo)
                BCOR=5mm,paper=a4,fontsize=11pt,%13pt,a4paper,%
                ngerman,spanish,%
                ]{scrreprt}

%********************************************************************
% Note: Make all your adjustments in here
%*******************************************************
% ****************************************************************************************************
% classicthesis-config.tex
% formerly known as loadpackages.sty, classicthesis-ldpkg.sty, and classicthesis-preamble.sty
% Use it at the beginning of your ClassicThesis.tex, or as a LaTeX Preamble
% in your ClassicThesis.{tex,lyx} with % ****************************************************************************************************
% classicthesis-config.tex
% formerly known as loadpackages.sty, classicthesis-ldpkg.sty, and classicthesis-preamble.sty
% Use it at the beginning of your ClassicThesis.tex, or as a LaTeX Preamble
% in your ClassicThesis.{tex,lyx} with % ****************************************************************************************************
% classicthesis-config.tex
% formerly known as loadpackages.sty, classicthesis-ldpkg.sty, and classicthesis-preamble.sty
% Use it at the beginning of your ClassicThesis.tex, or as a LaTeX Preamble
% in your ClassicThesis.{tex,lyx} with \input{classicthesis-config}
% ****************************************************************************************************
% If you like the classicthesis, then I would appreciate a postcard.
% My address can be found in the file ClassicThesis.pdf. A collection
% of the postcards I received so far is available online at
% http://postcards.miede.de
% ****************************************************************************************************


% ****************************************************************************************************
% 0. Set the encoding of your files. UTF-8 is the only sensible encoding nowadays. If you can't read
% äöüßáéçèê∂åëæƒÏ€ then change the encoding setting in your editor, not the line below. If your editor
% does not support utf8 use another editor!
% ****************************************************************************************************
\PassOptionsToPackage{utf8}{inputenc}
	\usepackage{inputenc}

% ****************************************************************************************************
% 1. Configure classicthesis for your needs here, e.g., remove "drafting" below
% in order to deactivate the time-stamp on the pages
% ****************************************************************************************************
\PassOptionsToPackage{eulerchapternumbers,listings,%
					 pdfspacing,%floatperchapter,%linedheaders,%
					 subfig,beramono,eulermath,parts}{classicthesis}
% ********************************************************************
% Available options for classicthesis.sty
% (see ClassicThesis.pdf for more information):
% drafting
% parts nochapters linedheaders
% eulerchapternumbers beramono eulermath pdfspacing minionprospacing
% tocaligned dottedtoc manychapters
% listings floatperchapter subfig
% ********************************************************************


% ****************************************************************************************************
% 2. Personal data and user ad-hoc commands
% ****************************************************************************************************
\newcommand{\myTitle}{Introducción a la Homología Singular. Resolución del problema del clique máximo.\xspace}
%\newcommand{\mySubtitle}{An Homage to The Elements of Typographic Style\xspace}
%\newcommand{\myDegree}{Doktor-Ingenieur (Dr.-Ing.)\xspace}
\newcommand{\myName}{José Carlos Entrena Jiménez\xspace}
\newcommand{\myProf}{Francisco Urbano Pérez-Aranda\xspace}
\newcommand{\myOtherProf}{Salvador García López\xspace}
\newcommand{\mySupervisor}{Put name here\xspace}
\newcommand{\myFaculty}{Put data here\xspace}
\newcommand{\myDepartment}{Put data here\xspace}
\newcommand{\myUni}{Universidad de Granada\xspace}
\newcommand{\myLocation}{Granada\xspace}
\newcommand{\myTime}{Septiembre 2017\xspace}

% ********************************************************************
% Setup, finetuning, and useful commands
% ********************************************************************
\newcounter{dummy} % necessary for correct hyperlinks (to index, bib, etc.)
\newlength{\abcd} % for ab..z string length calculation
\providecommand{\mLyX}{L\kern-.1667em\lower.25em\hbox{Y}\kern-.125emX\@}
\newcommand{\ie}{i.\,e.}
\newcommand{\Ie}{I.\,e.}
\newcommand{\eg}{e.\,g.}
\newcommand{\Eg}{E.\,g.}
% ****************************************************************************************************


% ****************************************************************************************************
% 3. Loading some handy packages
% ****************************************************************************************************
% ********************************************************************
% Packages with options that might require adjustments
% ********************************************************************
%\PassOptionsToPackage{ngerman,american}{babel}   % change this to your language(s)
% Spanish languages need extra options in order to work with this template
\PassOptionsToPackage{spanish,es-lcroman}{babel}
	\usepackage{babel}

\usepackage{csquotes}
\PassOptionsToPackage{%
    %backend=biber, %instead of bibtex
	backend=bibtex8,bibencoding=ascii,%
	language=spanish,%
	style=numeric-comp,%
    %style=authoryear-comp, % Author 1999, 2010
    %bibstyle=authoryear,dashed=false, % dashed: substitute rep. author with ---
    sorting=nyt, % name, year, title
    maxbibnames=10, % default: 3, et al.
    %backref=true,%
    natbib=true % natbib compatibility mode (\citep and \citet still work)
}{biblatex}
    \usepackage{biblatex}

\PassOptionsToPackage{fleqn}{amsmath}       % math environments and more by the AMS
    \usepackage{amsmath}

% ********************************************************************
% General useful packages
% ********************************************************************
\PassOptionsToPackage{T1}{fontenc} % T2A for cyrillics
    \usepackage{fontenc}
\usepackage{textcomp} % fix warning with missing font shapes
\usepackage{scrhack} % fix warnings when using KOMA with listings package
\usepackage{xspace} % to get the spacing after macros right
\usepackage{mparhack} % get marginpar right
\usepackage{fixltx2e} % fixes some LaTeX stuff --> since 2015 in the LaTeX kernel (see below)
%\usepackage[latest]{latexrelease} % will be used once available in more distributions (ISSUE #107)
\PassOptionsToPackage{printonlyused,smaller}{acronym}
    \usepackage{acronym} % nice macros for handling all acronyms in the thesis
    %\renewcommand{\bflabel}[1]{{#1}\hfill} % fix the list of acronyms --> no longer working
    %\renewcommand*{\acsfont}[1]{\textsc{#1}}
    \renewcommand*{\aclabelfont}[1]{\acsfont{#1}}
% ****************************************************************************************************


% ****************************************************************************************************
% 4. Setup floats: tables, (sub)figures, and captions
% ****************************************************************************************************
\usepackage{tabularx} % better tables
    \setlength{\extrarowheight}{3pt} % increase table row height
\newcommand{\tableheadline}[1]{\multicolumn{1}{c}{\spacedlowsmallcaps{#1}}}
\newcommand{\myfloatalign}{\centering} % to be used with each float for alignment
\usepackage{caption}
% Thanks to cgnieder and Claus Lahiri
% http://tex.stackexchange.com/questions/69349/spacedlowsmallcaps-in-caption-label
% [REMOVED DUE TO OTHER PROBLEMS, SEE ISSUE #82]
%\DeclareCaptionLabelFormat{smallcaps}{\bothIfFirst{#1}{~}\MakeTextLowercase{\textsc{#2}}}
%\captionsetup{font=small,labelformat=smallcaps} % format=hang,
\captionsetup{font=small} % format=hang,
\usepackage{subfig}
% ****************************************************************************************************


% ****************************************************************************************************
% 5. Setup code listings
% ****************************************************************************************************
\usepackage{listings}
%\lstset{emph={trueIndex,root},emphstyle=\color{BlueViolet}}%\underbar} % for special keywords
\lstset{language=[LaTeX]Tex,%C++,
    morekeywords={PassOptionsToPackage,selectlanguage},
    keywordstyle=\color{RoyalBlue},%\bfseries,
    basicstyle=\small\ttfamily,
    %identifierstyle=\color{NavyBlue},
    commentstyle=\color{Green}\ttfamily,
    stringstyle=\rmfamily,
    numbers=none,%left,%
    numberstyle=\scriptsize,%\tiny
    stepnumber=5,
    numbersep=8pt,
    showstringspaces=false,
    breaklines=true,
    %frameround=ftff,
    %frame=single,
    belowcaptionskip=.75\baselineskip
    %frame=L
}
% ****************************************************************************************************


% ****************************************************************************************************
% 6. PDFLaTeX, hyperreferences and citation backreferences
% ****************************************************************************************************
% ********************************************************************
% Using PDFLaTeX
% ********************************************************************
\PassOptionsToPackage{pdftex,hyperfootnotes=false,pdfpagelabels}{hyperref}
    \usepackage{hyperref}  % backref linktocpage pagebackref
\pdfcompresslevel=9
\pdfadjustspacing=1
\PassOptionsToPackage{pdftex}{graphicx}
    \usepackage{graphicx}


% ********************************************************************
% Hyperreferences
% ********************************************************************
\hypersetup{%
    %draft, % = no hyperlinking at all (useful in b/w printouts)
    colorlinks=true, linktocpage=true, pdfstartpage=3, pdfstartview=FitV,%
    % uncomment the following line if you want to have black links (e.g., for printing)
    %colorlinks=false, linktocpage=false, pdfstartpage=3, pdfstartview=FitV, pdfborder={0 0 0},%
    breaklinks=true, pdfpagemode=UseNone, pageanchor=true, pdfpagemode=UseOutlines,%
    plainpages=false, bookmarksnumbered, bookmarksopen=true, bookmarksopenlevel=1,%
    hypertexnames=true, pdfhighlight=/O,%nesting=true,%frenchlinks,%
    urlcolor=webbrown, linkcolor=RoyalBlue, citecolor=webgreen, %pagecolor=RoyalBlue,%
    %urlcolor=Black, linkcolor=Black, citecolor=Black, %pagecolor=Black,%
    pdftitle={\myTitle},%
    pdfauthor={\textcopyright\ \myName, \myUni, \myFaculty},%
    pdfsubject={},%
    pdfkeywords={},%
    pdfcreator={pdfLaTeX},%
    pdfproducer={LaTeX with hyperref and classicthesis}%
}

% ********************************************************************
% Setup autoreferences
% ********************************************************************
% There are some issues regarding autorefnames
% http://www.ureader.de/msg/136221647.aspx
% http://www.tex.ac.uk/cgi-bin/texfaq2html?label=latexwords
% you have to redefine the makros for the
% language you use, e.g., american, ngerman
% (as chosen when loading babel/AtBeginDocument)
% ********************************************************************
\makeatletter
\@ifpackageloaded{babel}%
    {%
       \addto\extrasamerican{%
			\renewcommand*{\figureautorefname}{Figure}%
			\renewcommand*{\tableautorefname}{Table}%
			\renewcommand*{\partautorefname}{Part}%
			\renewcommand*{\chapterautorefname}{Chapter}%
			\renewcommand*{\sectionautorefname}{Section}%
			\renewcommand*{\subsectionautorefname}{Section}%
			\renewcommand*{\subsubsectionautorefname}{Section}%
                }%
       \addto\extrasngerman{%
			\renewcommand*{\paragraphautorefname}{Absatz}%
			\renewcommand*{\subparagraphautorefname}{Unterabsatz}%
			\renewcommand*{\footnoteautorefname}{Fu\"snote}%
			\renewcommand*{\FancyVerbLineautorefname}{Zeile}%
			\renewcommand*{\theoremautorefname}{Theorem}%
			\renewcommand*{\appendixautorefname}{Anhang}%
			\renewcommand*{\equationautorefname}{Gleichung}%
			\renewcommand*{\itemautorefname}{Punkt}%
                }%
            % Fix to getting autorefs for subfigures right (thanks to Belinda Vogt for changing the definition)
            \providecommand{\subfigureautorefname}{\figureautorefname}%
    }{\relax}
\makeatother


% ****************************************************************************************************
% 7. Last calls before the bar closes
% ****************************************************************************************************
% ********************************************************************
% Development Stuff
% ********************************************************************
\listfiles
%\PassOptionsToPackage{l2tabu,orthodox,abort}{nag}
%   \usepackage{nag}
%\PassOptionsToPackage{warning, all}{onlyamsmath}
%   \usepackage{onlyamsmath}

% ********************************************************************
% Last, but not least...
% ********************************************************************
\usepackage{classicthesis}
% ****************************************************************************************************


% ****************************************************************************************************
% 8. Further adjustments (experimental)
% ****************************************************************************************************
% ********************************************************************
% Changing the text area
% ********************************************************************
%\linespread{1.05} % a bit more for Palatino
%\areaset[current]{312pt}{761pt} % 686 (factor 2.2) + 33 head + 42 head \the\footskip
%\setlength{\marginparwidth}{7em}%
%\setlength{\marginparsep}{2em}%

% ********************************************************************
% Using different fonts
% ********************************************************************
%\usepackage[oldstylenums]{kpfonts} % oldstyle notextcomp
%\usepackage[osf]{libertine}
%\usepackage[light,condensed,math]{iwona}
%\renewcommand{\sfdefault}{iwona}
%\usepackage{lmodern} % <-- no osf support :-(
%\usepackage{cfr-lm} %
%\usepackage[urw-garamond]{mathdesign} <-- no osf support :-(
%\usepackage[default,osfigures]{opensans} % scale=0.95
%\usepackage[sfdefault]{FiraSans}
% ****************************************************************************************************


% ****************************************************************************************************
% 9. Custom settings
% ****************************************************************************************************

% Espacios en el documento:
\linespread{1.01}                        % Espacio entre líneas.
\setlength\parindent{0pt}               % Selecciona la indentación para cada inicio de párrafo.
\usepackage[parfill]{parskip}						% Línea en blanco después de párrafo

% Color
\usepackage{xcolor}

% Creative commons
\usepackage[scale=1.3]{ccicons}


\usepackage{chngcntr}
\counterwithin{figure}{chapter}

\usepackage{amsthm}

\newtheoremstyle{break}
  {1em} % measure of space to leave above the theorem.
	{\topsep} % measure of space to leave below the theorem.
  {\itshape} % name of font to use in the body of the theorem
	{} % measure of space to indent
  {\bfseries} % name of head font
	{} % punctuation between head and body
  {\newline} % space after theorem head; " " = normal interword space
	{} %

\theoremstyle{break}
\newtheorem{theorem}{Teorema}[chapter]
\newtheorem{corollary}{Corolario}[chapter]
\newtheorem{proposition}{Proposición}[chapter]
\newtheorem{lemma}{Lema}[chapter]
\newtheorem{definition}{Definición}[chapter]

\theoremstyle{remark}
\newtheorem{remark}{Nota}[chapter]
\newtheorem{example}{Ejemplo}[chapter]


\def\vs{\vspace{5mm}}

\makeatletter
\let\c@proposition\c@theorem
\let\c@corollary\c@theorem
\let\c@lemma\c@theorem
\let\c@definition\c@theorem
\let\c@example\c@theorem
\let\c@remark\c@theorem
\makeatother

\def\M{{\mathcal M}}
\def\N{{\mathcal N}}
\def\C{{\mathcal C}}
\def\Khat{\widehat{K}}
\def\phat{\widehat{p}}
\def\that{T}

\usepackage{amssymb}
% \newenvironment{proof}{\noindent{\textbf{Proof.}}}{\hfill$\blacksquare$\vs}

% \newenvironment{remark}{\stepcounter{theorem} \noindent {\textbf{Remark
% \thetheorem}.}
% }{}

% Custom commands - Added by Alejandro García
\newcommand{\R}{\mathbb{R}}
\newcommand{\Z}{\mathbb{Z}}
% \newcommand{\H}{\mathbb{H}}

\def\Q{{\mathbb Q}}
\def\N{{\mathbb N}}
% \def\C{{\mathbb C}}
\def\S{{\mathbb S}}
\def\L{{\mathbb L}}
\def\H{{\mathbb H}}
\def\K{{\mathbb K}}
\def\X{{\mathbb X}}
\def\Y{{\mathbb Y}}
% \def\Z{{\mathbb Z}}
\def\A{{\mathbb A}}
\def\J{{\mathbb J}}
\def\I{{\mathbb I}}
\def\T{{\mathbb T}}


\makeatletter
\newcommand*{\defeq}{\mathrel{\rlap{%
                     \raisebox{0.3ex}{$\m@th\cdot$}}%
                     \raisebox{-0.3ex}{$\m@th\cdot$}}%
                     =}
\makeatother

\usepackage{tikz}
\usepackage{tikz-cd}
\tikzcdset{every label/.append style = {font = \normalsize}}

\usetikzlibrary{matrix,shapes,arrows}

%\usepackage{todonotes}

% To do notes, taken from http://tex.stackexchange.com/a/178806
\usepackage{xargs}                      % Use more than one optional parameter in a new commands
\usepackage[colorinlistoftodos,prependcaption,textsize=small]{todonotes}
\newcommandx{\fixme}[2][1=]{{\let\marginpar\oldmarginpar \todo[#1]{#2}}}
\newcommandx{\unsure}[2][1=]{\fixme[linecolor=red,backgroundcolor=red!25,bordercolor=red,#1]{#2}}
\newcommandx{\change}[2][1=]{\fixme[linecolor=blue,backgroundcolor=blue!25,bordercolor=blue,#1]{#2}}
\newcommandx{\info}[2][1=]{\fixme[linecolor=OliveGreen,backgroundcolor=OliveGreen!25,bordercolor=OliveGreen,#1]{#2}}
\newcommandx{\improvement}[2][1=]{\fixme[linecolor=Plum,backgroundcolor=Plum!25,bordercolor=Plum,#1]{#2}}

% Autoref commands
\newcommand*{\corollaryautorefname}{Corolario}
\newcommand*{\propositionautorefname}{Proposición}
\newcommand*{\lemmaautorefname}{Lema}
\newcommand*{\definitionautorefname}{Definición}
\newcommand*{\exampleautorefname}{Ejemplo}
\newcommand*{\remarkautorefname}{Nota}
\newcommand*{\algorithmautorefname}{Algoritmo}

% Numbering of equations following chapter.number convention
\numberwithin{equation}{chapter}

% Black filled squares are much better than empty white ones
% That is YOUR opinion -JC-
%\renewcommand{\qedsymbol}{$\blacksquare$}

% Partial derivatives
\newcommand*{\pd}[3][]{\ensuremath{\frac{\partial^{#1} #2}{\partial #3}}}

\newcommandx{\Ric}{{\operatorname{Ric}}}

% Defines xmapsto, xarrows
\usepackage{mathtools}

\usepackage[nice]{nicefrac}

% Defines a pointing arrow to a symbol with a note below.
% From http://tex.stackexchange.com/a/117408
\newcommand{\eqnote}[2]{%
	\underset{\substack{\uparrow\\\mathrlap{\text{\hspace{-1em}#2}}}}{#1}}

% Algorithms
\usepackage{algorithm}
\usepackage[noend]{algpseudocode}

% ****************************************************************************************************
% If you like the classicthesis, then I would appreciate a postcard.
% My address can be found in the file ClassicThesis.pdf. A collection
% of the postcards I received so far is available online at
% http://postcards.miede.de
% ****************************************************************************************************


% ****************************************************************************************************
% 0. Set the encoding of your files. UTF-8 is the only sensible encoding nowadays. If you can't read
% äöüßáéçèê∂åëæƒÏ€ then change the encoding setting in your editor, not the line below. If your editor
% does not support utf8 use another editor!
% ****************************************************************************************************
\PassOptionsToPackage{utf8}{inputenc}
	\usepackage{inputenc}

% ****************************************************************************************************
% 1. Configure classicthesis for your needs here, e.g., remove "drafting" below
% in order to deactivate the time-stamp on the pages
% ****************************************************************************************************
\PassOptionsToPackage{eulerchapternumbers,listings,%
					 pdfspacing,%floatperchapter,%linedheaders,%
					 subfig,beramono,eulermath,parts}{classicthesis}
% ********************************************************************
% Available options for classicthesis.sty
% (see ClassicThesis.pdf for more information):
% drafting
% parts nochapters linedheaders
% eulerchapternumbers beramono eulermath pdfspacing minionprospacing
% tocaligned dottedtoc manychapters
% listings floatperchapter subfig
% ********************************************************************


% ****************************************************************************************************
% 2. Personal data and user ad-hoc commands
% ****************************************************************************************************
\newcommand{\myTitle}{Introducción a la Homología Singular. Resolución del problema del clique máximo.\xspace}
%\newcommand{\mySubtitle}{An Homage to The Elements of Typographic Style\xspace}
%\newcommand{\myDegree}{Doktor-Ingenieur (Dr.-Ing.)\xspace}
\newcommand{\myName}{José Carlos Entrena Jiménez\xspace}
\newcommand{\myProf}{Francisco Urbano Pérez-Aranda\xspace}
\newcommand{\myOtherProf}{Salvador García López\xspace}
\newcommand{\mySupervisor}{Put name here\xspace}
\newcommand{\myFaculty}{Put data here\xspace}
\newcommand{\myDepartment}{Put data here\xspace}
\newcommand{\myUni}{Universidad de Granada\xspace}
\newcommand{\myLocation}{Granada\xspace}
\newcommand{\myTime}{Septiembre 2017\xspace}

% ********************************************************************
% Setup, finetuning, and useful commands
% ********************************************************************
\newcounter{dummy} % necessary for correct hyperlinks (to index, bib, etc.)
\newlength{\abcd} % for ab..z string length calculation
\providecommand{\mLyX}{L\kern-.1667em\lower.25em\hbox{Y}\kern-.125emX\@}
\newcommand{\ie}{i.\,e.}
\newcommand{\Ie}{I.\,e.}
\newcommand{\eg}{e.\,g.}
\newcommand{\Eg}{E.\,g.}
% ****************************************************************************************************


% ****************************************************************************************************
% 3. Loading some handy packages
% ****************************************************************************************************
% ********************************************************************
% Packages with options that might require adjustments
% ********************************************************************
%\PassOptionsToPackage{ngerman,american}{babel}   % change this to your language(s)
% Spanish languages need extra options in order to work with this template
\PassOptionsToPackage{spanish,es-lcroman}{babel}
	\usepackage{babel}

\usepackage{csquotes}
\PassOptionsToPackage{%
    %backend=biber, %instead of bibtex
	backend=bibtex8,bibencoding=ascii,%
	language=spanish,%
	style=numeric-comp,%
    %style=authoryear-comp, % Author 1999, 2010
    %bibstyle=authoryear,dashed=false, % dashed: substitute rep. author with ---
    sorting=nyt, % name, year, title
    maxbibnames=10, % default: 3, et al.
    %backref=true,%
    natbib=true % natbib compatibility mode (\citep and \citet still work)
}{biblatex}
    \usepackage{biblatex}

\PassOptionsToPackage{fleqn}{amsmath}       % math environments and more by the AMS
    \usepackage{amsmath}

% ********************************************************************
% General useful packages
% ********************************************************************
\PassOptionsToPackage{T1}{fontenc} % T2A for cyrillics
    \usepackage{fontenc}
\usepackage{textcomp} % fix warning with missing font shapes
\usepackage{scrhack} % fix warnings when using KOMA with listings package
\usepackage{xspace} % to get the spacing after macros right
\usepackage{mparhack} % get marginpar right
\usepackage{fixltx2e} % fixes some LaTeX stuff --> since 2015 in the LaTeX kernel (see below)
%\usepackage[latest]{latexrelease} % will be used once available in more distributions (ISSUE #107)
\PassOptionsToPackage{printonlyused,smaller}{acronym}
    \usepackage{acronym} % nice macros for handling all acronyms in the thesis
    %\renewcommand{\bflabel}[1]{{#1}\hfill} % fix the list of acronyms --> no longer working
    %\renewcommand*{\acsfont}[1]{\textsc{#1}}
    \renewcommand*{\aclabelfont}[1]{\acsfont{#1}}
% ****************************************************************************************************


% ****************************************************************************************************
% 4. Setup floats: tables, (sub)figures, and captions
% ****************************************************************************************************
\usepackage{tabularx} % better tables
    \setlength{\extrarowheight}{3pt} % increase table row height
\newcommand{\tableheadline}[1]{\multicolumn{1}{c}{\spacedlowsmallcaps{#1}}}
\newcommand{\myfloatalign}{\centering} % to be used with each float for alignment
\usepackage{caption}
% Thanks to cgnieder and Claus Lahiri
% http://tex.stackexchange.com/questions/69349/spacedlowsmallcaps-in-caption-label
% [REMOVED DUE TO OTHER PROBLEMS, SEE ISSUE #82]
%\DeclareCaptionLabelFormat{smallcaps}{\bothIfFirst{#1}{~}\MakeTextLowercase{\textsc{#2}}}
%\captionsetup{font=small,labelformat=smallcaps} % format=hang,
\captionsetup{font=small} % format=hang,
\usepackage{subfig}
% ****************************************************************************************************


% ****************************************************************************************************
% 5. Setup code listings
% ****************************************************************************************************
\usepackage{listings}
%\lstset{emph={trueIndex,root},emphstyle=\color{BlueViolet}}%\underbar} % for special keywords
\lstset{language=[LaTeX]Tex,%C++,
    morekeywords={PassOptionsToPackage,selectlanguage},
    keywordstyle=\color{RoyalBlue},%\bfseries,
    basicstyle=\small\ttfamily,
    %identifierstyle=\color{NavyBlue},
    commentstyle=\color{Green}\ttfamily,
    stringstyle=\rmfamily,
    numbers=none,%left,%
    numberstyle=\scriptsize,%\tiny
    stepnumber=5,
    numbersep=8pt,
    showstringspaces=false,
    breaklines=true,
    %frameround=ftff,
    %frame=single,
    belowcaptionskip=.75\baselineskip
    %frame=L
}
% ****************************************************************************************************


% ****************************************************************************************************
% 6. PDFLaTeX, hyperreferences and citation backreferences
% ****************************************************************************************************
% ********************************************************************
% Using PDFLaTeX
% ********************************************************************
\PassOptionsToPackage{pdftex,hyperfootnotes=false,pdfpagelabels}{hyperref}
    \usepackage{hyperref}  % backref linktocpage pagebackref
\pdfcompresslevel=9
\pdfadjustspacing=1
\PassOptionsToPackage{pdftex}{graphicx}
    \usepackage{graphicx}


% ********************************************************************
% Hyperreferences
% ********************************************************************
\hypersetup{%
    %draft, % = no hyperlinking at all (useful in b/w printouts)
    colorlinks=true, linktocpage=true, pdfstartpage=3, pdfstartview=FitV,%
    % uncomment the following line if you want to have black links (e.g., for printing)
    %colorlinks=false, linktocpage=false, pdfstartpage=3, pdfstartview=FitV, pdfborder={0 0 0},%
    breaklinks=true, pdfpagemode=UseNone, pageanchor=true, pdfpagemode=UseOutlines,%
    plainpages=false, bookmarksnumbered, bookmarksopen=true, bookmarksopenlevel=1,%
    hypertexnames=true, pdfhighlight=/O,%nesting=true,%frenchlinks,%
    urlcolor=webbrown, linkcolor=RoyalBlue, citecolor=webgreen, %pagecolor=RoyalBlue,%
    %urlcolor=Black, linkcolor=Black, citecolor=Black, %pagecolor=Black,%
    pdftitle={\myTitle},%
    pdfauthor={\textcopyright\ \myName, \myUni, \myFaculty},%
    pdfsubject={},%
    pdfkeywords={},%
    pdfcreator={pdfLaTeX},%
    pdfproducer={LaTeX with hyperref and classicthesis}%
}

% ********************************************************************
% Setup autoreferences
% ********************************************************************
% There are some issues regarding autorefnames
% http://www.ureader.de/msg/136221647.aspx
% http://www.tex.ac.uk/cgi-bin/texfaq2html?label=latexwords
% you have to redefine the makros for the
% language you use, e.g., american, ngerman
% (as chosen when loading babel/AtBeginDocument)
% ********************************************************************
\makeatletter
\@ifpackageloaded{babel}%
    {%
       \addto\extrasamerican{%
			\renewcommand*{\figureautorefname}{Figure}%
			\renewcommand*{\tableautorefname}{Table}%
			\renewcommand*{\partautorefname}{Part}%
			\renewcommand*{\chapterautorefname}{Chapter}%
			\renewcommand*{\sectionautorefname}{Section}%
			\renewcommand*{\subsectionautorefname}{Section}%
			\renewcommand*{\subsubsectionautorefname}{Section}%
                }%
       \addto\extrasngerman{%
			\renewcommand*{\paragraphautorefname}{Absatz}%
			\renewcommand*{\subparagraphautorefname}{Unterabsatz}%
			\renewcommand*{\footnoteautorefname}{Fu\"snote}%
			\renewcommand*{\FancyVerbLineautorefname}{Zeile}%
			\renewcommand*{\theoremautorefname}{Theorem}%
			\renewcommand*{\appendixautorefname}{Anhang}%
			\renewcommand*{\equationautorefname}{Gleichung}%
			\renewcommand*{\itemautorefname}{Punkt}%
                }%
            % Fix to getting autorefs for subfigures right (thanks to Belinda Vogt for changing the definition)
            \providecommand{\subfigureautorefname}{\figureautorefname}%
    }{\relax}
\makeatother


% ****************************************************************************************************
% 7. Last calls before the bar closes
% ****************************************************************************************************
% ********************************************************************
% Development Stuff
% ********************************************************************
\listfiles
%\PassOptionsToPackage{l2tabu,orthodox,abort}{nag}
%   \usepackage{nag}
%\PassOptionsToPackage{warning, all}{onlyamsmath}
%   \usepackage{onlyamsmath}

% ********************************************************************
% Last, but not least...
% ********************************************************************
\usepackage{classicthesis}
% ****************************************************************************************************


% ****************************************************************************************************
% 8. Further adjustments (experimental)
% ****************************************************************************************************
% ********************************************************************
% Changing the text area
% ********************************************************************
%\linespread{1.05} % a bit more for Palatino
%\areaset[current]{312pt}{761pt} % 686 (factor 2.2) + 33 head + 42 head \the\footskip
%\setlength{\marginparwidth}{7em}%
%\setlength{\marginparsep}{2em}%

% ********************************************************************
% Using different fonts
% ********************************************************************
%\usepackage[oldstylenums]{kpfonts} % oldstyle notextcomp
%\usepackage[osf]{libertine}
%\usepackage[light,condensed,math]{iwona}
%\renewcommand{\sfdefault}{iwona}
%\usepackage{lmodern} % <-- no osf support :-(
%\usepackage{cfr-lm} %
%\usepackage[urw-garamond]{mathdesign} <-- no osf support :-(
%\usepackage[default,osfigures]{opensans} % scale=0.95
%\usepackage[sfdefault]{FiraSans}
% ****************************************************************************************************


% ****************************************************************************************************
% 9. Custom settings
% ****************************************************************************************************

% Espacios en el documento:
\linespread{1.01}                        % Espacio entre líneas.
\setlength\parindent{0pt}               % Selecciona la indentación para cada inicio de párrafo.
\usepackage[parfill]{parskip}						% Línea en blanco después de párrafo

% Color
\usepackage{xcolor}

% Creative commons
\usepackage[scale=1.3]{ccicons}


\usepackage{chngcntr}
\counterwithin{figure}{chapter}

\usepackage{amsthm}

\newtheoremstyle{break}
  {1em} % measure of space to leave above the theorem.
	{\topsep} % measure of space to leave below the theorem.
  {\itshape} % name of font to use in the body of the theorem
	{} % measure of space to indent
  {\bfseries} % name of head font
	{} % punctuation between head and body
  {\newline} % space after theorem head; " " = normal interword space
	{} %

\theoremstyle{break}
\newtheorem{theorem}{Teorema}[chapter]
\newtheorem{corollary}{Corolario}[chapter]
\newtheorem{proposition}{Proposición}[chapter]
\newtheorem{lemma}{Lema}[chapter]
\newtheorem{definition}{Definición}[chapter]

\theoremstyle{remark}
\newtheorem{remark}{Nota}[chapter]
\newtheorem{example}{Ejemplo}[chapter]


\def\vs{\vspace{5mm}}

\makeatletter
\let\c@proposition\c@theorem
\let\c@corollary\c@theorem
\let\c@lemma\c@theorem
\let\c@definition\c@theorem
\let\c@example\c@theorem
\let\c@remark\c@theorem
\makeatother

\def\M{{\mathcal M}}
\def\N{{\mathcal N}}
\def\C{{\mathcal C}}
\def\Khat{\widehat{K}}
\def\phat{\widehat{p}}
\def\that{T}

\usepackage{amssymb}
% \newenvironment{proof}{\noindent{\textbf{Proof.}}}{\hfill$\blacksquare$\vs}

% \newenvironment{remark}{\stepcounter{theorem} \noindent {\textbf{Remark
% \thetheorem}.}
% }{}

% Custom commands - Added by Alejandro García
\newcommand{\R}{\mathbb{R}}
\newcommand{\Z}{\mathbb{Z}}
% \newcommand{\H}{\mathbb{H}}

\def\Q{{\mathbb Q}}
\def\N{{\mathbb N}}
% \def\C{{\mathbb C}}
\def\S{{\mathbb S}}
\def\L{{\mathbb L}}
\def\H{{\mathbb H}}
\def\K{{\mathbb K}}
\def\X{{\mathbb X}}
\def\Y{{\mathbb Y}}
% \def\Z{{\mathbb Z}}
\def\A{{\mathbb A}}
\def\J{{\mathbb J}}
\def\I{{\mathbb I}}
\def\T{{\mathbb T}}


\makeatletter
\newcommand*{\defeq}{\mathrel{\rlap{%
                     \raisebox{0.3ex}{$\m@th\cdot$}}%
                     \raisebox{-0.3ex}{$\m@th\cdot$}}%
                     =}
\makeatother

\usepackage{tikz}
\usepackage{tikz-cd}
\tikzcdset{every label/.append style = {font = \normalsize}}

\usetikzlibrary{matrix,shapes,arrows}

%\usepackage{todonotes}

% To do notes, taken from http://tex.stackexchange.com/a/178806
\usepackage{xargs}                      % Use more than one optional parameter in a new commands
\usepackage[colorinlistoftodos,prependcaption,textsize=small]{todonotes}
\newcommandx{\fixme}[2][1=]{{\let\marginpar\oldmarginpar \todo[#1]{#2}}}
\newcommandx{\unsure}[2][1=]{\fixme[linecolor=red,backgroundcolor=red!25,bordercolor=red,#1]{#2}}
\newcommandx{\change}[2][1=]{\fixme[linecolor=blue,backgroundcolor=blue!25,bordercolor=blue,#1]{#2}}
\newcommandx{\info}[2][1=]{\fixme[linecolor=OliveGreen,backgroundcolor=OliveGreen!25,bordercolor=OliveGreen,#1]{#2}}
\newcommandx{\improvement}[2][1=]{\fixme[linecolor=Plum,backgroundcolor=Plum!25,bordercolor=Plum,#1]{#2}}

% Autoref commands
\newcommand*{\corollaryautorefname}{Corolario}
\newcommand*{\propositionautorefname}{Proposición}
\newcommand*{\lemmaautorefname}{Lema}
\newcommand*{\definitionautorefname}{Definición}
\newcommand*{\exampleautorefname}{Ejemplo}
\newcommand*{\remarkautorefname}{Nota}
\newcommand*{\algorithmautorefname}{Algoritmo}

% Numbering of equations following chapter.number convention
\numberwithin{equation}{chapter}

% Black filled squares are much better than empty white ones
% That is YOUR opinion -JC-
%\renewcommand{\qedsymbol}{$\blacksquare$}

% Partial derivatives
\newcommand*{\pd}[3][]{\ensuremath{\frac{\partial^{#1} #2}{\partial #3}}}

\newcommandx{\Ric}{{\operatorname{Ric}}}

% Defines xmapsto, xarrows
\usepackage{mathtools}

\usepackage[nice]{nicefrac}

% Defines a pointing arrow to a symbol with a note below.
% From http://tex.stackexchange.com/a/117408
\newcommand{\eqnote}[2]{%
	\underset{\substack{\uparrow\\\mathrlap{\text{\hspace{-1em}#2}}}}{#1}}

% Algorithms
\usepackage{algorithm}
\usepackage[noend]{algpseudocode}

% ****************************************************************************************************
% If you like the classicthesis, then I would appreciate a postcard.
% My address can be found in the file ClassicThesis.pdf. A collection
% of the postcards I received so far is available online at
% http://postcards.miede.de
% ****************************************************************************************************


% ****************************************************************************************************
% 0. Set the encoding of your files. UTF-8 is the only sensible encoding nowadays. If you can't read
% äöüßáéçèê∂åëæƒÏ€ then change the encoding setting in your editor, not the line below. If your editor
% does not support utf8 use another editor!
% ****************************************************************************************************
\PassOptionsToPackage{utf8}{inputenc}
	\usepackage{inputenc}

% ****************************************************************************************************
% 1. Configure classicthesis for your needs here, e.g., remove "drafting" below
% in order to deactivate the time-stamp on the pages
% ****************************************************************************************************
\PassOptionsToPackage{eulerchapternumbers,listings,%
					 pdfspacing,%floatperchapter,%linedheaders,%
					 subfig,beramono,eulermath,parts}{classicthesis}
% ********************************************************************
% Available options for classicthesis.sty
% (see ClassicThesis.pdf for more information):
% drafting
% parts nochapters linedheaders
% eulerchapternumbers beramono eulermath pdfspacing minionprospacing
% tocaligned dottedtoc manychapters
% listings floatperchapter subfig
% ********************************************************************


% ****************************************************************************************************
% 2. Personal data and user ad-hoc commands
% ****************************************************************************************************
\newcommand{\myTitle}{Introducción a la Homología Singular. Resolución del problema del clique máximo.\xspace}
%\newcommand{\mySubtitle}{An Homage to The Elements of Typographic Style\xspace}
%\newcommand{\myDegree}{Doktor-Ingenieur (Dr.-Ing.)\xspace}
\newcommand{\myName}{José Carlos Entrena Jiménez\xspace}
\newcommand{\myProf}{Francisco Urbano Pérez-Aranda\xspace}
\newcommand{\myOtherProf}{Salvador García López\xspace}
\newcommand{\mySupervisor}{Put name here\xspace}
\newcommand{\myFaculty}{Put data here\xspace}
\newcommand{\myDepartment}{Put data here\xspace}
\newcommand{\myUni}{Universidad de Granada\xspace}
\newcommand{\myLocation}{Granada\xspace}
\newcommand{\myTime}{Septiembre 2017\xspace}

% ********************************************************************
% Setup, finetuning, and useful commands
% ********************************************************************
\newcounter{dummy} % necessary for correct hyperlinks (to index, bib, etc.)
\newlength{\abcd} % for ab..z string length calculation
\providecommand{\mLyX}{L\kern-.1667em\lower.25em\hbox{Y}\kern-.125emX\@}
\newcommand{\ie}{i.\,e.}
\newcommand{\Ie}{I.\,e.}
\newcommand{\eg}{e.\,g.}
\newcommand{\Eg}{E.\,g.}
% ****************************************************************************************************


% ****************************************************************************************************
% 3. Loading some handy packages
% ****************************************************************************************************
% ********************************************************************
% Packages with options that might require adjustments
% ********************************************************************
%\PassOptionsToPackage{ngerman,american}{babel}   % change this to your language(s)
% Spanish languages need extra options in order to work with this template
\PassOptionsToPackage{spanish,es-lcroman}{babel}
	\usepackage{babel}

\usepackage{csquotes}
\PassOptionsToPackage{%
    %backend=biber, %instead of bibtex
	backend=bibtex8,bibencoding=ascii,%
	language=spanish,%
	style=numeric-comp,%
    %style=authoryear-comp, % Author 1999, 2010
    %bibstyle=authoryear,dashed=false, % dashed: substitute rep. author with ---
    sorting=nyt, % name, year, title
    maxbibnames=10, % default: 3, et al.
    %backref=true,%
    natbib=true % natbib compatibility mode (\citep and \citet still work)
}{biblatex}
    \usepackage{biblatex}

\PassOptionsToPackage{fleqn}{amsmath}       % math environments and more by the AMS
    \usepackage{amsmath}

% ********************************************************************
% General useful packages
% ********************************************************************
\PassOptionsToPackage{T1}{fontenc} % T2A for cyrillics
    \usepackage{fontenc}
\usepackage{textcomp} % fix warning with missing font shapes
\usepackage{scrhack} % fix warnings when using KOMA with listings package
\usepackage{xspace} % to get the spacing after macros right
\usepackage{mparhack} % get marginpar right
\usepackage{fixltx2e} % fixes some LaTeX stuff --> since 2015 in the LaTeX kernel (see below)
%\usepackage[latest]{latexrelease} % will be used once available in more distributions (ISSUE #107)
\PassOptionsToPackage{printonlyused,smaller}{acronym}
    \usepackage{acronym} % nice macros for handling all acronyms in the thesis
    %\renewcommand{\bflabel}[1]{{#1}\hfill} % fix the list of acronyms --> no longer working
    %\renewcommand*{\acsfont}[1]{\textsc{#1}}
    \renewcommand*{\aclabelfont}[1]{\acsfont{#1}}
% ****************************************************************************************************


% ****************************************************************************************************
% 4. Setup floats: tables, (sub)figures, and captions
% ****************************************************************************************************
\usepackage{tabularx} % better tables
    \setlength{\extrarowheight}{3pt} % increase table row height
\newcommand{\tableheadline}[1]{\multicolumn{1}{c}{\spacedlowsmallcaps{#1}}}
\newcommand{\myfloatalign}{\centering} % to be used with each float for alignment
\usepackage{caption}
% Thanks to cgnieder and Claus Lahiri
% http://tex.stackexchange.com/questions/69349/spacedlowsmallcaps-in-caption-label
% [REMOVED DUE TO OTHER PROBLEMS, SEE ISSUE #82]
%\DeclareCaptionLabelFormat{smallcaps}{\bothIfFirst{#1}{~}\MakeTextLowercase{\textsc{#2}}}
%\captionsetup{font=small,labelformat=smallcaps} % format=hang,
\captionsetup{font=small} % format=hang,
\usepackage{subfig}
% ****************************************************************************************************


% ****************************************************************************************************
% 5. Setup code listings
% ****************************************************************************************************
\usepackage{listings}
%\lstset{emph={trueIndex,root},emphstyle=\color{BlueViolet}}%\underbar} % for special keywords
\lstset{language=[LaTeX]Tex,%C++,
    morekeywords={PassOptionsToPackage,selectlanguage},
    keywordstyle=\color{RoyalBlue},%\bfseries,
    basicstyle=\small\ttfamily,
    %identifierstyle=\color{NavyBlue},
    commentstyle=\color{Green}\ttfamily,
    stringstyle=\rmfamily,
    numbers=none,%left,%
    numberstyle=\scriptsize,%\tiny
    stepnumber=5,
    numbersep=8pt,
    showstringspaces=false,
    breaklines=true,
    %frameround=ftff,
    %frame=single,
    belowcaptionskip=.75\baselineskip
    %frame=L
}
% ****************************************************************************************************


% ****************************************************************************************************
% 6. PDFLaTeX, hyperreferences and citation backreferences
% ****************************************************************************************************
% ********************************************************************
% Using PDFLaTeX
% ********************************************************************
\PassOptionsToPackage{pdftex,hyperfootnotes=false,pdfpagelabels}{hyperref}
    \usepackage{hyperref}  % backref linktocpage pagebackref
\pdfcompresslevel=9
\pdfadjustspacing=1
\PassOptionsToPackage{pdftex}{graphicx}
    \usepackage{graphicx}


% ********************************************************************
% Hyperreferences
% ********************************************************************
\hypersetup{%
    %draft, % = no hyperlinking at all (useful in b/w printouts)
    colorlinks=true, linktocpage=true, pdfstartpage=3, pdfstartview=FitV,%
    % uncomment the following line if you want to have black links (e.g., for printing)
    %colorlinks=false, linktocpage=false, pdfstartpage=3, pdfstartview=FitV, pdfborder={0 0 0},%
    breaklinks=true, pdfpagemode=UseNone, pageanchor=true, pdfpagemode=UseOutlines,%
    plainpages=false, bookmarksnumbered, bookmarksopen=true, bookmarksopenlevel=1,%
    hypertexnames=true, pdfhighlight=/O,%nesting=true,%frenchlinks,%
    urlcolor=webbrown, linkcolor=RoyalBlue, citecolor=webgreen, %pagecolor=RoyalBlue,%
    %urlcolor=Black, linkcolor=Black, citecolor=Black, %pagecolor=Black,%
    pdftitle={\myTitle},%
    pdfauthor={\textcopyright\ \myName, \myUni, \myFaculty},%
    pdfsubject={},%
    pdfkeywords={},%
    pdfcreator={pdfLaTeX},%
    pdfproducer={LaTeX with hyperref and classicthesis}%
}

% ********************************************************************
% Setup autoreferences
% ********************************************************************
% There are some issues regarding autorefnames
% http://www.ureader.de/msg/136221647.aspx
% http://www.tex.ac.uk/cgi-bin/texfaq2html?label=latexwords
% you have to redefine the makros for the
% language you use, e.g., american, ngerman
% (as chosen when loading babel/AtBeginDocument)
% ********************************************************************
\makeatletter
\@ifpackageloaded{babel}%
    {%
       \addto\extrasamerican{%
			\renewcommand*{\figureautorefname}{Figure}%
			\renewcommand*{\tableautorefname}{Table}%
			\renewcommand*{\partautorefname}{Part}%
			\renewcommand*{\chapterautorefname}{Chapter}%
			\renewcommand*{\sectionautorefname}{Section}%
			\renewcommand*{\subsectionautorefname}{Section}%
			\renewcommand*{\subsubsectionautorefname}{Section}%
                }%
       \addto\extrasngerman{%
			\renewcommand*{\paragraphautorefname}{Absatz}%
			\renewcommand*{\subparagraphautorefname}{Unterabsatz}%
			\renewcommand*{\footnoteautorefname}{Fu\"snote}%
			\renewcommand*{\FancyVerbLineautorefname}{Zeile}%
			\renewcommand*{\theoremautorefname}{Theorem}%
			\renewcommand*{\appendixautorefname}{Anhang}%
			\renewcommand*{\equationautorefname}{Gleichung}%
			\renewcommand*{\itemautorefname}{Punkt}%
                }%
            % Fix to getting autorefs for subfigures right (thanks to Belinda Vogt for changing the definition)
            \providecommand{\subfigureautorefname}{\figureautorefname}%
    }{\relax}
\makeatother


% ****************************************************************************************************
% 7. Last calls before the bar closes
% ****************************************************************************************************
% ********************************************************************
% Development Stuff
% ********************************************************************
\listfiles
%\PassOptionsToPackage{l2tabu,orthodox,abort}{nag}
%   \usepackage{nag}
%\PassOptionsToPackage{warning, all}{onlyamsmath}
%   \usepackage{onlyamsmath}

% ********************************************************************
% Last, but not least...
% ********************************************************************
\usepackage{classicthesis}
% ****************************************************************************************************


% ****************************************************************************************************
% 8. Further adjustments (experimental)
% ****************************************************************************************************
% ********************************************************************
% Changing the text area
% ********************************************************************
%\linespread{1.05} % a bit more for Palatino
%\areaset[current]{312pt}{761pt} % 686 (factor 2.2) + 33 head + 42 head \the\footskip
%\setlength{\marginparwidth}{7em}%
%\setlength{\marginparsep}{2em}%

% ********************************************************************
% Using different fonts
% ********************************************************************
%\usepackage[oldstylenums]{kpfonts} % oldstyle notextcomp
%\usepackage[osf]{libertine}
%\usepackage[light,condensed,math]{iwona}
%\renewcommand{\sfdefault}{iwona}
%\usepackage{lmodern} % <-- no osf support :-(
%\usepackage{cfr-lm} %
%\usepackage[urw-garamond]{mathdesign} <-- no osf support :-(
%\usepackage[default,osfigures]{opensans} % scale=0.95
%\usepackage[sfdefault]{FiraSans}
% ****************************************************************************************************


% ****************************************************************************************************
% 9. Custom settings
% ****************************************************************************************************

% Espacios en el documento:
\linespread{1.01}                        % Espacio entre líneas.
\setlength\parindent{0pt}               % Selecciona la indentación para cada inicio de párrafo.
\usepackage[parfill]{parskip}						% Línea en blanco después de párrafo

% Color
\usepackage{xcolor}

% Creative commons
\usepackage[scale=1.3]{ccicons}


\usepackage{chngcntr}
\counterwithin{figure}{chapter}

\usepackage{amsthm}

\newtheoremstyle{break}
  {1em} % measure of space to leave above the theorem.
	{\topsep} % measure of space to leave below the theorem.
  {\itshape} % name of font to use in the body of the theorem
	{} % measure of space to indent
  {\bfseries} % name of head font
	{} % punctuation between head and body
  {\newline} % space after theorem head; " " = normal interword space
	{} %

\theoremstyle{break}
\newtheorem{theorem}{Teorema}[chapter]
\newtheorem{corollary}{Corolario}[chapter]
\newtheorem{proposition}{Proposición}[chapter]
\newtheorem{lemma}{Lema}[chapter]
\newtheorem{definition}{Definición}[chapter]

\theoremstyle{remark}
\newtheorem{remark}{Nota}[chapter]
\newtheorem{example}{Ejemplo}[chapter]


\def\vs{\vspace{5mm}}

\makeatletter
\let\c@proposition\c@theorem
\let\c@corollary\c@theorem
\let\c@lemma\c@theorem
\let\c@definition\c@theorem
\let\c@example\c@theorem
\let\c@remark\c@theorem
\makeatother

\def\M{{\mathcal M}}
\def\N{{\mathcal N}}
\def\C{{\mathcal C}}
\def\Khat{\widehat{K}}
\def\phat{\widehat{p}}
\def\that{T}

\usepackage{amssymb}
% \newenvironment{proof}{\noindent{\textbf{Proof.}}}{\hfill$\blacksquare$\vs}

% \newenvironment{remark}{\stepcounter{theorem} \noindent {\textbf{Remark
% \thetheorem}.}
% }{}

% Custom commands - Added by Alejandro García
\newcommand{\R}{\mathbb{R}}
\newcommand{\Z}{\mathbb{Z}}
% \newcommand{\H}{\mathbb{H}}

\def\Q{{\mathbb Q}}
\def\N{{\mathbb N}}
% \def\C{{\mathbb C}}
\def\S{{\mathbb S}}
\def\L{{\mathbb L}}
\def\H{{\mathbb H}}
\def\K{{\mathbb K}}
\def\X{{\mathbb X}}
\def\Y{{\mathbb Y}}
% \def\Z{{\mathbb Z}}
\def\A{{\mathbb A}}
\def\J{{\mathbb J}}
\def\I{{\mathbb I}}
\def\T{{\mathbb T}}


\makeatletter
\newcommand*{\defeq}{\mathrel{\rlap{%
                     \raisebox{0.3ex}{$\m@th\cdot$}}%
                     \raisebox{-0.3ex}{$\m@th\cdot$}}%
                     =}
\makeatother

\usepackage{tikz}
\usepackage{tikz-cd}
\tikzcdset{every label/.append style = {font = \normalsize}}

\usetikzlibrary{matrix,shapes,arrows}

%\usepackage{todonotes}

% To do notes, taken from http://tex.stackexchange.com/a/178806
\usepackage{xargs}                      % Use more than one optional parameter in a new commands
\usepackage[colorinlistoftodos,prependcaption,textsize=small]{todonotes}
\newcommandx{\fixme}[2][1=]{{\let\marginpar\oldmarginpar \todo[#1]{#2}}}
\newcommandx{\unsure}[2][1=]{\fixme[linecolor=red,backgroundcolor=red!25,bordercolor=red,#1]{#2}}
\newcommandx{\change}[2][1=]{\fixme[linecolor=blue,backgroundcolor=blue!25,bordercolor=blue,#1]{#2}}
\newcommandx{\info}[2][1=]{\fixme[linecolor=OliveGreen,backgroundcolor=OliveGreen!25,bordercolor=OliveGreen,#1]{#2}}
\newcommandx{\improvement}[2][1=]{\fixme[linecolor=Plum,backgroundcolor=Plum!25,bordercolor=Plum,#1]{#2}}

% Autoref commands
\newcommand*{\corollaryautorefname}{Corolario}
\newcommand*{\propositionautorefname}{Proposición}
\newcommand*{\lemmaautorefname}{Lema}
\newcommand*{\definitionautorefname}{Definición}
\newcommand*{\exampleautorefname}{Ejemplo}
\newcommand*{\remarkautorefname}{Nota}
\newcommand*{\algorithmautorefname}{Algoritmo}

% Numbering of equations following chapter.number convention
\numberwithin{equation}{chapter}

% Black filled squares are much better than empty white ones
% That is YOUR opinion -JC-
%\renewcommand{\qedsymbol}{$\blacksquare$}

% Partial derivatives
\newcommand*{\pd}[3][]{\ensuremath{\frac{\partial^{#1} #2}{\partial #3}}}

\newcommandx{\Ric}{{\operatorname{Ric}}}

% Defines xmapsto, xarrows
\usepackage{mathtools}

\usepackage[nice]{nicefrac}

% Defines a pointing arrow to a symbol with a note below.
% From http://tex.stackexchange.com/a/117408
\newcommand{\eqnote}[2]{%
	\underset{\substack{\uparrow\\\mathrlap{\text{\hspace{-1em}#2}}}}{#1}}

% Algorithms
\usepackage{algorithm}
\usepackage[noend]{algpseudocode}


%********************************************************************
% Bibliographies
%*******************************************************
\addbibresource{Bibliography.bib}

%********************************************************************
% Hyphenation
%*******************************************************
%\hyphenation{put special hyphenation here}

% ********************************************************************
% GO!GO!GO! MOVE IT!
%*******************************************************
\begin{document}
\frenchspacing
\raggedbottom
\selectlanguage{spanish} % Spanish
%\renewcommand*{\bibname}{new name}
%\setbibpreamble{}
\pagenumbering{roman}
\pagestyle{plain}
%********************************************************************
% Frontmatter
%*******************************************************
%\include{FrontBackmatter/DirtyTitlepage}
%*******************************************************
% Titlepage
%*******************************************************
\begin{titlepage}
    % if you want the titlepage to be centered, uncomment and fine-tune the line below (KOMA classes environment)
    % \begin{addmargin}[-1cm]{-3cm}
    \begin{center}
        \large

        \hfill

        \hfill

        \begingroup
            \color{Maroon}\spacedallcaps{\myTitle} \\ \bigskip
        \endgroup

        \spacedlowsmallcaps{\myName}

        \vfill

        \includegraphics[width=6cm]{gfx/UGRlogo.png} \\ \medskip

        %\mySubtitle \\ \medskip
        %\myDegree \\
        %\myDepartment \\
        %\myFaculty \\
        %\myUni \\ \bigskip

        \myTime\

        \vfill

        \myProf, \myOtherProf \\
        %\myFaculty \\
        \myUni \\ \bigskip


    \end{center}
%  \end{addmargin}

\end{titlepage}

\thispagestyle{empty}

\hfill

\vfill

\textbf{Licencia} \ccbysa

Este trabajo está bajo una licencia \hyperlink{https://creativecommons.org/licenses/by-sa/4.0/legalcode}{Creative Commons Attribution - ShareAlike 4.0 International License}.

Con esta licencia, puedes:
\begin{itemize}
	\item  Compartir, copiar, y redistribuir el contenido.
	\item  Modificar el material de cualquier forma, incluso con propósito comercial.
\end{itemize}

Con las siguientes obligaciones:
\begin{itemize}
	\item[\ccAttribution] Dar el crédito apropiado, aportando un enlace a la licencia, e indicar los cambios realizados. 
	\item[\ccShareAlike] Cualquier copia, o modificación debe ser distribuida con la misma licencia que el original.
	\item[] No aplicar medidas de carácter legal o tecnológico que impidan a otros hacer aquello que la licencia permita.
\end{itemize}
\clearpage

%\cleardoublepage\include{FrontBackmatter/Dedication}
%\cleardoublepage\include{FrontBackmatter/Foreword}
\cleardoublepage%*******************************************************
% Abstract
%*******************************************************
\begingroup
\let\clearpage\relax
\let\cleardoublepage\relax
\let\cleardoublepage\relax

\pdfbookmark[1]{Abstract}{Abstract}
\chapter*{Abstract}

In this work, we define the basics of singular homology theory, starting with the necessary geometric constructions
and the definition of the singular homology groups. Following the definitions and a proof of homotopy invariance,
we prove the baricentric subdivision theorem, the central result of this work, that leads us inmediatly to the
Mayer-Vietoris sequence. This and other results are applied to obtain the homology groups of the spheres and a number
of classical theorems, including the Brower fixed-point theorem and the Jordan-Brower separation theorem.

On the second part, we introduce the maximum clique problem, a well studied combinatorial optimization problem, known
to be NP-hard. Due to this, the use of metaheuristics is necessary to solve it in reasonable time. Several heuristics
are considered, explained in detail and applied to DIMACS instances. We discuss their behaviour and compare them,
showing which of the proposed algorithms can yield satisfactory results.


\paragraph{keywords} Homology theory, Homotopy invariance, Baricentric subdivision theorem, Mayer-Vietoris sequence,
Maximum clique, Combinatorial optimization problem, Metaheuristics, NP-hard, DIMACS.

\vfill

\endgroup

\vfill

%\cleardoublepage\include{FrontBackmatter/Publications}
%\cleardoublepage\include{FrontBackmatter/Acknowledgments}
\pagestyle{scrheadings}
\cleardoublepage%*******************************************************
% Table of Contents
%*******************************************************
%\phantomsection
\refstepcounter{dummy}
\pdfbookmark[1]{\contentsname}{tableofcontents}
\setcounter{tocdepth}{2} % <-- 2 includes up to subsections in the ToC
\setcounter{secnumdepth}{3} % <-- 3 numbers up to subsubsections
\manualmark
\markboth{\spacedlowsmallcaps{\contentsname}}{\spacedlowsmallcaps{\contentsname}}
\tableofcontents
\automark[section]{chapter}
\renewcommand{\chaptermark}[1]{\markboth{\spacedlowsmallcaps{#1}}{\spacedlowsmallcaps{#1}}}
\renewcommand{\sectionmark}[1]{\markright{\thesection\enspace\spacedlowsmallcaps{#1}}}
%*******************************************************
% List of Figures and of the Tables
%*******************************************************
\clearpage

\begingroup
    \let\clearpage\relax
    \let\cleardoublepage\relax
    \let\cleardoublepage\relax
    %*******************************************************
    % List of Figures
    %*******************************************************
    %\phantomsection
    \refstepcounter{dummy}
    %\addcontentsline{toc}{chapter}{\listfigurename}
    \pdfbookmark[1]{\listfigurename}{lof}
    \listoffigures

    \vspace{8ex}

    %*******************************************************
    % List of Tables
    %*******************************************************
    %\phantomsection
    \refstepcounter{dummy}
    %\addcontentsline{toc}{chapter}{\listtablename}
    \pdfbookmark[1]{\listtablename}{lot}
    \listoftables

    \vspace{8ex}
%   \newpage

    %*******************************************************
    % List of Listings
    %*******************************************************
      %\phantomsection
    %\refstepcounter{dummy}
    %\addcontentsline{toc}{chapter}{\lstlistlistingname}
    %\pdfbookmark[1]{\lstlistlistingname}{lol}
    %\lstlistoflistings

    %\vspace{8ex}

    %*******************************************************
    % Acronyms
    %*******************************************************
    %\phantomsection
    %\refstepcounter{dummy}
    %\pdfbookmark[1]{Acronyms}{acronyms}
    %\markboth{\spacedlowsmallcaps{Acronyms}}{\spacedlowsmallcaps{Acronyms}}
    %\chapter*{Acronyms}
    %\begin{acronym}[UMLX]
    %    \acro{DRY}{Don't Repeat Yourself}
    %    \acro{API}{Application Programming Interface}
    %    \acro{UML}{Unified Modeling Language}
    %\end{acronym}                     
\endgroup

%********************************************************************
% Mainmatter
%*******************************************************
\cleardoublepage\pagenumbering{arabic}
%\setcounter{page}{90}
% use \cleardoublepage here to avoid problems with pdfbookmark
\cleardoublepage

% Parte 1: Matemáticas.
% Homología Singular

\ctparttext{En esta parte, vamos a ver una introducción a la homología singular.
  Definiremos los grupos de homología singular de un espacio topológico y
  estudiaremos algunas de sus propiedades. Continuaremos demostrando el teorema de subdivisión
  baricéntrica, resultado que nos lleva a la obtención de herramientas para el
  cálculo de la homología, y las aplicaremos al caso particular de las esferas.
  Finalmente, demostraremos resultados clasicos haciéndonos valer de nuevas
  construcciones, unidas a los resultados ya demostrados.}
\part{Homología Singular}
%************************************************
\label{ch:introduccionHOM}
%************************************************

\section{Introducción al trabajo}

El propósito de este trabajo es introducir la teoría de homología singular sobre
espacios topológicos. Esto lo haremos de forma constructiva, introduciendo los conceptos
y estructuras necesarios para poder definir la homología singular sobre un espacio topológico.
Tras definir la homología, estudiaremos sus propiedades más inmediatas y algunos casos sencillos,
y extenderemos el concepto para pares topológicos, trasladando los resultados obtenidos
con anterioridad.

Una vez asentadas las bases, trataremos el teorema de subdivisión baricéntrica, el
resultado más importante del trabajo, que nos permite obtener herramientas para el
cálculo de la homología de espacios topológicos. Será necesario introducir nociones
geométricas y demostrar ciertas propiedades de estas nuevas construcciones para la
demostración del teorema.

Demostrado el teorema de subdivisión baricéntrica, trataremos sus consecuencias,
siendo las más importantes la sucesión de Mayer-Vietoris y el teorema de escisión.
Con ellas, seremos capaces de calcular la homología de las esferas, y lo haremos
de dos formas distintas; una mediante escisión, y otra usando la sucesión de Mayer-Vietoris.

Después de calcular la homología de las esferas demostraremos un número de resultados clásicos,
para lo que definiremos la homología local de un espacio topológico. Obtendremos resultados
como el teorema del punto fijo de Brower, el teorema de separación de Jordan-Brower y
ciertos teoremas de invarianza. Finalmente, apoyándonos en la definición del grado de
una aplicación continua sobre una esfera, seremos capaces de demostrar resultados
sobre campos definidos en esferas.

\section{Objetivos}
Los objetivos de este trabajo son los siguientes:

\begin{itemize}
  \item Construcción de la homología singular y estudio de sus propiedades.

  \item Cálculo de la homología singular de las esferas y obtención de aplicaciones.

\end{itemize}

El primer objetivo se ha cumplido satisfactoriamente, y está desarrollado tanto en el
\autoref{ch:introduccion} como en el \autoref{ch:calculoI}.

El segundo objetivo también se ha cumplido, con el cálculo de la homología de las esferas
al final del \autoref{ch:calculoI}, y la obtención de aplicaciones desarrollada en el
\autoref{ch:resultadosHom}.

\section{Organización de la memoria}

El \autoref{ch:introduccion} contiene todo lo referente a la definición de la homología
singular de un espacio topológico: definiciones, propiedades y construcciones algebraicas
y geométricas. También se demostrará la invarianza homotópica de la homología, y se
definirá la homología de los pares topológicos.

En el \autoref{ch:calculoI} se demostrará el teorema de subdivisión baricéntrica y
se obtendrán herramientas para el cálculo de la homología singular: el teorema
de escisión y la sucesión de Mayer-Vietoris. Estos serán utilizados para deducir
la homología de las esferas.

Finalmente, en el \autoref{ch:resultadosHom} obtendremos resultados clásicos, como el
teorema del punto fijo de Brower, el teorema de separación de Jordan-Brower y algunos
teoremas de invarianza. Para ello, nos haremos valer de los conceptos de homología
local y de grado de una aplicación continua definida en una esfera.

\section{Referencias}

Para la realización de este trabajo se ha utilizado principalmente el libro de
teoría de homología de James W. Vick \citep{vick:1973}, complementado por los títulos
de William S. Massey \citep{massey:1972} y Sze-Tsen Hu \citep{hu:1966}, que también
tratan la teoría de homología.


%************************************************
\chapter{Introducción}\label{ch:introduccion}
%************************************************

\section{Definiciones iniciales}

En este capítulo, introduciremos la \textbf{homología singular} de un espacio topológico. Comenzaremos con
un conjunto de definiciones iniciales, introduciendo la noción de \textbf{símplice singular}. Con ellos,
construiremos un conjunto al que dotaremos de estructura algebraica, y sobre el que definiremos una aplicación
que, gracias a sus propiedades, nos permitirá definir la homología singular del espacio. Seguidamente,
demostraremos algunos resultados sencillos sobre ciertos espacios topológicos, continuando con la invarianza
homotópica de la homología. Concluiremos el capítulo definiendo la homología sobre pares topológicos, y
viendo como se trasladan los resultados demostrados con anterioridad.

\begin{definition}
Un \textbf{p-símplice} $s_p$ es la envolvente conexa de $p+1$ puntos afínmente independientes en $\mathbb{R}^p$,
a los que llamaremos vértices, $\{x_0, \dots, x_n\}$, esto es
\[ s_p = \{\lambda_0 x_0 + \dots + \lambda_p x_p \mid \lambda_0 + \dots + \lambda_p = 1, \lambda_i \geq 0 \hspace{0.2em} \forall i \} \]
\end{definition}
% \underbrace{\smash{1}}_{i}
Cuando los vértices son $x_i = (0, \dots, 1, \dots, 0)$, con el $1$ en la posición $i$-ésima, le llamaremos el p-símplice \underline{standard},
y lo notaremos por $\sigma_p$.

Podemos establecer un homeomorfismo $h \colon \sigma_p \to s_p$ por
\[ h(t_0, \dots, t_p) = \sum\limits_{i = 0}^p t_i x_i \]

\begin{definition}
Sea $\X$ un espacio topológico. Un \textbf{p-símplice singular} en $\mathbb{X}$ es una aplicación continua $\phi \colon \sigma_p \to \X$ para $p \geq 0$.
\end{definition}

Representaremos el conjunto de p-símplices singulares de $\mathbb{X}$ como
\[F_p(\mathbb{X}) = \{\phi \colon \sigma_p \to \mathbb{X} \mid \phi \text{ es p-símplice singular}\}\]
Es claro que $\mathbb{F}_0(\mathbb{X}) = \mathbb{X}$. Además, $F_1(\mathbb{X})$ son los arcos en $\X$, identificables por
$\mathbb \mathcal{C}([0,1], \X)$.

Si tenemos una aplicación continua entre espacios topológicos $f \colon \X \to \Y$, podemos definir
\[f_{\#} \colon \FX{p} \to \FY{p}  \quad \text{ por }  \quad f_{\#} = f \circ \phi \colon \sigma_p \to Y \]
que es una aplicación entre los p-símplices singulares de ambos espacios.

Se verifica trivialmente que $I_\# = Id$, puesto que $I_\#(\phi) = I \circ \phi = \phi$.
Además, se tiene que $(g \circ f)_\# = g_\# \circ f_\#$.

\begin{definition}
  Sea $p \geq 1$, $0 \leq i < p$. Definimos $F_p^i$ como la aplicación de un $(p-1)$-símplice a un p-símplice de la siguiente forma:
  \begin{align*}
    F_p^i \colon \sigma_{p-1} &\to \sigma_p\\
    (t_0, \dots, t_{p-1}) &\mapsto (t_0, \dots, t_{i-1}, 0, t_i, \dots, t_{p-1})
  \end{align*}

  Podemos definir una nueva aplicación a partir de esta, a la que llamaremos \textbf{i-esima cara}:
  \begin{align*}
    \partial_i \colon \FX{p} &\to \FX{p-1}\\
    \partial_i(\phi) &= \phi \circ F_p^i
  \end{align*}
\end{definition}

Vamos a probar un resultado sobre la aplicación que acabamos de definir.

\begin{lemma}
  Se verifica que $F_p^i \circ F_{p-1}^j = F_p^j \circ F_{p-1}^{i-1} \quad \forall p \geq 2$, $\forall i, j \colon 0 \leq j < i \leq p$.
\end{lemma}

\begin{proof}
Lo comprobamos realizando los cálculos pertinentes
  \begin{align*}
    &F_p^i \circ F_{p-1}^j (t_0, \dots, t_{p-2}) = F_p^i (t_0, \dots, t_{j-1}, 0, t_j, \dots, t_{p-2}) \\
    &= (t_0, \dots, t_{j-1}, 0, t_j, \dots, t_{i-2}, 0, t_{i-1}, \dots, t_{p-2}) \\
    &= F_p^j (t_0, \dots, t_{i-2}, 0, t_{i-1}, \dots, t_{p-2}) = F_p^j \circ F_{p-1}^{i-1} (t_0, \dots, t_{p-2}).
  \end{align*}
\end{proof}

Se deduce directamente este corolario:

\begin{corollary}
  Para $p \geq 2, 0 \leq j < i \leq p$, se verifica $\partial_j \circ \partial_i = \partial_{i-1} \circ \partial_j$.
\end{corollary}

Debido a que $\FX{p}$ no está dotado de estructura algebraica, es posible realizar la siguiente construcción:

\begin{definition}
  Si $\X$ es un espacio topológico, y $p \geq 0$, se define el \textbf{grupo de p-cadenas singulares} de $\X$
  como el grupo abeliano libre generado por $\FX{p}$
  \[ S_p(\X) = \{\sum_{\phi \in \FX{p}} n_\phi \cdot \phi  \mid n_\phi \in \Z \text{, todos cero salvo un número finito} \}\]

  Podemos definir un operador \textbf{borde} $\partial$ con las aplicaciones cara $\partial_i$ de la siguiente forma:
  \begin{align*}
    \partial \colon S_p(\X) &\to S_{p-1}(\X)\\
    \partial(\phi) &= \sum_{i = 0}^p (-1)^i \partial_i(\phi)  \hspace{1cm} (p \geq 1)
  \end{align*}
  considerando $\partial_i$ como una extensión de la aplicación cara a un homomorfismo:
  \[\partial_i(\sum_{\phi \in F_p} n_\phi \cdot \phi) = \sum_{\phi \in F_p} n_\phi \cdot \partial_i(\phi) \]
\end{definition}

Por conveniencia, consideraremos $\partial(S_0(\X)) = S_{-1}(\X) = \{0\}$.

% Motivar (mejor) el resultado siguiente?

Veamos un resultado importante sobre la aplicación borde que acabamos de definir.

\begin{proposition}
  $\partial \circ \partial \colon \SX{n} \to \SX{n-2}$ es el homomorfismo cero.
\end{proposition}

Geométricamente, este resultado nos dice que el borde de una $p$-cadena es una $(p-1)$-cadena sin
borde. Esta propiedad nos permitirá definir los grupos de homología singular.

\begin{proof}
  Consideramos $S_p(\X) \xrightarrow{\partial} S_{p-1}(\X) \xrightarrow{\partial} S_{p-2}(\X)$. \\
  Es claro que si $p = 0, 1$, la cuestión es trivial, pues $S_{-1} = \{0\}$. Tomaremos $p > 1$.

  Sea $\phi \in \FX{p}$.
  \begin{align*}
    \partial^2(\phi) &= \partial(\sum_{i = 0}^p (-1)^i \partial_i(\phi)) = \sum_{i = 0}^p (-1)^i \partial(\partial_i(\phi)) \\
                     &= \sum_{i = 0}^p (-1)^i \sum_{j = 0}^{p-1} (-1)^j \partial_j \partial_i(\phi) \\
                     &= \sum_{0 \leq j < i \leq p} (-1)^{i + j} \partial_j \partial_i(\phi)
                        + \sum_{0 \leq i \leq j \leq p-1} (-1)^{i + j} \partial_j \partial_i(\phi) \\
                     &= \sum_{0 \leq j < i \leq p} (-1)^{i + j} \partial_{i-1} \partial_j(\phi)
                        + \sum_{0 \leq i \leq j \leq p-1} (-1)^{i + j} \partial_j \partial_i(\phi) \\
                     &= \sum_{0 \leq j \leq i-1 \leq p} (-1)^{i + j} \partial_{i-1} \partial_j(\phi)
                        + \sum_{0 \leq i \leq j \leq p-1} (-1)^{i + j} \partial_j \partial_i(\phi) = 0\\
                     &\text{pues son iguales salvo el signo.}
  \end{align*}
\end{proof}

Hasta ahora, hemos tomado un espacio topológico $\X$ y le hemos asignado un conjunto de p-símplices singulares $\FX{p}$, del que
hemos tomado un $\Z$-módulo, obteniendo el grupo abeliano libre de p-cadenas singulares, $S_p(\X)$. Sobre él, hemos definido un
operador borde $\partial$, que verifica $\partial^2 = 0$. A la familia $\{S_p(\X), \partial\}_{p \geq 0}$ se le llama el complejo de cadenas
singular asociado a $\X$.

Representaremos también por $f_\# \colon S_p(\X) \to S_p(\Y)$ a la composición entre
$f \colon \X \to \Y$, continua, y $\phi \in S_p(\X)$.
\[ f_\#(\sum_\phi n_\phi \cdot \phi) = \sum_\phi n_\phi \cdot (f \circ \phi) \]

Vamos a ver que esta aplicación conmuta con el borde, es decir, $f_\# \circ \partial = \partial \circ f_\#$. Esto hace conmutativo el siguiente diagrama:
\[
  \begin{tikzcd}
    \dots S_{p+1}(\X) \arrow{r}{\partial} \arrow{d}{f_\#} & S_p(\X) \arrow{r}{\partial} \arrow{d}{f_\#} & S_{p-1}(\X) \arrow{d}{f_\#} \dots \\
    \dots S_{p+1}(\Y) \arrow{r}{\partial}                 & S_p(\Y) \arrow{r}{\partial}                 & S_{p-1}(\Y) \dots
  \end{tikzcd}
\]

Su demostración es casi inmediata, pues se tiene que \[ f_\# \circ \partial_i(\phi) = f_\#(\phi \circ F_p^i) = f \circ \phi \circ F_p^i
                                             = \partial_i(f \circ \phi) = \partial_i \circ f_\#(\phi) \]

El resultado obtenido implica, por la definición de $\partial$, que $\partial f_\# = f_\# \partial$. De hecho, es más fuerte, pues vemos
como las aplicaciones cara conmutan una a una con el homomorfismo $f_\#$.

A esta aplicación entre complejos de cadenas se le llama en álgebra una aplicación de cadenas. Nos referiremos a ella como la aplicación de cadenas
inducida por la aplicación continua f.

\begin{definition}
  Diremos que $c \in S_p(\X)$ es un \textbf{p-ciclo} si $\partial(c) = 0$. \\
  Diremos que $c \in S_p(\X)$ es un \textbf{p-borde} si existe $d \in S_{p+1}(\X) \text{ tal que } \partial(d) = c$.
\end{definition}

% ¿Indentar los conjuntos?

Representaremos por: \\
$Z_p(\X) = \{c \in S_p(\X) \mid \text{ c es un p-ciclo}\}$ \\
$B_p(\X) = \{c \in S_p(\X) \mid \text{ c es un p-borde}\} $

Se verifica que: \\
$Z_p(\X) = \Ker (\partial \colon S_p(\X) \to S_{p-1}(\X)) $ \\
$B_p(\X) = \Img (\partial \colon S_{p+1}(\X) \to S_p(\X)) $

Ambos son subgrupos de $S_p(\X)$, llamados grupos de los ciclos y los bordes, respectivamente. Como hemos visto que $\partial^2 = 0$, se deduce
inmediatamente que $B_p(\X) \subseteq Z_p(\X) \hspace{0.2em}\forall p \geq 0$, luego se puede construir el grupo cociente.
\begin{equation*}
  H_p(\X) = \ddfrac{Z_p(\X)}{B_p(\X)}
\end{equation*}
al que llamaremos el \textbf{\underline{\smash{p-ésimo grupo de homología singular}}}, el cual es un grupo abeliano.

\begin{remark}
  Nótese que $Z_0(\X) = S_0(\X)$, pues $\partial(S_0(X)) = 0$.
\end{remark}

A los elementos de $H_p(\X)$ los notaremos $[c]$, con $\partial(c) = 0$, esto es, $c \in Z_p(\X)$, y los llamaremos clases de homología.

La motivación geométrica para la construcción de los grupos de homología singular se basa en querer
estudiar los ciclos en espacios topológicos. Como este conjunto tiene un tamaño demasiado grande para
poder estudiarlo de forma efectiva, restringimos nuestra atención a las clases de equivalencia, haciendo
que dos ciclos sean equivalentes si su diferencia es el borde de una cadena en una dimensión superior.

Tomemos nuevamente una función continua $f \colon \X \to \Y$, y definamos
\[ f_* \colon H_p(\X) \to H_p(\Y) \quad f_*([c]) = [f_\#(c)]  \]

Comprobemos que está bien definida, esto es, que no depende del representante que se elija.

Sean $[c_1] = [c_2]$, lo que implica $c_1 - c_2 \in B_p(\X)$. Por tanto, existe $d \in S_{p+1}(\X)$ tal que $\partial(d) = c_1 - c_2$.
Así, $f_\#(c_1) - f_\#(c_2) = f_\#(c_1 - c_2) = f_\#(\partial(d)) = \partial(f_\#(d))$ \\
Como $f_\#(d) \in S_{p+1}(\Y)$, se tiene $[f_\#(c_1)] = [f_\#(c_2)]$

Además, $\partial(f_\#(c)) = f_\#(\partial(c)) = 0$, por lo que $f_\#(c)$ es un ciclo en $\Y$. Con esto hemos visto que $f_*$ está bien definida,
y es un homomorfismo, al serlo $f_\#$, al que llamaremos homomorfismo inducido por f en la homología.

De las propiedades de $f_\#$ obtenemos directamente que:
\begin{itemize}
  \item $(g \circ f)_* = g_* \circ f_*$
  \item $I_* = Id$
\end{itemize}

Así, hemos construído un funtor covariante de la categoría de espacios topológico en la categoría de grupos abelianos,
el funtor de la homología singular.
\[  \begin{tikzcd}
  \X \arrow{r} \arrow{d}{f} & [5em] H_p(\X) \arrow{d}{f_*} \\
  \Y \arrow{r}  & H_p(\Y)
\end{tikzcd} \]

Si $\X$ e $\Y$ son homeomorfos, el homomorfismo que induce dicho homeomorfismo en la homología es un isomorfismo de grupos de homología, pues
basta considerar el inverso del homeomorfismo para construir el inverso del homomorfismo inducido.

\section{Homología de un punto}

Vamos a calcular los grupos de homología singular de los puntos, sin más que ver como se comporta la función borde
en ellos y utilizando la definición de los grupos de ciclos y bordes como núcleo e imagen de dicha aplicación.

Sea $\X$ un espacio topológico formado por un punto, $\X = \{x\}$.
$\FX{p} = \{ \phi \colon \sigma_p \to \{x\} \} = \{\phi_p\}$ donde $\phi_p$ es la función constante.

$S_p({\{x\}})$ son grupos abelianos libres con $\phi_p$ como generador. Además, $\partial \colon \SX{p} \to \SX{p-1}$ verifica:
\[\partial(\phi_p) = \sum_{i = 0}^p (-1)^i \partial_i(\phi_p) =  \sum_{i = 0}^p (-1)^i \phi_{p-1} = \begin{cases}
                                                                                                            0 & p \text{ es impar} \\
                                                                                                            \phi_{p-1} & p \text{ es par}
                                                                                                    \end{cases}   \]

Así, si p es impar, $S_{p+1}(\{x\}) \to S_p(\{x\}) \xrightarrow{0} S_{p-1}(\{x\})$ y $Z_p(\{x\}) = B_p(\{x\}) = S_p(\{x\})$, luego
$H_p(\{x\}) = 0$.

Si p es par, con $p > 0$, $Z_p(\{x\}) = B_p(\{x\}) = 0$, y $H_p(\{x\}) = 0$.

Por último, si $ p = 0, H_0(\{x\}) = \ddfrac{S_0(\{x\})}{B_0(\{x\})} = S_0(\{x\}) \cong \mathbb Z$.

Así, $H_p(\{x\}) = 0 \hspace{1em} \forall p \geq 1, \hspace{1em} H_0(\{x\}) \cong \mathbb Z$.

\section{El grupo $H_0(\X)$ y la arcoconexión}

En esta sección vamos a estudiar algunas de las propiedades de los grupos de homología singular, asumiendo que el espacio
sobre el que trabajamos es arcoconexo.

\begin{proposition}
  Sea $\X$ un espacio topológico arcoconexo. Entonces $H_0(\X) \cong \mathbb Z$.
\end{proposition}

\begin{proof}
  Consideramos $\SX{1} \xrightarrow{\partial} \SX{0} \to \{0\}$. Tomamos $H_0(\X)$.

  $H_0(\X) = \ddfrac{Z_0(\X)}{B_0(\X)} = \ddfrac{S_0(\X)}{\Img \partial}$

  Sea $\epsilon \colon S_0(\X) \to \Z$ definido por \[\epsilon(\sum\limits_{\phi} n_\phi \phi) = \epsilon(\sum\limits_{x} n_x \cdot x) = \sum\limits_{x} n_x \]
  suma finita, por cómo definimos $S_0$.

  Se verifica que $\epsilon$ es un epimorfismo de grupos, pues podemos definir el conveniente elemento de $S_0(\X)$ para obtener
  cualquer entero. Calculamos el núcleo de la aplicación:
  \[ \epsilon \circ \partial(\phi_1) = \epsilon(\sum_{i=0}^1 \partial_i(\phi_i)) = 1 - 1 = 0\]  Así, $\Img \partial \subseteq \Ker \epsilon$.

  Sea ahora $\sum n_x \cdot x \in \Ker \epsilon$, esto es, $\sum n_x = 0$. \\
  $\sum n_x \cdot x = n_1 x_1 + \dots + n_k x_k$ con $\sum\limits_{i=1}^k = 0, x_i \in \X$. \\
  Sean $\phi^i \colon \sigma_1 \to \X$ arcos tales que $\phi^i(e_1) = x_i, \phi^i(e_0) = x$,
  donde $e_1, e_0$ representan los vértices de $\sigma_1$ y $x \in \X$ es un punto fijo.
  Es posible hacer esta construcción gracias a la arcoconexión de $\X$. De esta forma:
  \begin{align*}
    &\sum_i n_i \phi^i \in \SX{1}, \text{ y se tiene} \\
    &\partial(\sum_i n_i \phi^i) = \sum_i n_i \partial(\phi^i) = \sum_i n_i(x_i - x) = (\sum_i n_i x_i) - (\sum_i n_i) x = \sum_i n_i x_i
  \end{align*}
  De aquí se deduce que $\Ker \epsilon \subseteq \Img \partial$, y por tanto tenemos la igualdad.

  Finalmente como $\Ker \epsilon = \Img \partial$, $H_0(\X) = \ddfrac{S_0(\X)}{\Ker \epsilon} \cong \Z$, al ser $\epsilon$ epimorfismo.
\end{proof}

\begin{proposition}
  Sea $\X$ un espacio topológico, $\X = \bigcup\limits_{\alpha \in A} \X_\alpha$ su descomposición en componentes arcoconexas. Entonces se verifican:
  \begin{itemize}
    \item[a)] $\forall \alpha \in A$, la función $i_{\alpha*} \colon H_p(\X_\alpha) \to H_p(\X)$ es un monomorfismo, donde la aplicación
              $i_\alpha \colon \X_\alpha \to \X$ representa la inclusión.
    \item[b)] $H_p(\X) = \bigoplus\limits_{\alpha \in A} i_{\alpha*}(H_p(\X_\alpha)) \hspace{0.5em} \forall p \geq 0$. En particular, por la inyectividad de $i_{\alpha*}$,
              \[ H_p(\X) \cong \bigoplus_{\alpha \in A} H_p(\X_\alpha) \]
  \end{itemize}
\end{proposition}

\begin{proof}
  Sea $\phi \colon \sigma_p \to \X$ un p-símplice singular. Como $\sigma_p$ es convexo y $\phi$ es continua, $\phi(\sigma_p) \subseteq \X_\alpha$, para algún
  $\alpha \in A$, por lo que $\phi = i_{\alpha*}(\phi)$ para dicho $\alpha$. \\
  Si $c = \sum n_\phi \cdot \phi$ es una p-cadena, poniendo cada $\phi$ de la forma anterior y agrupando por componentes conexas:
  \begin{equation}
    c = \sum\limits_{\alpha \in A} i_{\alpha*}(c_\alpha), \quad c_\alpha \in S_p(\X_\alpha)  \tag{*}
  \end{equation}

  Sea $[c]_\alpha \in H_p(\X_\alpha)$, y supongamos que $i_{\alpha*}[c]_\alpha = 0$ en $\H_p(\X)$. En ese caso, $[i_{\alpha*}(c)] = 0$ en $H_p(\X)$,
  luego existe $d \in \SX{p+1}$ tal que $i_{\alpha\#}(c) = \partial(d)$, y además:
  \[\partial(d) = \sum\limits_{\beta \in A} i_{\beta\#}(\partial(d_\beta)) = i_{\alpha\#}(c) \]
  por lo que se tiene $c = \partial(d_\alpha), \partial(d_\beta) = 0 \hspace{0.5em} \forall \beta \neq \alpha$. Así, $[c]_\alpha = 0$ y $i_{\alpha*}$ es un monomorfismo.
  Esto prueba a).

  Observamos que la suma en (*) es finita, pues si no, habría un número no finito de generadores en la expresión de c. Así:
  \[  [c] = \sum\limits_{\alpha \in A} [i_{\alpha\#}(c_\alpha)] =  \sum\limits_{\alpha \in A} i_{\alpha*}([c_\alpha]_\alpha)\]

  Como $\partial(c) = 0 \iff \partial(c_\alpha) = 0 \hspace{0.5em} \forall \alpha \in A$, se obtiene directamente b), como consecuencia de a).
\end{proof}

En particular, para $H_0(\X)$, por el resultado anterior, se tiene que $H_0(\X)$ es el grupo abeliano libre con tantos generadores como
componentes arcoconexas tenga $\X$.

\section{Invarianza homotópica de la homología singular}

En esta sección demostraremos la invarianza de los grupos de homología singular en espacios homotópicamente equivalentes, viendo antes como
dos aplicaciones homotópicas inducen el mismo homomorfismo en la homología.

Vamos a comenzar demostrando el llamado Lema de Poincaré.

\begin{lemma}[Lema de Poincaré]
  Sea $\X \subseteq \R^n$ estrellado desde algún punto. Entonces:
  \[ H_p(\X) = \begin{cases}  \Z & \text{ si } p = 0 \\
                              0  & \text{ si } p \neq 0
                            \end{cases} \]
\end{lemma}

\begin{proof}
  El caso $p = 0$ está claro, pues un conjunto estrellado es, en particular, arcoconexo, por lo que solo hemos de aplicar el resultado ya demostrado.
  Suponemos $p > 0$.

  Para demostrarlo, vamos a construir un homomorfismo $ T: \SX{p} \to \SX{p+1}$ tal que $ T \partial + \partial  T = Id_{\SX{p}}$ para $p > 0$.

  De darse la situación anterior, si $[c] \in H_p(\X), c \in Z_p(\X)$, se tiene $\partial(c) = 0$ y $c \in S_p(\X)$.
  En consecuencia, $c = (\partial  T +  T \partial)(c)  = \partial( T(c))$ con $ T(c) \in \SX{p+1}$.
  Por tanto, $c$ es un borde, luego $[c] = 0$, y eso implica que $H_p(\X) = 0 \quad \forall p > 0$.

  Definamos $T$: \\
  Sean $x_0 \in \X$ un punto de estrella, $\phi \colon \sigma_p \to \X$ un p-símplice singular. Tomemos $T(\phi) \colon \sigma_{p+1} \to \X$ dado por:
  \[ T(\phi)(t_0,\dots,t_{p+1}) = \begin{cases} x_0 & \text{ si } t_0 = 1 \\
                                                (1-t_0)\phi(\frac{t_1}{1-t_0},\dots,\frac{t_{p+1}}{1-t_0}) + t_0 x_0 & \text{ si } t_0 < 1
                                  \end{cases} \]
  Como $\X$ es estrellado, $\Img(T(\phi)) \subseteq \X$.

  Veamos que $T(\phi)$ es continuo. \\
  La continuidad está clara para $t_0 < 1$, por lo que la vemos para $t_0 = 1$. Se verifica que:
  \begin{align*}
    &|T(\phi)(t_0,\dots,t_{p+1}) - x_0| = |t_0 x_0 + (1-t_0)\phi(\frac{t_1}{1-t_0},\dots,\frac{t_{p+1}}{1-t_0}) - x_0| \\
    &= (1-t_0)|\phi(\frac{t_1}{1-t_0},\dots,\frac{t_{p+1}}{1-t_0}) - x_0| \leq (1-t_0)(c + |x_0|) \\
    & \text{ pues $\phi$ está acotada.}
  \end{align*}

  Así, $\lim\limits_{t_0 \to 1} T(\phi) = x_0$.

  Queda claro que $\partial_0(T(\phi)) = \phi$. Si seguimos llamando $T$ al homomorfismo extensión de $T$ a $\SX{p}$, tenemos construído
  \[  T \colon \SX{p} \to \SX{p+1} \hspace{0.5em} \text{tal que } \partial_0 \circ  T = I. \]
  Si $p = 0 \implies i = 1$ y $\partial_1( T(\phi)) = x_0$.

  Si $1 \leq i \leq p+1$, entonces:
  \begin{align*}
    &\partial_i( T(\phi))=(t_0,\dots,t_p)) =  T(\phi)(t_0, \dots, t_{i-1} \circ t_i, \dots, t_p) \\
    &= (1-t_0)\phi(\frac{t_1}{1-t_0}, \dots, \frac{t_{i-1}}{1-t_0}, \dots, \frac{t_p}{1-t_0}) + t_0 x_0 \\
    &= (1-t_0)\partial_{i-1}(\frac{t_1}{1-t_0}, \dots, \frac{t_p}{1-t_0}) + t_0 x_0 = T(\partial_{i-1} \phi)(t_0, \dots, t_p)
  \end{align*}
  Como consecuencia:
  \begin{align*}
    \partial( T(\phi)) &= \sum\limits_{i = 1}^{p+1} (-1)^i \partial_i  T(\phi) = \phi + \sum\limits_{i = 1}^{p+1} (-1)^i  T \partial_{i-1}(\phi) \\
    &= \phi - \sum\limits_{j = 0}^p (-1)^j  T \partial_j(\phi) = \phi -  T \partial(\phi)
  \end{align*}

  De donde se obtiene directamente que $\partial  T +  T \partial = Id$.
\end{proof}

\begin{theorem}[Invarianza homotópica de la homología singular]
  Sean $\X, \Y$ espacios topológicos y $f, g \colon \X \to \Y$ aplicaciones continuas. Si $f$ es homotópica a $g$,
  entonces $f_* = g_* \colon H_p(\X) \to H_p(\Y) \hspace{0.5em} \forall p \geq 0$.
\end{theorem}

\begin{proof}
  Supongamos que $\forall p \geq 0$ podemos definir homomorfismos \[T: \SX{p} \to \SY{p+1} \] tales que $T\partial + \partial T = f_\# - g_\#$,
  esto es, $f$ y $g$ son homotópicas como aplicaciones de cadenas. Entonces, $\forall [c] \in H_p(\X)$ se tiene:
  \[ (f_* - g_*)([c]) = [f_\#(c) - g_\#(c)] = [(T\partial + \partial T)(c)] = [\partial(T(c))] = 0\]
  de donde se deduce $f_* = g_*$.

  Veamos la definición de dicho $T$.

  Sea $H \colon \X \times I \to \Y$ la homotopía entre f y g, esto es,
  \[H(x, 0) = f(x), \quad H(x, 1) = g(x)\]
  donde $I$ representa el intervalo $[0, 1]$.

  Sean $i_0, i_1 \colon \X \to \X \times I$ dados por $i_0(x) = (x, 0), i_1(x) = (x, 1)$. Es claro que $f = H \circ i_0, g = H \circ i_1$, de donde
  se tiene \[ f_\# - g_\# = H_\# \circ i_{0\#} - H_\# \circ i_{1\#} = H_\# \circ (i_{0\#} - i_{1\#}) \]

  Si somos capaces de construir $\tau \colon \SX{p} \to S_{p+1}(\X \times I)$ tal que $\partial \tau + \tau \partial = i_{0\#} - i_{1\#}$, definiendo
  $T = H_\# \circ \tau$, se verifica:
  \[ \partial \tau + \tau \partial = \partial(H_\# \circ \tau) + (H_\# \circ \tau)\partial = H_\#(i_{0\#} - i_{1\#}) = f_\# - g_\#  \]
  lo que demostraría el teorema.

  Construiremos $\tau$ por inducción. Para ello, vamos a probar el siguiente resultado:
  \begin{lemma}
    Para todo $p \geq 0$ y para todo espacio topológico $\X$, existe un homomorfismo de grupos $T^\X \colon \SX{p} \to S_{p+1}(\X \times I)$ verificando:
    \begin{itemize}
      \item[a)] $\partial T^\X + T^\X \partial = (i_0^\X)_\# - (i_1^\X)_\#$
      \item[b)] (Condición de naturalidad) Para todo espacio topológico $\Y$ y toda aplicación continua $h \colon \X \to \Y$, el diagrama
      \[
      \begin{tikzcd}
        \SX{p} \arrow{r}{T^\X} \arrow{d}{h_\#} & [2em] S_{p+1}(\X \times I) \arrow{d}{(h \times Id)_\#} \\
        \SY{p} \arrow{r}{T^\Y}                 & S_{p+1}(\X \times I)
      \end{tikzcd}
      \quad \text{es conmutativo.}
      \]
    \end{itemize}
    Nótese que en b) pedimos más de lo requerido, pero nos sirve para definir $T^\X$.
  \end{lemma}

  Demostrémoslo para $p = 0$.

  Tomamos $\X = \sigma_0$ e intentamos definir $T^{\sigma_0} \colon S_0(\sigma_0) \to S_1(\sigma_0 \times I)$ \\
  Si $\tau_0 \colon \sigma_0 \to \sigma_0$ es la identidad, entonces $\tau_0$ genera $S_0(\sigma_0)$. Como $T^{\sigma_0}$
  debe verificar $\partial T^\sigma_0 + T^\sigma_0 \partial = (i_0^{\sigma_0})_\# - (i_1^{\sigma_0})_\#$, se deduce
  $\partial T^{\sigma_0}(\sigma_0) = (i_0^{\sigma_0})_\# (\tau_0) - (i_1^{\sigma_0})_\# (\tau_0)$

  Definiendo:
  \begin{align*}
    T^{\sigma_0}(\tau_0) \colon \sigma_1 &\to \sigma_0 \times I \\
                       (t, 1-t) &\mapsto (\sigma_0, 1-t)
  \end{align*}
  es claro que verifica lo anterior.

  Sea ahora $\X$ espacio topológico, y sea $\phi \in S_0(\X)$ un generador, $\phi \colon \sigma_0 \to \X$.
  Entonces, como la construcción debe ser natural, el diagrama:
  \[
  \begin{tikzcd}
    S_0(\sigma_0) \arrow{r}{T^{\sigma_0}} \arrow{d}{\phi_\#} & [2em] S_1(\sigma_0 \times I) \arrow{d}{(\phi \times Id)_\#} \\
    \SX{0} \arrow{r}{T^\X}                 & S_1(\X \times I)
  \end{tikzcd}
  \]
  debe ser conmutativo, luego $T^\X \circ \phi_\# = (\phi \times Id)_\# \circ T^{\sigma_0}$.

  Como $\phi_\#(\tau_0) = \phi$, se tiene $T^\X(\phi) = (\phi \times Id)_\# \circ T^{\sigma_0}(\tau_0)$.

  Definimos $T^\X$ de esta manera sobre los generadores, y lo extendemos por linealidad. Veamos que verifica a) y b).

  a)
  \begin{align*}
    \partial T^\X(\phi) &= \partial(\phi \times Id)_\# T^{\sigma_0}(\tau_0) = (\phi \times Id)_\# \partial T{ \sigma_0}(\tau_0) \\
                        &= (\phi \times Id)_\# ((i_0)_\#^{\sigma_0}(\tau_0) - (i_1)_\#^{\sigma_0}(\tau_0)) \\
                        &= ((\phi \times Id) (i_0)_\#^{\sigma_0}) - (\phi \times Id) (i_1)_\#^{\sigma_0}))(\tau_0) \\
                        &= ((i_0^\X \circ \phi)_\# - (i_1^\X \circ \phi)_\#) (\tau_0) = ((i_0^\X)_\# - (i_1^\X)_\#)(\phi)
  \end{align*}
  de donde se obtiene $\partial T^\X + T^\X \partial = (i_0^\X)_\# - (i_1^\X)_\#$.

  b) Si $h \colon \X \to \Y$ es continua, entonces:
  \begin{align*}
    &(T^\Y \circ h_\#)(\phi) = T^\Y(h \circ \phi) = (h\phi \times Id)_\# T^{\sigma_0}(\tau_0) \\
    &= ((h \times Id)_\# \circ (\phi \times Id)_\#) (T^{\sigma_0}(\tau_0)) = ((h \times Id)_\# \circ T^\X)(\phi)
  \end{align*}
  lo que finaliza la demostración para $p = 0$.

  Suponemos el resultado cierto para $1, \dots, p-1$ y lo probamos para $p$. El razonamiento es similar al anterior.

  Tomemos $\X = \sigma_p$, y vamos a tratar de definir $T^{\sigma_p} \colon S_p(\sigma_p) \to S_{p+1}(\sigma_p \times I)$.
  De nuevo, como $T^{\sigma_p}$ debe verificar a), se tiene que $\forall \phi \in F_p(\sigma_p), T^{\sigma_p}(\phi)$ ha de satisfacer
  \[ \partial(T^{\sigma_p}(\phi)) = (i_0^{\sigma_p})_\# (\phi) - (i_1^{\sigma_p})_\# (\phi) - T^{\sigma_p}(\partial(\phi)) \]
  donde $T^{\sigma_p}(\partial(\phi))$ está bien definida por hipótesis de inducción.

  Sea $d = (i_0^{\sigma_p})_\# (\phi) - (i_1^{\sigma_p})_\# (\phi) - T^{\sigma_p}(\partial(\phi)) \in S_p(\sigma_p \times I)$.
  Entonces:
  \begin{align*}
    &\partial(\phi) = (i_0^{\sigma_p})_\# (\partial(\phi)) - (i_1^{\sigma_p})_\# (\partial(\phi)) - \partial T^{\sigma_p}(\partial(\phi)) \\
    &\stackrel{\text{h.i.}}{=} (i_0^{\sigma_p})_\#(\partial(\phi)) - (i_1^{\sigma_p})_\#(\partial(\phi)) - [ - T^{\sigma_p}(\partial^2(\phi)) +
    (i_0^{\sigma_p})_\#(\partial(\phi)) - (i_1^{\sigma_p})_\#(\partial(\phi)) ] = 0
  \end{align*}
  y por tanto d es un ciclo.

  Como $p \geq 1$ y $\sigma_p \times I$ es convexo, por el lema de Poincaré, $d$ es un borde, luego existe $e \in S_{p+1}(\sigma_p \times I)$
  tal que $\partial(e) = d$. Definimos $T^{\sigma_p}(\phi) = e$. Notemos que si bien $e$ no tiene por qué ser único, basta cond hacer una elección
  entre los existentes.

  Extendiendo por linealidad, conseguimos construir $T^{\sigma_p}$ verificando a).

  Sean ahora $\X, \phi \colon \sigma_p \to \X$ u p-símplice singular. De nuevo, al tener que verificar b), el diagrama
  \[
  \begin{tikzcd}
    S_p(\sigma_p) \arrow{r}{T^{\sigma_p}} \arrow{d}{\phi_\#} & [2em] S_{p+1}(\sigma_p \times I) \arrow{d}{(\phi \times Id)_\#} \\
    \SX{p} \arrow{r}{T^\X}                 & S_{p+1}(\X \times I)
  \end{tikzcd}
  \]
  ha de ser conmutativo.

  Si $\tau_p \colon \sigma_p \to \sigma_p$ es la identidad (que es p-símplice singular en $\sigma_p$), $\phi_\#(\tau_p) = \phi$,
  de donde definimos $T^\X(\phi) = (\phi \times Id)_\# T^{\sigma_p}(\tau_p)$. \\
  Como los $\phi$ son la base de $\SX{p}, T^\X$ se extiende por linealidad a un homomorfismo. Veamos que verifica a) y b).

  a)
  \begin{align*}
    &(\partial T^\X + T^\X \partial)(\phi) = \partial((\phi \times Id)_\# T^{\sigma_p}) + T^\X \partial(\phi_\# \tau_p) \\
    &= (\phi \times Id)_\# \partial T^\sigma_p(\tau_p) + T^\X \phi_\#(\partial(\tau_p)) \\
    &= (\phi \times Id)_\# (- T^{\sigma_p} \partial(\tau_p) + (i_0^{\sigma_p})_\# (\tau_p) - (i_1^{\sigma_p})_\# (\tau_p)) + (\phi \times I)_\# T^{\sigma_p}\partial(\tau_p) \\
    &= (\phi \times Id)_\# (i_0^{\sigma_p})_\#(\tau_p) - (\phi \times Id)_\# (i_1^{\sigma_p})_\#(\tau_p) \\
    &= (i_0^\X)_\# \phi_\#(\tau_p) - (i_1^\X)_\# \phi_\#(\tau_p) = ((i_0^\X)_\# - (i_1^\X)_\#) (\phi).
  \end{align*}

  b) Sea $h\colon \X \to \Y$ continua y $\phi \colon \sigma_p \to \X$ un generador de $\SX{p}$. Entonces:
  \begin{align*}
    &T^\Y(h_\#(\phi)) = T^\Y(h \circ \phi) = ((h \circ \phi) \times Id)_\# T^{\sigma_p}(\tau_p) \\
    &= (h \times Id)_\# (\phi \times Id)_\# T^{\sigma_p} (\tau_p) = (h \times Id)_\# T^\X(\phi)
  \end{align*}
\end{proof}

\begin{corollary}
  Si $\X$ e $\Y$ son espacios homotópicamente equivalentes, entonces \[H_p(\X) = H_p(\Y) \quad \forall p \geq 0.\]
\end{corollary}

\begin{proof}
  Sea $f$ la equivalencia homotópica de $\X$ en $\Y$. Entonces existe su inversa, $g$, que verifica $g \circ f = Id_\X, f \circ g = Id_\Y$.
  Así, $(f \circ g)_* = Id_{H_p(\Y)},$ y $(g \circ f)_* = Id_{H_p(\X)}$. Como $(g \circ f)_* = g_* \circ f_*, f_*$ tiene inverso, y por tanto
  $f_*$ es un isomorfismo entre $H_p(\X)$ y $H_p(\Y)$.
\end{proof}

\section{Homología singular de un par}

Sea $A \subseteq \X$. Vamos a tratar de cuantificar la relación existente entre $H_*(A)$ y $H_*(\X)$.

Para ello, vamos a definir la homología singular del par $(\X, A)$, que será una extensión de la homología singular de un espacio, en el sentido
$H_*(\X, \emptyset) = H_*(\X)$.

Consideremos $i \colon A \xhookrightarrow{} \X$ la inclusión, $i_\# \colon S_p(A) \to \SX{p}$ el homomorfismo inducido.
Entonces $i\#(\sum\limits_{j = 1}^n n_j \phi_j) = \sum\limits_{j = 1}^n n_j i_\#(\phi_j) = \sum\limits_{j = 1}^n n_j (i \circ \phi_j) = \sum\limits_{j = 1}^n n_j \phi_j$,
luego $i_\#$ es un monomorfismo.

Definimos $S_p(\X, A) = \frac{\SX{p}}{i_\#(S_p(A))} = \frac{\SX{p}}{S_p(A)} = \{\overline{c} \mid c \in \SX{p} \}$ con $\overline{c} = c + S_p(A)$.
Este grupo se corresponde con el grupo abeliano libre generado por $\{\overline{\phi} \mid \phi \in \FX{p}\}$, ya que
$ \overline{c} = \overline{\sum_i n_i \phi_i} = \sum_i n_i \overline{\phi_i}, \quad c \in \SX{p}$.

Llamamos a $S_p(\X, A)$ el grupo de p-cadenas singulares del par $(\X, A)$. \\
Podemos definir un operador borde, que seguiremos notando $\partial$:
\begin{align*}
  \partial \colon S_p(\X, A) &\to S_{p-1}(\X, A) \\
  \partial \overline{c} &= \overline{\partial c}
\end{align*}
Este borde está bien definido, pues si $\overline{c} = \overline{d}$; entonces $d = c + a, a \in \S_p(A)$ y
$\partial d = \partial c + \partial a \implies \overline{\partial d} = \overline{\partial c} + \overline{\partial a} = \overline{\partial c}$ ya que
$\partial a \in S_{p-1}(A) \implies \overline{\partial a} = 0$.

Es claro, por la definición que hemos hecho, que el operador borde es un homomorfismo de clases y que $\partial^2 = 0$.

Con esto, hemos construído un complejo de cadenas $\{S_p(\X, A), \partial\}_{p \geq 0}$ al que llamaremos el complejo de cadenas singulares del par $(\X, A)$.

Para cada $p \geq 0$ definimos: \\
$Z_p(\X, A) = \{\overline{c} \in S_p(\X, A) \mid \partial\overline{c} = 0\}$ p-ciclos. \\
$B_p(\X, A) = \{\overline{c} \in S_p(\X, A) \mid \overline{c} = \partial\overline{d} \}$ p-bordes.

Así, dado que de nuevo $B_p(\X, A) \subseteq Z_p(\X, A)$, podemos considerar el cociente
\[H_p(\X, A) = \frac{Z_p(\X, A)}{B_p(\X, A)} \quad \forall p \geq 0 \]
que llamamos el \textbf{\underline{\smash{p-ésimo grupo de homología singular}}} del par $(\X, A)$.

Dados dos pares $(\X, A)$ y $(\Y, B)$, una aplicación continua $f \colon (\X, A) \to (\Y, B)$ es una aplicación continua
$f \colon \X \to \Y$ tal que $f(A) \subseteq B$. Llamaremos a $f$ una aplicación de pares entre $(\X, A)$ e $(\Y, B)$. \\
Si $f$ es tal aplicación, definimos:
\begin{align*}
  f_\# \colon S_p(\X, A) &\to S_p(\Y, B) \hspace{4em} \text{que está bien definida.} \\
  f_\#(\overline{c}) &= \overline{f_\#(c)}
\end{align*}

Con esto, hemos hecho conmutativo el diagrama:
\[
\begin{tikzcd}
    S_p(\X) \arrow{r}{f_\#} \arrow{d}{\pi} & S_p(\Y) \arrow{d}{\pi} \\
    S_p(\X, A) \arrow{r}{f_\#} & S_p(\Y, B)
  \end{tikzcd}
  \hspace{2em} \text{con } \pi(c) = \overline{c}.
\]
Es claro que $\partial f_\#(\overline{c}) = \partial \overline{f_\#(c)} = \overline{\partial f_\#(c)} = f_\# \partial(c) = f_\# \partial(\overline{c})$,
de donde se deduce $\partial f_\# = f_\# \partial$.

A la función $f_\# \colon S_*(\X, A) \to S_*(\Y, B)$ la llamamos la aplicación de cadenas inducida por la aplicación de pares $f$. \\
Nuevamente, se verifica que $(g \circ f)_\# = g_\# \circ f_\#$ y $Id_\# = Id$.

Definimos $f_* \colon H_p(\X, A) \to H_p(\Y, B)$ por $f_*([\overline{c}]) = [f_\#(\overline{c})]$, que está bien definida. \\
Claramente, $f_*$ es un homomorfismo de grupos, que verifica $(g \circ f)_* = g_* \circ f_*$ y $Id_* = Id$.

Así, hemos construído un funtor covariante de la categoría de pares y aplicaciones entre pares en la de grupos graduados y homomorfismos,
el funtor de la homología singular.
\[
\begin{tikzcd}
  (\X, A) \arrow{r} \arrow{d}{f} & [3em] \{H_p(\X, A)\}_{p \geq 0} \arrow{d}{f_*} \\
  (\Y, B) \arrow{r} & \{H_p(\Y, B)\}_{p \geq 0}
\end{tikzcd}
\]
Tenemos un resultado de obtención directa:

\begin{corollary}
  Si $h \colon (\X, A) \to (\Y, B)$ es un homeomorfismo de $\X$ sobre $\Y$ que verifica $h(A) = B$, entonces
  $h_* \colon H_p(\X, A) \to H_p(\Y, B)$ es un isomorfismo de grupos $\forall p \geq 0$.
\end{corollary}

\begin{proof}
  Basta considerar el inverso de $h$, que induce una aplicación inversa en la homología.
\end{proof}

Vamos a trasladar los resultados que vimos de arcoconexión para el caso de pares.

\begin{proposition}
  Si $\X$ es arcoconexo y $A \neq \emptyset$, se tiene $H_0(\X, A) = 0$.
\end{proposition}

\begin{proof}
  Sea $[\overline{c}] \in H_0(\X, A)$. $\partial\overline{c} = 0$, con $c = \sum\limits_{i = 1}^n n_i x_i \quad x_i \in \X$. \\
  $\overline{c} = \sum\limits_{i = 1}^n n_i \overline{x_i}$ y así, tenemos $[\overline{c}] = \sum\limits_{i = 1}^n n_i [\overline{x_i}]$. \\
  Veamos que $\forall x \in \X, [\overline{x}] = 0$.

  Como $A \neq \emptyset$, sea $a \in A$. Por arcoconexión, tomamos $\alpha \colon \sigma_1 \to \X$ verificando $\alpha(0, 1) = x, \alpha(1, 0) = a$.
  Se cumple que $\overline{\alpha} \in S_1(\X, A)$ y $\partial\overline{\alpha} = \overline{\partial\alpha} = \overline{x - a} = \overline{x} - \overline{a}
  = \overline{x}$, pues $a \in S_0(A) \implies \overline{a} = 0$. \\
  De esta forma, $[\overline{x}] = [\partial\overline{a}] = 0$, ya que es un borde.
\end{proof}

\begin{proposition}
  Sea $X = \bigcup\limits_{\alpha \in A} \X_\alpha$ la descomposición de $\X$ en componentes arcoconexas, y sea $A \subseteq \X$. Entonces se verifica:
  \begin{itemize}
    \item[a)] $\forall \alpha, i_{\alpha *} \colon H_p(\X_\alpha, \X_\alpha \cap A) \to H_p(\X, A)$ es un monomorfismo.
    \item[b)] $H_p(\X, A) = \bigoplus\limits_{\alpha \in A}  i_{\alpha *}(H_p(\X_\alpha, \X_\alpha \cap A))$.
  \end{itemize}
\end{proposition}

La demostración es similar al caso visto para espacios simples. No la repetimos, pues no aporta nada nuevo.

Tenemos también un resultado equivalente al ya visto para aplicación homotópicas.

\begin{proposition}
  Sean $f, g \colon (\X, A) \to (\Y, B)$ aplicaciones de pares. Si existe $H \colon \X \times I \to \Y$ continua tal que
  $H(\_, 0) = f, H(\_, 1) = g$ y $H(A, t) \subseteq B \hspace{0.2em} \forall t \in I$, entonces $g_* = f_*$. En tal caso, se dice que
  f y g son homotópicas.
\end{proposition}

La demostración es similar a la ya vista, por lo que no se incluye.

Consideremos ahora la proyección $\pi \colon \SX{p} \to S_p(\X, A) \cong \frac{\SX{p}}{i_*(S_p(A))}$ \\
Tenemos que $\pi(c) = \overline{c}$, y la aplicación de cadenas inducida por la inclusión,
$i_\# \colon S_p(A) \to \SX{p}$, que es claramente un monomorfismo. \\
$\pi$ es un epimorfismo, y se verifica que $\Img i_\# \cong \Ker \pi = i_\#(S_p(A))$ por la definición de $S_p(\X, A)$.
Esto nos permite definir la sucesión:
\[0 \to S_p(A) \xrightarrow{i_\#} \SX{p} \xrightarrow{\pi} S_p(\X, A) \to 0 \]
que verifica $\partial i_\# = i_\# \partial$ y $\partial \pi = \pi \partial$. \\
A estas sucesiones, en las que cada triple $C \xrightarrow{f} D \xrightarrow{g} E$ verifica $\Img f = \Ker g$, las llamaremos exactas.
Si además tienen solo tres elementos, como en el caso anterior, las llamaremos sucesiones exactas cortas.

Podemos extender el resultado anterior, y obtenemos una sucesión exacta corta de complejos de cadenas:
\[0 \to S_*(A) \xrightarrow{i_\#} S_*(\X) \xrightarrow{\pi} S_*(\X, A) \to 0 \]

Sea ahora $\pi_* \colon H_*(\X) \to H_*(\X, A)$ el homomorfismo inducido por $\pi$. \\
$\pi_*([c]) = [\overline{c}]$ que está bien definido, pues si $[c] = [d] \implies d = c + \partial h \implies \overline{d} = \overline{c} + \partial \overline{h}
\implies [\overline{c}] = [\overline{d}]$.

Veamos que $\Img i_* = \Ker \pi_*$.

Se verifica $\pi_* i_* ([a]) = \pi_*[i_\#(a)] = [\pi i_\#(a)] = 0$, pues $\pi i_\# = 0$. Así, $\Img i_* \subseteq \Ker \pi_*$.

Si $[c] \in \Ker \pi_*, \pi_*([c]) = [\overline{c}] = 0 \implies \overline{c} = \overline{\partial d}, c = \partial d + i_\#(a)$,
luego $\partial c = \partial^2 d + \partial i_\#(a) \implies \partial c = \partial^2 d + i_\# \partial a$, y se deduce que $\partial a = 0$.
$a$ es un ciclo y genera $[a]$. Como $c = \partial d + i_\#(a), [c] = [\partial d] + [i_\#(a)] = i_*[a] \implies \Ker \pi_* \subseteq \Img i_*$,
lo que demuestra la igualdad.

Generalmente $i_*$ no tiene por qué ser un monomorfismo, y $\pi_*$ no tiene por qué ser un epimorfismo, aunque se verifica el siguiente resultado.

\begin{proposition}
  $\pi_* \colon H_0(\X) \to H_0(\X, A)$ es un epimorfismo.
\end{proposition}

\begin{proof}
  Si tomamos $[\overline{c}] \in H_0(\X, A)$ entonces $c \in S_0(\X) = Z_0(\X)$, y por tanto $[\overline{c}] \in H_0(\X)$ y
  $\pi_*([c]) = [\overline{c}]$.
\end{proof}

\textit{Para p mayor que 1 no podemos asegurar que sea un ciclo, por lo que el resultado solo vale para el caso visto.}

\begin{proposition}
  Para todo $p \geq 1$ existe un homomorfismo $\Delta \colon H_p(\X, A) \to H_{p-1}(A)$ tal que $\Ker \Delta = \Img \pi_*$ y $\Img \Delta = \Ker i_*$

  A la sucesión exacta así construída:
  \[\dots H_{p+1}(\X, A) \xrightarrow{\Delta} H_p(A) \xrightarrow{i_*} H_p(\X) \xrightarrow{\pi_*} H_p(\X, A) \xrightarrow{\Delta} H_{p-1}(A) \dots\]
  se le llama la \textbf{\underline{\smash{sucesión exacta}}} del par $(\X, A)$.
\end{proposition}

\begin{proof}
  En primer lugar, vamos a definir $\Delta$.

  Sea $[\overline{c}] \in H_p(\X, A)$, con $\overline{c} \in S_p(\X, A)$ y $\partial \overline{c} = 0$. Como $\partial \overline{c} = \overline{\partial c} = 0,
  \partial c \in i_\#(S_{p-1}(A))$. Por la inyectividad de $i_\#$, existe un único $a \in S_{p-1}(A)$ tal que $\partial c = i_\#(a)$. \\
  Además, $i_\# \partial a = \partial i_\#(a) = \partial^2 c = 0$, y al ser $i_\#$ monomorfismo, $\partial a = 0$ luego $a \in Z_{p-1}(A), [a] \in H_{p-1}(A)$.

  Definimos pues $\Delta([\overline{c}]) = [a]$, que está bien definido, ya que:
  \begin{align*}
    [\overline{c}] = [\overline{d}] &\implies \overline{d} = \overline{c} + \partial \overline{h} \implies d = c + \partial h + i_\#(e) \\
    &\implies \partial d = \partial c + i_\#(\partial e) = i_\#(a) + i_\#(\partial e) = i_\#(a + \partial e) \\
    &\implies \Delta([\overline{d}]) = [a + \partial e] = [a] = \Delta([\overline{c}])
  \end{align*}

  Veamos que $\Delta$ es homomorfismo, es decir, $\Delta([\overline{c_1}] + [\overline{c_2}]) = \Delta([\overline{c_1 + c_2}])$.

  $\Delta([\overline{c_1}]) = [a_1]$ con $\partial c_1 = i_\#(a_1)$ \\
  $\Delta([\overline{c_2}]) = [a_2]$ con $\partial c_2 = i_\#(a_2)$

  De esta forma, $\partial(c_1 + c_2) = i_\#(a_1 + a_2)$ y $\Delta([\overline{c_1 + c_2}]) = [a_1 + a_2]$.

  Veamos que $\Img \Delta = \Ker i_*$.

  $i_*(\Delta([\overline{c}])) = i_*[a] = [i_\#(a)] = [\partial c] = 0$, de donde $\Img \Delta \subseteq \Ker i_*$

  Si $[a] \in H_{p-1}(A)$ verificando $i_*(a) = 0$, entonces $[i_\#(a)] = 0$, y eso nos dice que $i_\#(a) = \partial b, b \in S_p{\X}$.
  Además, $\overline{\partial b} = \overline{i_\#(a)} = 0$, y como $\overline{\partial b} = \partial \overline{b}, b \in Z_p(\X, A)$ con
  $\partial b = i_\#(a)$. Por definición, $\Delta([\overline{b}]) = [a]$. Así se da la igualdad.

  Veamos que $\Ker \Delta = \Img \pi_*$.

  $\Delta\pi_*([c]) = \Delta([\overline{c}]) = [a]$, con $\partial c = i_\#(a)$. \\
  Como $c$ es un ciclo, $i_\#(a) = 0$, y al ser $i_\#$ un monomorfismo, $a = 0$, luego $[a] = 0$. Así, $\Img \pi_* \subseteq \Ker \Delta$.

  Si $[\overline{c}] \in \Ker \Delta, \Delta[\overline{c}] = [\overline{a}] = 0$, con $\partial c = i_\#(a)$. Como $[a] = 0, a = \partial b, b \in \S_p(A)$.
  Así, $\partial(c - i_\#(b)) = 0$, y $c - i_\#(b) \in Z_p(\X)$, de donde obtenemos $[c - i_\#(b)]$. \\
  $\pi_*([c . i_\#(b)]) = \pi_*([c]) + e = [\overline{c}]$. Con esto, $\Img \pi_* = \Ker \Delta$.
\end{proof}

\begin{proposition}[Naturalidad de la sucesión exacta]
  Sea $f \colon (\X, A) \to (\Y, B)$ una aplicación de pares. Entonces el siguiente diagrama es conmutativo:
  \[  \begin{tikzcd}
    \dots H_p(A) \arrow{r}{i_*} \arrow{d}{(f_{|A})_*} & H_p(\X) \arrow{r}{\pi_*} \arrow{d}{f_*} & H_p(\X, A) \arrow{r}{\Delta} \arrow{d}{f_*} & H_{p-1}(A) \dots \arrow{d}{(f_{|A})_*} \\
    \dots H_p(B) \arrow{r}{i_*} & H_p(Y) \arrow{r}{\pi_*} & H_p(\Y, B) \arrow{r}{\Delta} & H_{p-1}(B) \dots
  \end{tikzcd} \]
\end{proposition}

\begin{proof}
  Nos basta con probar que $\Delta \circ f_* = f_{|A} \circ \Delta$, pues ya conocemos el resto.

  Realizamos el cálculo:
  \[ (\Delta \circ f_*)([c]) = \Delta([f_\#(\overline{c}]) = \Delta([\overline{f_\#(c)}]) \stackrel{(*)}{=} [(f_{|A})_\#(a)]
    = (f_{|A})_*[a] = (f_{|A})_* \Delta[\overline{c}] \]

  Donde en $(*)$ estamos usando que $\partial(f_\#(c)) = f_\# \partial c = f_\# i_\#(a) = i_\#(f_{|A})_\#(a)$, con $\partial c = i_\#(a).$
\end{proof}


%************************************************
\chapter{Cálculo de la homología singular}\label{ch:introduccion}
%************************************************

En esta parte, vamos a enunciar y demostrar el teorema de subdivisión baricéntrica,
el cual nos permite calcular directa o indirectamente (a través de sus consecuencias)
el grupo de homología de ciertos espacios topológicos, como haremos con las esferas
n-dimensionales.

\section{Teorema de subdivisión baricéntrica}

\begin{definition}
  Sea $\X$ un espacio topológico. Diremos que $\U \subseteq \mathcal P(\X)$ es un recubrimiento de $\X$ si $\X = \bigcup\limits_{U \in \U} U$.
\end{definition}

Tomemos ahora $\X$ un espacio topológico y $\U$ un recubrimiento suyo. Para $ p \geq 0$, podemos definir
\[ F_p^\U(\X) = \{\phi \in F_p(\X) \mid \exists U \in \U \text{ tal que } \Img \phi \subseteq U\} \]
que es, por definición, subconjunto de $F_p(\X)$.

Si $S_p^\U(\X)$ es el grupo abeliano libre generado por $F_p^\U(\X)$, por $F_p^\U(\X) \subseteq F_p(\X)$, se tiene
que $S_p^\U(\X)$ es un subgrupo de $S_p(\X)$. Si consideramos la inclusión $i \colon S_p^\U(\X) \to S_p(\X)$, si
$\phi \in F_p^\U(\X)$, entonces $\partial_i(\phi) \in F_{p-1}^\U(\X)$ y así, $\partial(\phi) \in S_{p-1}^\U(\X)$. \\
En consecuencia, la restricción del borde a $S_p^\U(\X)$ define un homomorfismo borde, y así, construímos un complejo
de cadenas $S_*^\U(\X) = \{S_p^\U(\X), \partial\}_{p \geq 0}$.

Este complejo es un subcomplejo de $S_*(\X)$. Si tomamos la inclusión $i \colon S_*^\U(\X) \to S_*(\X)$, induce un
homomorfismo en la homología, $i_* \colon H_p(S_*^\U(\X)) \to H_p(\X) \hspace{0.2em} \forall p \geq 0$, donde
$H_p(S_*^\U(\X)) = \frac{\Ker \partial}{\Img \partial}$

Sea ahora $\Y$ otro espacio topológico, y $\V$ un recubrimiento suyo. Diremos que una función continua $f \colon \X \to \Y$ es compatible
con $\U, \V$ si $\forall U \in \U \exists V \in \V$ tal que $f(U) \subseteq V$. En tal caso, puede definirse un homomorfismo
$f_\# \colon S_p^\U(\X) \to S_p^\V(\Y)$ dado por la restricción de $f_\# \colon \SX{p} \to \SY{p} a S_p^\U(\X)$.

Esto hace que el siguiente diagrama sea conmutativo:
\[ \begin{tikzcd}
  S_p^\U(\X) \arrow{r}{f_\#} \arrow{d}{i} & S_p^\V(\Y) \arrow{d}{i} \\
  \SX{p} \arrow{r}{f_\#} & \SY{p}
\end{tikzcd} \]

\begin{remark}
  El diagrama conmutativo está bien definido, pues si $\phi \in F_p^\U(\X)$, entonces existe $U \in \U$ con $\Img \phi \subseteq U$, de donde
  $f_\#(\phi) = f \circ \phi$ verifica $\Img(f \circ \phi) \subseteq f(U) \subseteq V \in \V$.
\end{remark}

Como $f_\# \partial = \partial f_\#$, $f$ induce $f_* \colon H_p(S_p^\U(\X)) \to H_p(S_p^\V(\Y))$.

\begin{theorem}[Teorema de subdivisión baricéntrica]
  Sean $\X$ un espacio topológico y $\U$ un recubrimiento de $\X$ tal que $\interior{\U} = \{\interior{U} \mid U \in \U\}$
  también recubra a $\X$. Entonces la aplicación
  \[ i_* \colon H_p(S_*^\U(\X)) \to H_p(\X) \]
  es un isomorfismo, $\forall p \geq 0$.

  Además, si $(\Y, \V)$ es otro par de las mismas características, y $f \colon \X \to \Y$ es una función continua
  compatible con $\U, \V$, entonces el siguiente diagrama es conmutativo:
  \[ \begin{tikzcd}
    H_p(S_*^\U(\X)) \arrow{r}{i_*} \arrow{d}{f_*} & H_p(\X) \arrow{d}{f_*} \\
    H_p(S_*^\V(\Y)) \arrow{r}{i_*} & H_p(\Y)
  \end{tikzcd} \]

\end{theorem}

Antes de pasar a la demostración del teorema, necesitaremos introducir nuevos conceptos. La idea es definir la subdivisión
de símplices singulares mediante la subdivisión baricéntrica de símplices, que nos llevan a poder dar una especie de
inversa homotópica de la inclusión.

\section{Subdivisión de símplices}

\begin{definition}
  Sea $s_p = (a_0, \dots, a_p)$ un p-símplice. Definimos el baricentro de $s_p$ como
  \[b(s_p) = \frac{a_0 + \dots + a_p}{p+1}\]
\end{definition}

Los puntos $\{b(s_p), a_0, \dots, a_{i-1}, a_{i+1}, \dots, a_p\}$ genera un p-símplice $\forall i = 0, \dots, p$.
Definimos la subdivisión baricéntrica del símplice $s_p$ por inducción sobre $p$, y la notamos $Sd(s_p)$.
\begin{itemize}
  \item $Sd(s_0) = s_0$
  \item $Sd(s_p) = (b(s_p), t_0, \dots, t_{p-1}) \mid (t_0, \dots, t_{p-1}) \in Sd(\U_{p-1}), \U_{p-1} \in \interior{S_p}$
        siendo $\interior{S_p}$ el conjunto de caras de orden $p-1$ de $S_p$.
\end{itemize}

\begin{remark}
  Es importante destacar que el orden en el que se ponen los puntos es importante, pues si el baricentro fuera en otro lugar,
  la orientación del símplice cambiaría, y afectaría al cálculo del borde.
\end{remark}

Si $K$ es un conjunto de símplices, la subdivisión de $K$ vienen dada por $Sd(K) = \bigcup\limits_{s_p \in K} Sd(s_p)$.

Se define la n-ésima subdivisión como \[Sd^n(s_p) = Sd(Sd^{n-1}(s_p)), \text{ con } Sd^1 = Sd \]

\begin{proposition}
  Sea $t_p \in Sd(s_p)$ un p-símplice de la subdivisión. Entonces $diam(t_p) \leq \frac{p}{p+1} diam(s_p)$.
\end{proposition}

\begin{proof}
  Haremos la demostración por inducción sobre $p$.

  Para $p = 0$, se tiene $0 = 0$.

  Supuesto cierto hasta $p-1$, lo probamos para $p$.

  Sea $s_p = (a_0, \dots, a_p)$. Se verifica que $diam(s_p) = \max\limits_{i,j} |a_i - a_j|$.
  Como $t_p \in Sd(s_p), t_p = (b(s_p), t_0, \dots, t_{p-1}) \in Sd(\U_{p-1})$.
  Entonces, $diam(t_p) = max\{|b(s_p) - t_i|,\hspace{0.2em} |t_i - t_j|, i,j = 0, \dots, p-1\}$.

  Por h. de inducción sabemos que $|t_i - t_j| \leq diam(t_0, \dots, t_{p-1}) \leq \frac{p-1}{p} diam(\U_{p-1})$.
  Como $U_{p-1} \subseteq s_p$, entonces $diam(U_{p-1}) \leq diam(s_p)$, y como $\frac{p-1}{p} \leq \frac{p}{p+1}$
  se tiene $|t_i - t_j| \leq \frac{p}{p+1} diam(s_p)$

  Además, $|b(s_p) - t_i| \leq \max\limits_i |b(s_p) - a_i| = \max\limits_i |\frac{a_0 + \dots + a_p}{p-1} - a_i| =
  \max\limits_i | \frac{\sum\limits_{i \neq j} |a_j - a_i|}{p+1} \leq \frac{p}{p+1} diam(s_p)$
\end{proof}

\begin{definition}
  Sea $\varphi$ una colección de subconjuntos acotados de un espacio euclídeo. Definimos $malla \varphi = \sup\limits{c \in \varphi} diam(c)$
\end{definition}

\begin{corollary}
  Si $K$ es una colección de p-símplices, se verifica \[malla(Sd(K)) \leq \frac{p}{p+1} malla(K) \]
\end{corollary}

\begin{corollary}
  Si $K$ es una colección finita de p-símplices, entonces $\forall \epsilon > 0, \exists n \in \N$ tal que \\$malla(Sd^n(k)) < \epsilon$
\end{corollary}

\begin{proof}
  Simplemente aplicamos $\lim\limits_{n \to \infty}(\frac{p}{p+1})^n = 0$
\end{proof}

\section{Subdivisión baricéntrica de símplices singulares afines}

Vamos a tratar ahora la subdivisión de símplices singulares afines. Estos nos permitirán aproximar la
subdivisión de cualquier símplice singular, por lo que su estudio nos resulta de gran interés.

Sea $C$ un subconjunto convexo del espacio euclídoe en el que trabajamos. Para todo $p \geq 0$, sea
\[F_p^A(C) = \{ \phi \in F_p(C) \mid \phi \text{ es una aplicación afín}\} \]
es decir, $\phi \colon \sigma_p \to C \quad$ verifica \[\phi(ta + (1-t)b) = t\phi(a) + (1-t)\phi(b) \hspace{1em} \forall t \in [0,1]\]

Sea $A_p(C)$ el grupo abeliano libre generado por $F_p^A(C)$, subgrupo libre de $S_p(C)$. Si $\phi \in F_p^A(C), \partial_i \phi = \phi \circ F_p^i$,
y como $\phi$ es afín y $F_p^i$ también, $\partial_i \phi$ es afín. Así:
\[\partial \phi = \sum\limits_{i = 0}^p (-1)^1 \partial_i \phi \hspace{0.5em} \in A_{p-1}(C)\]

Es consecuencia, podemos restringir el borde a $A_p(C)$, haciendo conmutativo
\[ \begin{tikzcd}
      A_p(C) \arrow{r}{\partial} \arrow{d}{i} & A_{p-1}(C) \arrow{d}{i} \\
      S_p(C) \arrow{r}{\partial} & S_{p-1}(C)
\end{tikzcd} \]

Tenemos pues un complejo de cadenas $A_*(C) = \{A_p(C), \partial\}_{p \geq 0}$, para todo convexo $C$.

Si $\phi \in F_p^A(C), \phi \colon \sigma_p \to C$ afín. Notando $\sigma_p = (v_0, \dots, v_p)$, entonces
${x_i = \phi(v_i) \in C}$ determinan unívocamente a $\phi$, pues si $v = \sum\limits_{i = 0}^p \alpha_i v_i$, entonces
$\phi(v) = \sum\limits_{i = 0}^p \alpha_i \phi(v_i)$ por afinidad. Por tanto, todo elemento de $\phi \in F_p^A(C)$ se
identifica con $p+1$ puntos de $C$.
\[\phi \in F_p^A(C) \equiv \{x_0, \dots, x_p\} \subseteq C\]

Además, $\partial \phi = \sum\limits_{i = 0}^p (-1)^i \partial_i \phi$, con $\partial_i \phi = \{x_0, \dots, \hat{x_i}, \dots, x_p\}$

En consecuencia, $A_p(C)$ es el grupo abeliano libre sobre:
\[\{(x_0, \dots, x_p) \mid x_i \in C\} \equiv C \times \dots \times C \]
y el operador borde lo da $\partial(\{x_0, \dots, x_p\}) = \sum\limits_{i = 0}^p (-1)^i \{x_0, \dots, \hat{x_i}, \dots, x_p\}$

Si tenemos $C, C'$ convexos, y $f \colon C \to C'$ es una aplicación afín (y por tanto continua), al tomar la restricción de $f_\#$
a $A_p(C)$, se tiene que si ${\phi \in F_p^A(C)}$ entonces $f_\#(\phi) = {f \circ \phi \in F_p^A(C')} \implies {f_\#(A_p(C)) \subseteq A_p(C')}$

En consecuencia, $\partial f_\# = f_\# \partial$ y $f_\# \colon A_*(C) \to A_*(C')$ es una aplicación de cadenas. Además,
${f_\#(\{x_0, \dots, x_p\}) = \{f(x_0), \dots, f(x_p)\}}$

\begin{definition}
  Sea $x \in C$ Definimos
  \begin{align*}
    \C_x \colon A_p(C) &\to A_{p+1}(C) \\
    \{x_0, \dots, x_p\} &\mapsto \{x, x_0, \dots, x_p\}
  \end{align*}
  y la extendemos por linealidad a un homomorfismo, pues tenemos los valores de los generadores.
\end{definition}

\begin{proposition}
  Veamos que:
  \begin{itemize}
    \item[a)] $\partial \C_x + \C_x \partial = Id_{A_p(C)} \quad \forall p \geq 1$
    \item[b)] $f_\# \C_x = \C_{f(x)} f_\# \quad \forall f \colon C \to C'$ afín.
  \end{itemize}
\end{proposition}

\begin{proof}
  a) $\partial \C_x(x_0, \dots, x_p) = \partial(x, x_0, \dots, x_p) = \{x_0, \dots, x_p\} + \\ \sum\limits_{i = 1}^{p+1} (-1)^i \partial_i (x, x_0, \dots, x_p)
  = \{x_0, \dots, x_p\} + \sum\limits_{j = 0}^{p} (-1)^{j+1} \partial_{j+1} (x, x_0, \dots, x_p)
  = \{x_0, \dots, x_p\} - \sum\limits_{j = 0}^{p} (-1)^j \C_x(x_0, \dots, \hat{x_j}, \dots, x_p) = \{x_0, \dots, x_p\} - \C_x \partial(x_0, \dots, x_p)$

  b) $f_\# \C_x(x_0, \dots, x_p) = f_\#(x, x_0, \dots, x_p) = (f(x), f(x_0), \dots, f(x_p)) = \C_{f(x)} f_\#(x_0, \dots, x_p)$
\end{proof}

Ahora podemos definir una subdivisión para símplices o cadenas afines. Lo haremos por inducción, definiendo ${Sd' \colon A_p(C) \to A_p(C)}$

$\fbox{p = 0}$
\[Sd' = Id_{A_p(C)}\]

$\fbox{p > 0}$, supuesto definido hasta $p-1$
\[ Sd'(\phi) = \C_{b(\phi)}(Sd' \circ \partial \phi) \]
donde $b(\phi) = \frac{x_0 + \dots + x_p}{p+1}$, siendo $\phi = \{x_0, \dots, x_p\}$ y extendido linealmente.

Veamos algunas propiedades. Debido a que la definición se hizo por inducción, todas las demostraciones serán por inducción.

a) $\partial \circ Sd' = Sd' \circ \partial$

\begin{proof}
  Para $p = 0$ la demostración es trivial, por la definición de $Sd'$. Suponemos $p \geq 1$.

  Sea $\phi \in F_p^A(C)$, entonces
  \begin{align*}
    \partial \circ Sd'(\phi) &= \partial \C_{b(\phi)} Sd' \partial \phi = Sd'(\partial \phi) - \C_{b(\phi)} \partial Sd'(\partial \phi) \\
    &= Sd'(\partial \phi) - \C_{b(\phi)} Sd'(\partial^2 \phi) = Sd'(\partial \phi)
  \end{align*}
\end{proof}

b) Si $f \colon C \to C'$ es afín, entonces $f_\# \circ Sd' = Sd' \circ f_\#$

\begin{proof}
  $(f_\# Sd')(\phi) = f_\#(\C_{b(\phi)} Sd'(\partial(\phi))) = \C_{f(b(\phi))} f_\# Sd'(\partial(\phi))\\ \stackrel{\text{h. de i}}{=}
  \C_{f(b(\phi))} Sd'(f_\#(\partial(\phi))) = \C_{f(b(\phi))} Sd'(\partial(f_\#(\phi))) = Sd'(f_\#(\phi))$
\end{proof}

c) Existe un homomorfismo $\tau' \colon A_p(C) \to A_{p+1}(C)$ tal que ${\partial \tau' + \tau' \partial = Sd' - Id_{A_p(C)}}$,
es decir, $Sd'$ es, como aplicación de cadenas, homotópica a la identidad.

\begin{proof}
  Para $p = 0$, tomamos $\tau' = 0$

  Supuesto cierto hasta $p-1$, buscamos $\tau' \colon A_p(C) \to A_{p+1}(C)$ verificando
  \[\partial \tau' = Sd' - I_{A_p(C)} - \tau' \partial\]

  Como el segundo miembro es conocido, entonces si $\phi \in F_p^A(C), \tau'(\phi)$ ha de verificar
  \[\partial(\tau'(\phi)) = Sd'(\phi) - \phi - \tau'(\partial(\phi))\]
  lo calculamos
  \begin{align*}
    &\partial(Sd'(\phi) - \phi - \tau'(\partial(\phi)) = Sd'(\partial(\phi)) - \partial(\phi) - \partial(\tau' \partial(\phi)) \\
    &= Sd'(\partial(\phi)) - \partial(\phi) + \tau'(\partial^2(\phi)) - Sd'(\partial(\phi)) + \partial(\phi) = 0
  \end{align*}
  Definiendo $\tau'(\phi) = \C_{b(\phi)}(Sd'(\phi) - \phi - \tau'(\partial(\phi))$, entonces
  \begin{align*}
    &\partial(\tau'(\phi)) = \partial \C_{b(\phi)}(Sd'(\phi) - \phi -\tau'(\partial(\phi))) = (Sd'(\phi) - \phi - \tau' \partial(\phi)) \\
    &- \C_{b(\phi)} \partial(Sd'(\phi) - \phi - \tau'(\partial(\phi))) = Sd'(\phi) - \phi - \tau'(\partial(\phi))
  \end{align*}
\end{proof}

d) $\tau'$ es natural, esto es, si $f \colon C \to C'$ es afín, entonces $f_\# \tau' = \tau' f_\#$

\begin{proof}
  Para $p = 0$ es trivial por definición.

  Suponemos cierto hasta $p-1$. Sea $\phi \in F_p^A(C)$, entonces
  \begin{align*}
    &f_\#(\tau'(\phi)) = f_\#(\C_{b(\phi)}(Sd'(\phi) - \phi - \tau'(\partial(\phi))) \\
    &= \C_{f(b(\phi))} f_\#(Sd'(\phi) - \phi - \tau(\partial(\phi))) \\
    &= \C_{f(b(\phi))}(Sd'(f_\# \phi) - f_\# \phi - \tau'(f_\# \partial(\phi))) = \tau'(f_\#(\phi))
  \end{align*}
\end{proof}

e) Propiedad técnica. Nos permitirá subdividir símplices hasta que estos queden dentro del recubrimiento.

Sea $c = \sigma_p$ y $\tau_p \colon \sigma_p \to \sigma_p$ la identidad ($\tau_p \in F_p^A(\sigma_p)$). \\
Si $Sd'(\tau_p) = \sum_i n_i \phi_i$ (que ocurre siempre), con $\phi_i \colon \sigma_p \to \sigma_p$ afines, entonces
$\forall i \hspace{0.3em} \exists t_p^i \in Sd(\sigma_p)$ tal que $\Img(\phi_i) \subseteq t_p^i$.

\begin{proof}
  Para $p = 0$, $Sd'(\tau_0) = \tau_0 = \sigma_0$. Siendo $t_0^i = \sigma_0 = Sd'(\sigma_0)$.

  Supuesto cierto hasta $p-1$, sea $Sd'(\tau_p) = \C_{b(\tau_p)} Sd'(\partial(\tau_p))$
  \begin{align*}
    &\partial(\tau_p) = \sum\limits_{i = 0}^p (-1)^i \partial_i(\tau_p) = \sum\limits_{i = 0}^p (-1)^i \tau_p \circ F_p^i = (por ser la identidad) \\
    &= \sum\limits_{i = 0}^p (-1)^i F_p^i \circ \tau_{p-1} = \sum\limits_{i = 0}^p (-1)^i(F_p^i)_\# (\tau_{p-1}) \text{. De esta forma:} \\
    &Sd'(\tau_p) = \C_{b(\tau_p)}(\sum\limits_{i = 0}^p (-1)^i Sd'(F_p^i)_\# (\tau_{p-1})) = \C_{b(\tau_p)}(\sum\limits_{i = 0}^p (-1)^i(F_p^i)_\# Sd'(\tau_{p-1})) \\
  \end{align*}
  Si $Sd'(\tau_{p-1}) = \sum\limits_j m_j \psi_j$, entonces, por inducción, $\forall j$ existe $t_{p-1}^j \in Sd(\sigma_{p-1})$ tal que
  $\Img(\psi_j) \subseteq t_{p-1}^j$.

  Por tanto, $Sd'(\tau_p) = \sum\limits_{i = 0}^p \sum\limits_j (-1)^i m_j C_{b(\tau_p)}(F_p^i)_\# (\psi_j)$

  Como $\Img(\psi_j) \subseteq t_{p-1}^j \in Sd(\sigma_{p-1})$
  $\Img((F_p^i)_\# (\psi_j)) = \Img(F_p^i \circ \psi_j) \subseteq F_p^i(\Img(\psi_j)) \subseteq F_p^i(t_{p-1}^j) \in Sd(\interior{\sigma_p})$,
  luego $\Img((F_p^i)_\# (\psi_j)) \subseteq U_{p-1}^{i,j} \in Sd(\interior{\sigma_p})$.

  Así, $\Img(\C_{b(\tau_p)}(F_p^i)_\# (\phi_j)) \subseteq \C_{b(\tau_p)}(U_{p-1}^{i, j}) = (b(\sigma_p), t_0, \dots, t_{p-1}) \in Sd(\sigma_p)$.\\
  (Donde $U_{p-1}^{i,j} = (t_0, \dots, t_p)$)
\end{proof}

\section{Subdivisión de símplices singulares}

% Introducción

Nos basaremos en el caso afín.

\begin{proposition}
  Sea $\X$ un espacio topológico. Entonces, para todo $p \geq 0$ existe
  \[Sd\colon S_p{\X} \to S_p{\X} \] homomorfismo, verificando:
  \begin{itemize}
    \item[a)] Si $\X$ es un convexo de un espacio euclídeo, $\frac{Sd}{A_p(C)} = Sd'$
    \item[b)] $Sd \circ \partial = \partial \circ Sd$
    \item[c)] $Sd \circ f_\# = f_\# \circ Sd \quad \forall f \colon \X \to \Y$ continua
    \item[d)] Existe un homomorfismo $\tau \colon \SX{p} \to S_{p+1}(\X)$ tal que ${\partial \tau + \tau \partial = Sd - Id_{S_p(\X)}}$
              $\forall p \geq 0$ y $\tau f_\# = f_\# \tau \quad \forall f \colon \X \to \Y$ continua
    \item[e)] Si $\phi \in F_p(\X)$ y $Sd(\phi) = \sum_i n_i \phi_i$, entonces $\forall i \exists t_p^i \in Sd(\sigma_p)$ tal que
              $\Img(\phi_i) \subseteq \phi(t_p^i)$

  \end{itemize}
\end{proposition}

\begin{proof}
  Sea $\phi \in F_p(\X)$. Entonces, si $\tau_p \colon \sigma_p \to sigma_p$ es el p-símplice identidad, se tiene que
  $\phi = \phi_\#(\tau_p)$. En consecuencia, $Sd(\phi)$ debe verificar $Sd(\phi_\#)(\tau_p) = \phi_\#(Sd(\tau_p)) = \phi_\# Sd'(\tau_p)$. \\
  Definimos $Sd(\phi) = \phi_\#(Sd'(\tau_p))$ y lo extendemos por linealidad.

  Veamos ahora las propiedades.

  a) Si $\X$ es convexo, $\phi \in F_p^A(\X)$, y entonces, $Sd(\phi) = \phi_\# Sd'(\tau_p) = (Sd' \circ \phi_\#)(\tau_p) = Sd'(\phi)$
  pues ya lo hemos probado en este caso.

  b) \begin{align*}
    &(\partial \circ Sd)(\phi) = \partial \phi_\# Sd'(\tau_p) = \phi_\# \partial Sd'(\tau_p) = \phi_\# Sd'(\partial \tau_p) \\
    &= \phi_\#(Sd'(\sum\limits_{i = 0}^p (-1)^i (F_p^i)_\# (\tau_{p-1}))) = \phi_\#(\sum\limits{i = 0}^p (-1)^i Sd' \circ (F_p^i)_\#(\tau_{p-1}))\\
    &= \text{(caso afín)} = \phi_\#(\sum\limits_{i = 0}^p (-1)^i (F_p^i)_\# Sd'(\tau_{p-1})) = \sum\limits_{i = 0}^p (-1)^i (\phi \circ F_p^i)_\# Sd'(\tau_{p-1}) \\
    &= (\partial_i \phi)_\# (\sum\limits_{i = 0}^p (-1)^i Sd'(\tau_{p-1})) = \sum\limits_{i = 0}^p (-1)^i Sd(\partial_i \phi) = Sd(\sum\limits_{i = 0}^p (-1)^i \partial_i \phi) \\
    &= (Sd \circ \partial)(\phi)
  \end{align*}

  c) $(f_\# \circ Sd)(\phi) = f_\# \phi_\# Sd'(\tau_p) = (f \circ \phi)_\# Sd'(\tau_p) = Sd(f_\#(\phi))$

  d) Definimos \[\tau(\phi) = \phi_\#(\tau'(\tau_p)) \quad \text{con} \quad \phi_\# \colon S_{p+1}(\sigma_p) \to S_{p+1}(\X) \]
  Entonces, se tiene:
  \begin{align*}
    &(\partial \tau + \tau \partial)(\phi) \partial(\phi_\# \tau'(\tau_p)) + \tau(\sum\limits_{i = 0}^p (-1)^i \partial_i \phi) \\
    &= \phi_\#(\partial \tau'(\tau_p)) + \sum\limits_{i = 0}^p (-1)^i((\partial_i \phi)_\#(\tau'(\tau_{p-1})) = \phi_\#(\partial \tau'(\tau_p)) \\
    &+ \phi_\#(\sum\limits_{i = 0}^p (-1)^i (F_p^i)_\# (\tau'(\tau_{p-1}))) = \text{caso afín} = \phi_\#(\partial \tau'(\tau_p)) \\
    &+ \phi_\#(\tau'(\sum\limits_{i = 0}^p (-1)^i (F_p^i)_\#(\tau_{p-1})) = \phi_\#(\partial \tau'(\tau_p)) + \phi_\#(\tau'\partial(\tau_p)) \\
    &= \phi_\#(\partial \tau' + \tau' \partial)(\tau_p) = \text{caso afín} = \phi_\#(Sd' - Id)(\tau_p) = Sd(\phi) - \phi
  \end{align*}

  e) Como $Sd(\phi) = \phi_\#(Sd'(\tau_p))$, entonces si $Sd'(\tau_p) = \sum_i n_i \psi_i, \phi_i = \phi_\#(\psi_i)$

  Como existe $t_p^i \in Sd(\sigma_p)$ tal que $\Img(\psi_i) \subseteq t_p^i$, entonces
  \[\Img(\phi_i) = \Img(\phi_\#(\psi_i)) = \phi(\Img(\psi_i)) \subseteq \phi(t_p^i)\]

  \begin{remark}
    Probar que si $\phi \in F_p(\X)$ y $Sd^n(\phi) = \sum_i n_i \phi_i \implies \forall i \exists t_p^i \in Sd^n(\sigma_p)$ tal que
    $\Img(\phi_i) \subseteq \phi(t_p^i)$, consiste en reiterar la propiedad e) ya demostrada.
  \end{remark}

\end{proof}

\section{Demostración del teorema}

Demostremos ya el teorema de subdivisión baricéntrica. Recordamos su enunciado.

\begin{theorem}[Teorema de subdivisión baricéntrica]
  Sean $\X$ un espacio topológico y $\U$ un recubrimiento de $\X$ tal que $\interior{\U} = \{\interior{U} \mid U \in \U\}$
  también recubra a $\X$. Entonces la aplicación
  \[ i_* \colon H_p(S_*^\U(\X)) \to H_p(\X) \]
  es un isomorfismo, $\forall p \geq 0$.

  Además, si $(\Y, \V)$ es otro par de las mismas características, y $f \colon \X \to \Y$ es una función continua
  compatible con $\U, \V$, entonces el diagrama que obtenemos considerando $f_* e i_*$ es conmutativo.
\end{theorem}

\begin{proof}
  Sea $\X$ un espacio topológico, y sea $\U$ un recubrimiento de $\X$ tal que $\interior{\U}$ tambień recubre a $\X$.

  Si $\phi \in F_p(\X)$, esto es, $\phi \colon \sigma_p \to \X$ es continua, sea $V = \{\phi^{-1}(U) \mid U \in \U\}$, recubrimiento
  abierto de $\sigma_p$ como espacio métrico compacto. Sea $\epsilon(\phi)$ el número de Lebesgue de dicho recubrimiento.
  Así, si $A \subseteq \sigma_p$ y $diam(A) < \epsilon(\phi)$ entonces existe $U \in \U$ tal que $A \subseteq \phi^{-1}(U)$ y entonces
  $\phi(A) \subseteq U$

  En consecuencia, $\forall \phi \in F_p(\X) \quad \exists \epsilon(\phi) > 0$ tal que si $A \subseteq \sigma_p, diam(A) < \epsilon(\phi)$,
  entonces existe $U \in \U$ tal que $\phi(A) \subseteq(U)$

  Por un corolario que ya hemos visto, dado $\epsilon(\phi)$, existe $n(\phi)$ tal que \\
  ${malla(Sd^{n(\phi)}(\sigma_p)) < \epsilon(\phi)}$. Si $t_p \in Sd^{n(\phi)}(\sigma_p), diam(t_p) < \epsilon(\phi)$,
  luego existe $U \in \U$ con $\phi(t_p) \subseteq U$. En consecuencia, si $Sd^{n(\phi)}(\phi) = \sum_i n_i \phi_i$,
  entonces $\Img(\phi_i) \subseteq \phi(t_p^i)$ con $t_p^i \in Sd^{n(\phi)}(\sigma_p)$,
  por lo que existirá $U_i \in \U$ con $\Img(\phi_i) \subseteq \phi(t_p^i) \subseteq U_i$ \\
  Así, $Sd^{n(\phi)}(\phi) \in S_p^\U(\X)$

  Sea $m(\phi) = \inf\{n(\phi) \mid Sd^{n(\phi)}(\phi) \in S_p^\U(\X)\}$, que puede ser $0$, si, por ejemplo, $\phi \in S_p^\U(\X)$

  Definimos $\psi \colon S_p(\X) \to S_p^\U(\X)$ por
  \[\psi(\phi) = Sd^{m(\phi)}(\phi) - \sum\limits_{i = 0}^p (-1)^i \tau(Sd^{m(\partial_i(\phi))} + \dots + Sd^{m(\phi) - 1}) \partial_i \phi \]
  y si $m(\partial_i \phi) = m (\phi)$, el sumando $i$ no aparece en la expresión de $\psi$.

  Hay que ver que está bien definido, para lo que comprobamos:
  \begin{itemize}
    \item[i)] $m(\partial_i(\phi)) \leq m(\phi)$
    \item[ii)] $\tau(S_{p-1}^\U(\X)) \subseteq S_p^\U(\X)$
  \end{itemize}

  i) $Sd^n(\partial_i(\phi)) = \sum_\alpha n_\alpha \psi_\alpha$, con $\Img(\psi_\alpha \subseteq (\partial_i \phi)(t_{p-1}^\alpha), t_{p-1}^\alpha \in Sd^n(\sigma_{p-1})$. \\
  Como $\partial_i \phi = \phi \circ F_p^i$, se tiene:
  \[ \Img(\phi_\alpha) \subseteq (\phi \circ F_p^i)(t_{p-1}^i) \subseteq \phi(\U_{p-1}^\alpha) \subseteq \phi((b(\sigma_p), \U_ {p-1}^\alpha)) = \phi(S_p^\alpha)\]
  con $S_p^\alpha \in Sd^n(\sigma_p)$

  Si $Sd^n(\phi) \in S_p^\U(\X)$, entonces $\phi(S_p^\alpha) \subseteq U_\alpha \in \U$, luego se verifica
  \[ \Img(\psi_\alpha) \subseteq U_\alpha \in \U \implies Sd^n(\partial_i \phi) \in S_{p-1}^\U(\X) \]
  De esta forma, $M(\partial_i \phi) = \inf\{n(\phi) \mid Sd^n(\partial_i \phi) \in S_{p-1}\U(\X) \leq m(\phi)$

  ii) Como $\tau(\phi) = \phi_\# \tau'(\tau_p))$, sea $\phi \colon \sigma_{p-1} \to \X$ continua con $\Img(\phi) \subseteq U \in \U$, esto es,
  un generador de $S_{p-1}^\U(\X)$. Entonces
  \[ \tau(\phi) = \phi_\#(\tau'(\tau_p)) = \phi_\#(\sum_i n_i \phi_i) = \sum_i n_i \phi_\#(\phi_i) \]
  Como $\Img(\phi_\#(\phi_)) \subseteq \Img(\phi) \subseteq U \in \U, \tau(\phi) \in S_p^\U(\X)$

  Una vez probada la buena definición de $\psi$, definimos:
  \begin{align*}
    T \colon S_p(\X) &\to \SX{p+1} \quad \text{ por}
    T(\phi) &= \tau(Sd^0 + Sd + Sd^2 + \dots + Sd^{m(\phi) - 1}
  \end{align*}
  Vamos a ver que:
  \begin{itemize}
    \item[a)] $\psi \circ i = Id_{S_p^\U(\X)}$
    \item[b)] $i \circ \psi - Id_{\SX{p}} = \partial T - T \partial$
  \end{itemize}
  En tal caso, $i \circ \psi$ y la identidad serían aplicaciones de cadenas homotópicas, e inducirían los mismos homomorfismos en la homología.
  Esto, junto a que $\psi \circ i = Id$, nos dice que $i_* \colon H_p^{\U(\X)} \to H_p(\X)$ es un isomorfismo $\hspace{0.2em} \forall p$, con inverso $\psi_*$.
  Concluímos la demostración probando a) y b).

  a) $(\psi \circ i)(\phi) = \psi(\phi)$ con $\phi \in F_p^\U(\X)$. Como $m(\phi) = 0$ en este caso, $\psi(\phi) = \phi - 0 = \phi$

  b) Sabemos que $\partial \tau + \tau \partial = Sd - Id$, por lo que \\
  \[ \partial \tau Sd - \tau Sd \partial = \partial \tau Sd - \tau \partial Sd = Sd^2 - Sd \]
  Repitiendo este proceso sucesivamente, se tiene
  \[ \partial \tau Sd^{m - 1} - \tau Sd^{m-1} \partial = Sd^m - Sd^{m-1} \]
  Si los sumamos todos, $\partial \tau(I  Sd + \dots + Sd^{m-1}) - \tau(I  Sd + \dots + Sd^{m-1}) \partial = Sd^m - I$ \\
  De esta forma
  \begin{align*}
    &(\partial T + T \partial)(\phi) = \partial \tau(I  Sd + \dots + Sd^{m-1}) \phi + T(\sum\limits_{i = 0}^p (-1)^i \partial_i(\phi)) \\
    &= \partial \tau(I  Sd + \dots + Sd^{m(\phi)-1}) \phi + \sum\limits_{i = 0}^p (-1)^i \tau(I + \dots + Sd^{m(\partial_i(\phi)) - 1} \partial_i \phi \\
    &= Sd^{m(\phi)}(\phi) - \phi - \tau(I + Sd + \dots + Sd^{m(\phi)-1})(\partial \phi)  \\
    &+ \sum\limits_{i = 0}^p (-1)^i \tau(I + \dots + Sd^{m(\partial_i(\phi)) - 1} \partial_i \phi \\
    &= Sd^{m(\phi)}(\phi) - \phi - \sum\limits_{i = 0}^p (-1)^i \tau(I + \dots + Sd^{m(\phi) - 1}) \partial_i \phi \\
    &+ \sum\limits_{i = 0}^p (-1)^i \tau(I + \dots + Sd^{m(\partial_i(\phi)) - 1}) \partial_i \phi \\
    &= Sd^{m(\phi)}(\phi) - \phi - \sum\limits_{i = 0}^p (-1)^i \tau(Sd^{m(\partial_i(\phi))} + \dots + Sd^{m(\phi) - 1}) \partial_i \phi = \psi(\phi) - \phi
  \end{align*}
  de donde se sigue b).

  Para finalizar, veremos que la hipótesis de que $\interior{\U}$ sea recubrimiento es esencial.

  Sea $\X = S^1, \U = \{\{N\}, S^1 - \{N\}\}$ donde $N$ representa el polo norte, $(0, 1)$. Es claro que $\U$ recubre a $\X$, pero no su interior.

  Veamos que $i_* \colon H_1(S_*^\U(S^1)) \to H_1(S^1)$ no es un isomorfismo. \\
  Sea $[c] \in H_1(S_*^\U(S^1)), c \in S_1^\U(\S^1), \partial c = 0$. Entonces, $c = c_1 + c_2$, y si $\phi \colon \sigma_1 \to \{N\}$ es la función
  (única) continua, entonces, $c_1 = n\phi, n \in \Z$. Así, $\partial = \partial c = n \partial \phi + \partial c_2 = \partial c_2$, con $c_2$ una
  1-cadena en $S^1 - \{N\} \equiv \R$, que es un ciclo.

  Como $H_1(\R) = 0$, existe $d_2$, una 2-cadena en $S^1 - \{N\}$ tal que $\partial d_2 = c_2$. Así, $d_2$ es también cadena en $S_2^\U(S^1)$, y se tiene
  $c = n\phi + \partial d_2 = n \partial \overline{\phi} + \partial d_2 = \partial(n\overline{\phi}  d_2) \implies [c] = 0$ por ser un borde.

  Por tanto, $H_1(S_*^\U(S^1)) = 0$ y $H_1(S^1) \cong \Z$, que no son isomorfos.
\end{proof}

\section{Consecuencias del teorema de subdivisión baricéntrica}

Vamos a tratar ahora algunas de las principales consecuencias del teorema de subdivisión baricéntrica.

\begin{theorem}[Sucesión de Mayer-Vietoris]
  Sea $\X$ un espacio topológico y $\{U, V\}$ un recubrimiento de $\X$ tal que $\interior{U} \cup \interior{V} = \X$. Consideramos las inclusiones
  \begin{align*}
    & i \colon U \cap V \to U &\quad k \colon U \to \X \\
    & j \colon U \cap V \to V &\quad l \colon V \to \X
  \end{align*}
  En este caso, existe un homomorfismo $\Delta \colon H_p(\X) \to H_{p-1}(U \cap V)$ tal que la siguiente sucesión es exacta para todo $p \geq 1$:
  \[ \dots \to H_p(U \cap V) \xrightarrow{G} H_p(U) \oplus H_p(V) \xrightarrow{H} H_p(\X) \xrightarrow{\Delta} H_{p-1}(U \cap V) \to \dots \]
  donde $G([c]) = (i_*[c], j_*[c]), \quad H([c_1], [c_2]) = k_*([c_1]) - l_*([c_2])$

  Además, esta construcción es natural, pues si $(\Y, \{R, S\})$ está en las condiciones del teorema, y $f \colon \X \to \Y$ es continua con
  $f(U) \subseteq R, f(V) \subseteq S$, entonces el diagrama resultante es conmutativo.
  \[ \begin{tikzcd}
    H_p(U \cap V) \arrow{r}{G} \arrow{d}{f*} & H_p(U) \oplus H_p(V) \arrow{r}{H} \arrow{d}{f_*  \oplus f_*} & H_p(\X) \arrow{r}{\Delta} \arrow{d}{f_*} & H_{p-1}(U \cap V) \arrow{d}{f_*} \\
    H_p(R \cap S) \arrow{r}  & H_p(R) \oplus H_p(S) \arrow{r} & H_p(\Y) \arrow{r} & H_{p-1}(R \cap S)
  \end{tikzcd} \]

\end{theorem}

\begin{proof}
  Consideramos la siguiente sucesión:
  \begin{align*}
    0 \to S_p(U \cap V) &\to S_p(U) \oplus S_p(V) \to S_p^U(\X) \to 0 \\
                    c &\mapsto (i_\#(c), j_\#(c)) =  (c_1, c_2) \mapsto k_\#(c_1) - l_\#(c_2)
  \end{align*}

  Veamos que es exacta.

  Claramente, $(i_\#, j_\#)$ es inyectiva, pues las aplicaciones $i_\#, j_\#$ lo son. \\
  Si $d \in S_p^U(\X)$, podemos expresar $d$ como $d = d_1 + d_2$, con $d_1 = \sum_i n_i \phi_i, d_2 = \sum_j m_j \psi_j,
  \Img \phi_i \subseteq U, \Img \psi_j \subseteq V$. De esta forma, $d_1 = k_\#(d_1), -d_2 = l_\#(d_2)$, y $d = k_\#(d_1) - l_\#(-d_2)$.
  Es claro que $k_\#(i_\#(c)) - l_\#(j_\#(c)) = 0$, y, por último, si $k_\#(c_1) - l_\#(c_2) = 0$

\end{proof}


%************************************************
\chapter{Algunos resultados clásicos}\label{ch:resultados}
%************************************************

\section{Teorema del punto fijo de Brower}

Vamos a utilizar el cálculo de la homología de las esferas para obtener resultados
de carácter topológico.

\begin{proposition}
  a) Esferas de distinta dimensión no tienen el mismo tipo de homotopía. En particular, no son homeomorfas.

  b) Espacios euclídeos de distinta dimensión no son homeomorfos.
\end{proposition}

\begin{proof}
  a) Sean $S^n, S^m$ con $n \neq m$. Si $S^n \cong S^m$, entonces $H_n(S^n) \cong H_n(S^m)$, pero hemos visto que eso no sucede.

  b) Si $n \neq m$, y $h \colon \R^n \to \R^m$ es un homeomorfismo, entonces su restricción a $R^n - \{0\}$ también lo es.
  Como $R^n - \{0\} \cong S^n$, y $R^m - \{f(0)\} \cong S^m$, llegamos de nuevo a una contradicción.
\end{proof}

\begin{remark}
  $\R^n$ y $\R^m$ tienen el mismo tipo de homotopía. Si $n \leq m$, podemos ver $\R^m$ como \\
  $\R^m = \{(x, y) \mid x \in \R^n, y \in \R^{m-n}\}$, y se define
  \[ H((x, y), t) = (1-t)(x, y) + t(x, 0) \]
  homotopía entre ambos espacios, siendo $\R^n$ un retracto de $\R^m$.
\end{remark}

\begin{lemma}
  No hay ninguna retracción de $D^n$ a $S^{n-1}$.
\end{lemma}

\begin{proof}
  Sean $D^n = \{x \in \R^n \mid |x| \leq 1\}, S^{n-1} = \{x \in \R^n \mid |x| = 1\} \subseteq D^n$. Si existiera dicha retracción,
  $r \colon D^n \to S^{n-1}$, tenemos el siguiente diagrama conmutativo:
  \[ \begin{tikzcd}
    S^{n-1} \arrow{dr}{i} \arrow{rr}{Id} & & S^{n-1} \\
    & D^n \arrow{ur}{r}
  \end{tikzcd} \]
  de donde se obtiene un diagrama, también conmutativo, en la homología:
  \[ \begin{tikzcd}
    Z \cong H_{n-1}(S^{n-1}) \arrow{dr}{i_*} \arrow{rr}{Id_*} & & H_{n-1}(S^{n-1}) \\
    & H_{n-1}(D^n) \arrow{ur}{r_*}
  \end{tikzcd} \]
  lo cual es imposible, pues $H_{n-1}(D^n) \cong 0$, y tendríamos una factorización de la identidad que incluye a la aplicación nula,
  que en este caso sería $i_*$.
\end{proof}

\begin{corollary}
  Sea $f \colon D^n \to \R^n$ continua, tal que la restricción de $f$ a $S^{n-1}$ sea la identidad.
  Entonces existe $x \in D^n$ tal que $f(x) = 0$.
\end{corollary}

\begin{proof}
  Supongamos que no existe dicho $x$. En tal caso, es posible definir $h(x) = \frac{f(x)}{|f(x)|}$
  continua, que contradice el resultado anterior.
\end{proof}

Vamos a demostrar un teorema que nos dará como consecuencia directa un resultado clásico, el teorema
del punto fijo de Brower.

\begin{theorem}
  Sea $f \colon D^n \to \R^n$ continua. Entonces o bien existe $y \in D^n$ tal que $f(y) = 0$, o existe
  $z \in S^{n-1}$ con $f(z) = \lambda z, \quad \lambda < 0$.
\end{theorem}

\begin{proof}
  Sea $h \colon D^n \to \R^n$, definida por:
  \[ h(x) = \begin{cases} f(2x) &\text{ si } |x| \leq \frac{1}{2} \\
                          (2-2|x|)f(\frac{x}{|x|}) + (2|x| + 1)x &\text{ si } |x| \geq \frac{1}{2}  \end{cases} \]
  que es continua, ya que en $\frac{1}{2}$ ambas definiciones coinciden.

  Si $|x| = 1, h(x) = x$, y aplicando el corolario anterior, existe $y \in D^n$ tal que $h(y) = 0$. Distinguimos dos casos:

  Si $|y| \leq \frac{1}{2}, 0 = f(2y)$, y estamos en el primer caso.

  Si $|y| > \frac{1}{2}, f(\frac{x}{|x|}) = \frac{(-2|x| + 1)x}{2-2|x|}$. Si llamamos $z = \frac{x}{|x|}$,
  $f(z) = \frac{(-2|x|+1)|x|}{2-2|x|} z$, por lo que estamos en el segundo caso del teorema.
\end{proof}

\begin{corollary}
  Sea $f \colon D^n \to \R^n$ continua. Entonces existe $y \in D^n$ tal que $f(y) = y$ o existe $z \in S^{n-1}$ con $f(z) = \mu z, \mu > 1$.
\end{corollary}

\begin{proof}
  Aplicamos el teorema anterior a $-f + Id$.
\end{proof}

\begin{corollary}[Teorema del punto fijo de Brower]
  Toda aplicación continua de $D^n$ en sí mismo tiene un punto fijo.
\end{corollary}

\begin{proof}
  Basta ver que no puede darse la segunda opción del corolario anterior, pues la imagen de la función caería fuera de $D^n$.
\end{proof}

\section{Teoremas de invarianza}

Definamos el grupo de homología local, herramienta que nos servirá para demostrar varios resultados de invarianza.

\begin{definition}
  Sean $\X$ un espacio topológico, $x \in \X$. A $H_p(\X, \X - \{x\})$ lo llamaremos el grupo de homología local de $\X$ en $x$.
\end{definition}

\begin{proposition}
  Sea $x \in \X$, y supongamos que $\{x\}$ es un cerrado. Entonces, para todo entorno $V$ de $x$, se verifica:
  \[ H_p(\X, \X - \{x\}) \cong H_p(V, V - \{x\}) \quad \forall p \geq 0 \]
\end{proposition}

\begin{proof}
  Sea $U = \X - V \subseteq \X - \{x\}$. Si tomamos el cierre, se tiene
  \[ \overline{U} = \overline{\X - V} = \X - \interior{V} \subseteq \X - x = \X - \overline{x} = \interior{(\X - x)} .\]
  Aplicando el teorema de escisión, $i \colon (\X - U, (\X - \{x\}) - U) \to (\X, \X - \{x\})$ induce un isomorfismo en la homología.
  Como $\X - U = V, (\X - \{x\}) - U = V - \{x\}$, obtenemos el resultado.
\end{proof}

\begin{proposition}
  Sean $\X, \Y$ espacios topológicos, $x \in \X, y \in \Y$ cerrados. Si existen entornos $x \in V, y \in W$ tales que $(V, \{x\}) \cong (W, \{y\})$,
  entonces $H_*(\X, \X - \{x\}) \cong H_*(\Y, \Y - \{y\})$.
\end{proposition}

\begin{proof}
  Sea $h \colon V \to W$ el homeomorfismo entre $V$ y $W$, con $h(x) = y$. La restricción $h \colon V - \{x\} \to W - \{y\}$ también es un homeomorfismo,
  por lo que $H_*(V, V - \{x\}) \cong H_*(W, W - \{y\})$: Basta aplicar la proposición anterior para obtener el resultado.
\end{proof}

Vamos a calcular la homología local de $\R^n$.

\begin{proposition}
  $H_p(\R^n, \R^n - \{x\}) \cong \begin{cases} 0 &\text{ si } n \neq p \\ \Z &\text{ si } n = p \end{cases}$
\end{proposition}

\begin{proof}
  Si $p = 0$, ya sabemos que $H_0(\R^n, \R^n - \{x\}) = 0$. Supondremos $p \geq 1$.

  Consideramos la sucesión de homología del par, teniendo en cuenta que $\R^n - \{x\}$ y $S^{n-1}$ tienen el mismo tipo de homotopía. \\
  Si $p = 1$:
  \[ 0 = H_1(\R^n) \xrightarrow{\pi_*} H_1(\R^n, \R^n - \{x\}) \xrightarrow{\Delta} H_0(\R^n - \{x\}) \xrightarrow{i_*} H_0(\R^n) = \Z \to 0 \]
  Si $n \geq 2, H_0(\R^n - \{x\}) \cong \Z$. \\
  Como $\Ker \Delta = \Img \pi_* = 0, \Img \Delta = \Ker i_* = 0, H_1(\R^n, \R^n - \{x\}) = 0$. \\
  Si $n = 1, H_0(\R^n - \{x\}) \cong \Z \oplus \Z$, luego $H_1(\R^n, \R^n - \{x\}) \cong \Z$.

  Si $p \geq 2$:
  \[ 0 = H_p(\R^n) \xrightarrow{\pi_*} H_p(\R^n, \R^n - \{x\}) \xrightarrow{\Delta} H_{p-1}(\R^n - \{x\}) \xrightarrow{i_*} H_{p-1}(\R^n) = 0 \]
  Si $p \neq n, H_{p-1}(\R^n - \{x\} = 0$, luego $H_p(\R^n, \R^n - \{x\}) = 0$. \\
  Si $p = n, H_{p-1}(\R^n - \{x\} \cong \Z$, luego  $H_p(\R^n, \R^n - \{x\}) = \Z$.
\end{proof}

\begin{theorem}[Invarianza de la dimensión]
  Sean $O_1 \subseteq \R^n, O_2 \subseteq \R^m$ abiertos homeomorfos. Entonces $n = m$.
\end{theorem}

\begin{proof}
  Sea $\phi \colon O_1 \to O_2$ el homeomorfismo entre los abiertos, y sea $x \in O_1$. Como $O_1, O_2$ son entornos de $x, \phi(x)$,
  por la proposición anterior:
  \[H_*(\R^n, \R^n - \{x\}) \cong H_*(\R^m, \R^m - \{\phi(x)\}) \]
  y en particular
  \[ \Z \cong H_n(\R^n, \R^n - \{x\}) \cong H_n(\R^m, \R^m - \{\phi(x)\}) \]
  de donde se deduce $n = m$.
\end{proof}

\begin{corollary}
  Si $M$ es un espacio topológico $T_2$, conexo, y para cada $p \in M$ existen $V_p$ entorno de $p$, un entero
  $n(p) \geq 1$ y un homeomorfismo $\phi_p \colon V_p \to O_p \subseteq \R^{n(p)}$ con $O_p$ abierto, entonces
  $n(p) = n$ constante y $M$ es una variedad topológica $n$-dimensional.
\end{corollary}

\begin{proof}
  Sea $p_0 \in M$ fijo, $A = \{p \in M \mid n(p) = n(p_0)\}$ no vacío, pues $p_0 \in A$. Es claro que si $q \in V_p$,
  $n(q) = n(p)$, pues si no fuera así, se contradice el teorema de invarianza de la dimensión. Así, si $p \in A, V_p \subseteq A$
  y por tanto $A$ es un abierto.

  Si $p \in \overline{A}$, existe $q \in A \cap V_p$ con $n(q) = n(p_0)$ por $q \in A$, y $n(q) = n(p)$ por $q \in V_p$.
  Por tanto, $n(p) = n(p_0), p \in A$, y $A$ es cerrado.

  En consecuencia, $A = M$.
\end{proof}

Pasemos ahora a un resultado que nos permitirá demostrar el teorema de invarianza del borde.

\begin{lemma}
  Sean $H^n = \{x \in \R^n \mid x_n \geq 0 \}$ semiespacio de $\R^n$, $\partial H^n = Fr H^n = \{x \in \R^n \mid x_n = 0\}$.
  Consideremos también $\interior{H^n} = \{x \in \R^n \mid x_n > 0\}$. Entonces, se verifica
  \[H_p(H^n, H^n - \{x\}) =
    \begin{cases} \Z &\text{ si } p = n, x \in \interior{H^n} \\
                  0  &\text{ si } p \neq n, x \in \interior{H^n} \\
                  0  &\text{ si } x \in \partial H^n  \end{cases} \]
\end{lemma}

\begin{proof}
  Si $x \in \interior{H^n}$, por una proposición anterior, sabemos que
  $H_p(H^n, H^n - \{x\}) \cong H_p(\interior{H^n}, \interior{H^n} - \{x\}) \cong H_p(\R^n, \R^n - \{x\})$, pues
  $\interior{H^n} \cong \R^n$. Con esto obtenemos los valores para $x \in \interior{H^n}$.

  Si $x \in \partial H^n, H^n - \{x\}$ es estrellado desde un punto. Usando esto y la sucesión de homología del par $(H^n, H^n - \{x\})$,
  se obtiene la última parte.
\end{proof}

\begin{corollary}[Invarianza del borde]
  Sean $O_1 \subseteq \R^n, O_2 \subseteq \R^m$ abiertos homeomorfos. Entonces
  \begin{itemize}
    \item [a)] Si $\phi \colon O_1 \to O_2$ es el homeomorfismo, $\phi(O_1 \cap \partial H^n) = O_2 \cap \partial H^m$.
    \item [b)] $n = m$.
  \end{itemize}
\end{corollary}

\begin{proof}
  Si $x \in O_1 \cap \partial H^n$, por $x \in O_1$, se tiene $H_*(H^n, H^n - \{x\}) = 0$. \\
  $H_*(H^n, H^n - \{x\}) \cong H_*(O_1, O_1 - \{x\}) \cong H_*(O_2, O_2 - \{\phi(x)\}) \cong H_*(H^m, H^m - \{\phi(x)\})$.
  Así, $x \in O_1 \cap \partial H^n \iff \phi(x) \in O_2 \cap \partial H^m$, de donde se deduce a).

  Para demostrar b), solo hemos de aplicar a) para obtener que $\phi(O_1 \cap \interior{H^n}) = O_2 \cap \interior{H^m}$, y como
  $\interior{H^n}, \interior{H^m}$ son abiertos en $\R^n, \R^m$ respectivamente, el teorema de invarianza de la dimensión
  nos dice que $n = m$.
\end{proof}

\begin{definition}
  Una variedad topológica con borde es un espacio topológico Hausdorff $M$ verificando que $\hspace{0.2em} \forall p \in M$ existen
  $(V_p, \phi_p), n(p) \geq 1$ con $V_p$ entorno abierto de $p$, $\phi_p \colon V_p \to O_p \subseteq H^{n(p)}$ homeomorfismo
  sobre un abierto $O_p$ de $H^{n(p)}$.
\end{definition}

\begin{proposition}[Consecuencias]
  Se verifican los siguientes resultados sobre variedades topológicas con borde:
  \begin{itemize}
    \item[a)] En cada componente conexa de $M, n(p)$ es constante, y la llamaremos la dimensión de dicha componente.

    \item[b)] Si $\partial M = \{p \in M \mid H_*(M, M - \{p\}) = 0\}$, entonces
    \[ \partial M = \{p \in M \mid \forall (V_p, \phi_p), \phi_p(p) \in \partial H^n \cap O_p\} \]

    \item[c)] Si $\interior{M} = M - \partial M$, entonces $\interior{M}$ tiene estructura de variedad topológica
    (sin borde) $n$-dimensional, y $\partial M$ tiene estructura de variedad topológica con borde $(n-1)$-dimensional.
    Dicha estructura es la única que convierte a las inclusiones $\interior{M} \xhookrightarrow{} M, \partial M \xhookrightarrow{} M$
    en embebimientos. Además, $\partial M$ es cerrado de $M$ y, por tanto, $\interior{M}$ es abierto de $M$.
  \end{itemize}
\end{proposition}

\begin{proof}
  a) es trivial por una proposición ya vista.

  b) Si $(V_p, \phi_p)$ es una carta en $p$, entonces
  \[ H_*(M, M - \{p\}) \cong H_*(V_p, V_P - \{p\}) \cong H_*(O_p, O_p - \{\phi_p(p)\}) \cong H_*(H^n, H^n - \{\phi_p(p)\}) \]
  que vale $0$ sí y solo sí $\phi_p(p) \in \partial H^n$.

  c) Si $\phi \colon V \to O \subseteq H^n$ es una carta de M, entonces $\phi(V \cap \partial M) = O \cap \partial H^n$, y se tiene que
  $\phi(V \cap \interior{M}) = O \cap \interior{H^n}$, abierto en $O$, luego $V \cap \interior{M}$ es abierto en $V$, y por tanto en $M$.
  Así, $\interior{M}$ es abierto en $M$.

  Las estructuras se definen por:
  \begin{align*}
    \interior{M} &\colon \{(V \cap \interior{M}, \phi_{|V \cap \interior{M}}) \mid (V, \phi) \text{ es una carta en }M\}. \\
    \partial M &\colon   \{(V \cap \partial M,  \phi_{|V \cap \partial M}), \mid (V, \phi) \text{ es una carta en }M\}.
  \end{align*}

  El resto es trivial.
\end{proof}

Vamos a definir el concepto de variedad diferenciable.

\begin{definition}
  Una variedad diferenciable con borde $n$-dimensional es una variedad topológica con borde $n$-dimensional en la que
  los cambios de coordenadas son diferenciables.
\end{definition}

\begin{remark}
  Podemos extender la proposición anterior al caso de variedades diferenciables. Para ello, únicamente hemos de extender
  c) al caso diferenciable.
\end{remark}

\section{Separación de Jordan-Brower. Invarianza del dominio}

\begin{proposition}
  Sea $A \subseteq S^n$ homeomorfo a $I^k$, con $0 \leq k \leq n$. Entonces:
  \[ H_p(S^n - A) \cong \begin{cases} 0 &\text{ si } p > 0 \\ \Z &\text{ si } p = 0 \end{cases} \]
\end{proposition}

\begin{proof}
  Lo demostraremos por inducción sobre $k$.

  Si $k = 0, I^k = \{p\}$, luego $S^n - A \cong \R^n$, y se tiene el resultado.

  Supongamos el resultado cierto hasta $k-1$, y probémoslo para $k$.

  Sea $\phi \colon I^k \to A$ el homeomorfismo, y definamos $I_+^k = \{x \in I^k \mid x_1 \geq \frac{1}{2}\}$, \\
  $I_-^k = \{x \in I^k \mid x_1 \leq \frac{1}{2}\}$. Sean $A^+, A^-$ las imágenes por $\phi$ de $I_+^k, I_-^k$,
  respectivamente. En $S^n - (A^+ \cap A^-)$ tomamos la descomposición $U = S^n - A^+, V = S^n - A^-$. Así,
  $U \cap V = S^n - A$, y por Mayer-Vietoris, se tiene:
  \[ H_p(S^n - (A^+ \cap A^-)) \to H_{p-1}(S^n - A) \xrightarrow{(i_*^+, i_*^-)} H_{p-1}(S^n - A^+) \oplus H_{p-1}(S^n - A^-) \dots \]

  Como $A^+ \cap A^- \cong I^{k-1}$, por hipótesis de inducción \\
  \[H_p(S^n - (A^+ \cap A^-)) \cong \begin{cases} 0 &\text{ si } p > 0 \\
                                                \Z &\text{ si } p = 0 \end{cases}\]
  Así, $(i_*^+, i_*^-)$ es un isomorfismo si $p > 0$, y un monomorfismo si $p = 0$.

  Supongamos ahora que el resultado no es cierto. Entonces, para $p > 0$ existe $[\alpha] \in H_p(S^n - A)$ con $[\alpha] \neq 0$, y
  si $p = 0$, existen $x, y \in S^n - A$ tales que $[x-y] \neq 0$.

  Para tratar ambos casos a la vez llamamos $\beta$ a $(x-y)$ o $\alpha$, indistintamente.

  Por lo que sabemos de $(i_*^+, i_*^-)$ se sigue que existe $A_1 \subseteq A, A_1 \cong I^k$ tal que
  $i_{1*}(\beta) \neq 0$, con $i_1 \colon S^n - A \xhookrightarrow{} S^n - A_1$.

  Nos basta con tomar $A_1 = A^+$ o $A_1 = A^-$, dependiendo de si $i_*^+(\beta) = 0$ o $i_*^-(\beta) = 0$.

  Aplicamos el mismo razonamiento a $A_1$, con lo que obtenemos una sucesión
  \[ A \supset A_1 \supset \dots \supset A_k \supset \dots \text{ con } A_i \cong I^k \hspace{0.2em} \forall i \]
  tal que $\bigcap\limits_{i = 1}^\infty A_i \cong I^{k-1}$, y tal que si $i_m \colon S^n - A \xhookrightarrow{} S^n - A_m$
  es la inclusión, entonces $i_{m*}(\beta) \neq 0$.

  Consideramos el diagrama de inclusiones:
  \[ \begin{tikzcd}
    S^n - A \arrow{dr}{i} \arrow{rr}{i_m} & & S^n - A_m \arrow{dl}{j_m} \\
    & S^n - \bigcup\limits_{i = 1}^\infty A_i
  \end{tikzcd} \]
  Como $H_p(S^n - \bigcap\limits_{i = 1}^\infty A_i) = \begin{cases} 0 &\text{ si } p > 0 \\ \Z &\text{ si } p = 0 \end{cases}$ por inducción, entonces \\
  $0 = i_*(\beta) = 0 = [i_\#(\beta)]$, luego existe $d \in S_{p+1}(\bigcup\limits_{i = 1}^\infty A_i)$ con $\partial d = i_\#(\beta)$.

  Si $d = \sum\limits_{i = 1}^r m_i \phi_i$, entonces $\bigcup\limits_{i = 1}^r \Img \phi_i$ es un compacto en $\bigcap\limits_{i = 1}^\infty A_i
  = \bigcup\limits_{i = 1}^\infty (S^n - A_i)$, luego por el teorema de Heine-Borel:
  \[\bigcup\limits_{i = 1}^r \Img \phi_i \subseteq (S^n - A_i) \cup \dots \cup (S^n - A_{ik}) = S^n - A_m \]
  con $m = \max\{i_1, \dots, i_k\}$.

  Así, existe $e \in S_{p+1}(S^n - A_m)$ con $j_{m\#(e)} = d$. Por tanto, $j_{m\#} \circ i_{m\#} (\beta) = i_\#(\beta) = \partial d
  = \partial j_{m\#}(e) = j_{m\#}(\partial e)$.

  Como $j_{m\#}$ es un monomorfismo, $i_{m\#}(\beta) = \partial e$, entonces $i_{m\#}(\beta) = 0$, lo que contradice lo afirmado.
\end{proof}

\begin{theorem}
  Sea $A \subseteq S^n$, $n \geq 2$, $A$ homeomorfo a $S^k$ con $0 \leq k < n$. Entonces:

  Si $k \leq n-2, \quad H_p(S^n - A) = \begin{cases} 0 &\text{ si } p \neq 0, n-k-1 \\ \Z &\text{ si } p = 0, p = n-k-1 \end{cases}$

  Si $k = n-1, \quad    H_p(S^n - A) = \begin{cases} 0 &\text{ si } p > 0 \\ \Z \oplus \Z &\text{ si } p = 0 \end{cases}$
\end{theorem}

\begin{proof}
  La probamos por inducción sobre $k$.

  Si $k = 0$, $S^0$ son dos puntos, luego $A$ es homeomorfo a dos puntos, y $S^n - A \cong \R^n - \{x\} \cong S^{n-1}$.
  Se comprueba que el resultado coincide con el que comocemos para el grupo de homología de las esferas.

  Supuesto cierto hasta $k - 1$, lo probamos para $k$.

  Sea $S^k = E_+^k \cap E_-^k$, con $E_+^k = \{x \in S^k \mid x_{k+1} \geq 0\}, E_-^k = \{x \in S^k \mid x_{k+1} \leq 0\}$.
  Entonces, $E_+^k$ y $E_-^k$ son homeomorfos a $I^k$, y $E_+^k \cap E_-^k \cong S^{k-1}$.

  Del resultado anterior y de la hipótesis de inducción, conocemos $H_*(S^k, E_+^k)$, $H_*(S^k, E_-^k)$ y $H_*(S^n - (E_+^k \cap E_-^k))$.
  Aplicando Mayer-Vietoris a $S^n - (E_+^k \cap E_-^k)$, con $U = S^n - E_+^k, V = S^n - E_-^k$, se tiene:
  \begin{align*}
    &H_p(S^n - E_+^k) \oplus H_p(S^n - E_-^k) \to H_p(S^n - (E_+^k \cap E_-^k)) \to H_{p-1}(S^n - S^k) \\
    &\to H_{p-1}(S^n - E_+^k) \oplus H_{p-1}(S^n - E_-^k)
  \end{align*}
  Si $p \geq 2$, se verifica $H_{p-1}(S^n - S^k) \cong H_p(S^n - (E_+^k \cap E_-^k)) \cong H_p(S^n - S^{k-1})$,
  que por inducción es $\begin{cases} 0 &\text{ si } p \neq 0, n-k \\ \Z &\text{ si } p = 0, n-k \end{cases}$

  Llamando $q = p-1$, si $k = n-1$, como $p \geq 2$, $H_q(S^n - A) = 0 \quad \forall q \geq 1$.\\
  Si $k \leq n-2, n-k \geq 2$, y $H_q(S^n - A) = \begin{cases} \Z &\quad q = n-k-1 \\ 0 &\quad \text{en otro caso.} \end{cases}$.

  Para $p = 1 (q = 0)$, tomamos la sucesión exacta
  \[0 \to H_1(S^n - (E_+^k \cap E_-^k)) \to H_0(S^n - S^k) \to \Z \oplus \Z \to \Z \to 0 \]
  y se deduce $H_0(S^n - S^k) \cong \Z \oplus H_1(S^n - (E_+^k \cap E_-^k)) = \begin{cases} \Z \oplus \Z &\quad k = n-1 \\
                                                                                                    \Z &\quad k < n - 1  \end{cases} $
\end{proof}

\begin{theorem}[Separación de Jordan-Brower]
  Sea $f \colon S^{n-1} \to S^n$ un embebimiento. Entonces $S^n - f(S^{n-1})$ tiene dos componentes conexas,
  cuya frontera común es $f(S^{n-1})$.
\end{theorem}

\begin{proof}
  Sea $A = f(S^{n-1})$. $H_0(S^n - A) \cong \Z \oplus \Z$, por lo que $S^n - A$ tiene dos componentes arcoconexas,
  a las que llamamos $\Omega_1$ y $\Omega_2$.

  Como $S^n - A$ es un abierto de $S^n$, es localmente arcoconexo, luego sus componentes arcoconexas y conexas coinciden,
  y por tanto $\Omega_1$ y $\Omega_2$ son las componentes conexas de $S^n - A$.

  Como $S^n - A$ es localmente arcoconexo, sus componentes conexas son abiertos en $S^n -A$, y por tanto en $S^n$.
  De esta forma, $\Omega_i \cup A = S^n - \Omega_j, i \neq j$ es un cerrado en $S^n$, de donde se obtiene
  $\overline{\Omega_i} \subseteq \Omega_i \cup A$, y por tanto
  \[ Fr \Omega_i = \overline{\Omega_i} \cap \overline{S^n - \Omega_i} \subseteq (\Omega_i \cup A) \cap (S^n - \Omega_i) =
  (\Omega_i \cup A) \cap (\Omega_j \cup A) = A. \]
  Así, $Fr \Omega_i \subseteq A$.

  Veamos que $A \subseteq Fr \Omega_i$, lo que completa la demostración.

  Sean $a \in A, V$ entonrno abierto de $a$ en $S^n$. Usando que $S^n$ es una variedad $n$-dimensional, se puede descomponer
  $A = A^1 \cup A^2$, con $a \in \interior{A^2} \subseteq V \cap A, A^1, A^2 \cong I^{n-1}$.

  Sean $p_i \in \Omega_i, p_i \in S^n - A \subseteq S^n - A^1$. Como $H_0(S^n - A^1) \cong \Z$, $S^n - A^1$ es arcoconexo,
  luego existe un arco $\gamma \colon [0,1] \to S^n - A^1$ que va desde $p_1$ a $p_2$.

  Por tanto, existe un $t \in [0,1]$ con $\gamma(t) \in A^2$, pues si no, la curva estaría en $S^n - (A^1 \cup A^2) = S^n - A$
  y uniría dos componentes conexas distintas.

  Como $\gamma^{-1}(A^2) \neg \emptyset$ y es cerrado, es un compacto de $[0,1]$, del que podemos tomar mínimo y máximo,
  a los que llamamos $t_m$ y $t_M$ respectivamente. Tomando $\epsilon$ suficientemente pequeño, $\gamma(t_m - \epsilon) \in \Omega_1 \cap V$,
  y $\gamma(t_M + \epsilon) \in \Omega_2 \cap V$, luego $a \in \overline{Omega_1} \cap \overline{Omega_2}, a \in Fr \Omega_i$.
\end{proof}

Este teorema nos permite obtener directamente un resultado que puede verse como una generalización del teorema
de la curva de Jordan a dimensiones superiores.

\begin{corollary}
  Sea $f \colon S^{n-1} \to \R^n$ un embebimiento. Entonces, $\R^n - f(S^{n-1})$ tiene dos componentes conexas,
  y $f(S^{n-1})$ es la frontera común de ambas.
\end{corollary}

\begin{theorem}[Invarianza del dominio]
  Sean $A_1, A_2 \subseteq S^n$ subconjuntos de $S^n$ homeomorfos por $h \colon A_1 \to A_2$. Entonces,
  $A_1$ es abierto de $S^n \iff A_2$ es abierto de $S_n$.
\end{theorem}

\begin{remark}
  Nótese que este hecho no es trivial. Lo sería en caso de que $A_1, A_2$ fueran cerrados, o si $h$ estuviera definida
  en todo $S^n$ y fuera un homeomorfismo entre ambos conjuntos.
\end{remark}

\begin{proof}
  Supongamos $A_1$ abierto, y veamos que $A_2$ también lo es. Sea $a_2 \in A_2$, y sea $a_1 \in A_1$ tal que $h(a_1) = a_2$.
  Existe un entorno de $a_1$, al que notamos $V_1 \subseteq S^n$, tal que $V_1 \subseteq A_1$ y $V_1 \cong D^n$. Sean
  $V_2 = h(V_1), W_2 = h(\partial V_1) (\partial V_1 \cong S^{n-1})$. Se verifica que $a_2 \in V_2 - W_2 \subseteq A_2$.

  Por el teorema de separación de Jordan-Brower, $S^n - W_2$ tiene dos componentes conexas, y al ser $V_2 \cong D^n$,
  $S^n - V_2$ es conexo. Además, $S^n - W_2 = (V_2 - W_3) \cup (S^n - V_2)$, siendo esta unión disjunta, luego $V_2 - W_2$
  y $S^n - V_2$ son las componentes conexas de $S^n - W_2$, y por tanto abiertos en $S^n - W_2$ (y en $S^n$).

  Como $a_2 \in V_2  W_2 \subseteq A_2$, $A_2$ es abierto.

  Para demostrar la otra implicación basta con intercambiar los papeles de $A_1$ y $A_2$ y considerar el homeomorfismo inverso.
\end{proof}

\section{Grado de una aplicación continua}

Vamos a definir el grado de una aplicación continua de una esfera en sí misma, aprovechando lo que sabemos
de los grupos de homología de las esferas.

Sea $f \colon S^n \to S^n$ una aplicaciones continua. Los únicos homomorfismos que induce f y que son no triviales son
$f_* \colon H_0(S^n) \to H_0(S^n)$ y $f_* \colon H_n(S^n) \to H_n(S^n)$.

Se cumple que el primero de ellos es siempre la identidad. En el segundo caso, como $H_n(S^n) \cong \Z$, si $[\alpha]$ es un
generador de $H_n(S^n)$, necesariamente se tiene que $f_*([\alpha]) = n \times [\alpha] \quad n \in Z$.

Definimos el grado de $f$, y lo notamos $\Deg f$, a dicho entero $n$.
\[ f_*([\alpha]) = \Deg f \times [\alpha] \]
Veamos algunas de sus propiedades.

\begin{proposition}[Propiedades]
  Estando en el escenario anterior, se verifican:
  \begin{itemize}
    \item[a)] $\Deg Id = 1$.
    \item[b)] $\Deg (g \circ f) = \Deg g \times \Deg f$.
    \item[c)] $f \simeq g \implies \Deg f = \Deg g$. La implicación contraria también es cierta, el llamado Teorema de Hopf.
    \item[d)] Si $f$ es una equivalencia homotópica, $\Deg f = \pm 1$ (en particular, es cierto si $f$ es homeomorfismo).
    \item[e)] Si $\Deg f \neq 0 \implies f$ es sobreyectiva.
  \end{itemize}
\end{proposition}

\begin{proof}
  Las demostraciones son sencillas usando la definición. Probamos únicamente e).

  Supongamos que $f$ no es sobreyectiva, luego existe $x \in S^n$ con $x \notin \Img f$. Así, tenemos el diagrama
  \[ \begin{tikzcd}
    S^n  \arrow{dr}{f'} \arrow{rr}{f} & & S^n  \\
    & S^n - \{x\} \arrow[hookrightarrow]{ur}{i}
  \end{tikzcd} \]
  que induce el siguiente diagrama en la homología:
  \[ \begin{tikzcd}
    H_n(S^n)  \arrow{dr}{f'_*} \arrow{rr}{f_*} & & H_n(S^n)  \\
    & H_n(S^n - \{x\}) \cong 0 \arrow[hookrightarrow]{ur}{i_*}
  \end{tikzcd} \]
  que nos da una descomposición de $f_*$ que tiene al la función nula como factor, lo cual es imposible.
\end{proof}


\cleardoublepage
\ctparttext{En esta parte, analizaremos y resolveremos mediante distintas técnicas el
          problema del clique máximo, comparando sus resultados. Introduciremos el
          trabajo realizado y los recursos utilizados, pasando a una descripción de
          los algoritmos. Continuaremos con detalles específicos de la implementación,
          y expondremos los resultados obtenidos, realizando un análisis individual
          y conjunto de las técnicas incluidas.}
\part{Análisis y resolución del problema del Clique Máximo}

%************************************************
\chapter{Introducción}\label{ch:introduccionMCP}
%************************************************

\section{Introducción al problema}

Dado un grafo $G = (V, E)$, donde $V$ representa el conjunto de vértices
y $E$ el conjunto de aristas, llamaremos un clique a un subgrafo de $G$
en el que todos sus vértices están conectados entre sí. Llamaremos a los
cliques \textbf{maximales} si no son subgrafo de un clique de mayor tamaño,
y \textbf{máximos} si no existe ningún clique de mayor tamaño en el grafo. \\
Es claro que todos los cliques máximos son maximales, pero no al contrario.

El problema del clique máximo, o \textit{maximum clique problem} en inglés (MCP),
consiste en encontrar el clique de mayor número de vértices dentro de un
grafo $G$ cualquiera. Este problema es equivalente a otros dos también muy conocidos,
que son el problema del conjunto independiente, en el que se intenta buscar el mayor
conjunto de nodos no conectados entre sí, y el problema del recubrimiento por vértices,
en el que buscamos el mínimo conjunto de nodos que cubran todas la aristas.

El problema del clique máximo lleva siendo estudiado desde segunda mitad del siglo XX, y fue sujeto
de estudio en el segundo desafío de implementación de DIMACS, que tuvo lugar entre 1992
y 1993, junto con el problema del coloreo de grafos y de satisfabilidad booleana (SAT).
A día de hoy, se encuentran disponibles de forma libre las instancias utilizadas y los
datos de los mejores cliques conocidos hasta la fecha.

En 1977, Garey y Johnson demostraron que el problema es NP-difícil. En 1999, Hastad
probó que, salvo que NP = ZPP, no existe ningún algoritmo capaz de aproximar MCP con
factor $n^{1-\epsilon} \quad \forall \epsilon$, donde la clase ZPP (zero-error
probabilistic polynomial time) es la clase de complejidad de los problemas
reconocibles por algoritmos aleatorizados con complejidad polinómica.

Debido a esto, no es posible obtener algoritmos exactos ni que den una solución
aproximada al problema, de ahí la necesidad del uso de algoritmos heurísticos para
resolverlo.


\section{Trabajo realizado}

En este trabajo he implementado distintos algoritmos, basados en diferentes
metaheurísticas, para comprobar sus eficacia a la hora de resolver el problema.
En varios de ellos se incluyen distintas variantes de la misma heurística, para
analizar cómo afectan cambios en la toma de ciertas decisiones, o la estructura
de elementos del problema.

El objetivo final del trabajo es comparar todos los algoritmos implementados
para comprobar empíricamente su eficacia, basándonos fundamentalmente en la
calidad de las soluciones obtenidas por cada uno de ellos. Veremos como actúan
sobre instancias similares del problema, y discutiremos su viabilidad a la
hora de resolverlo, pues cabe esperar que ciertos algoritmos se comporten
mejor que otros.

Además de la calidad de las soluciones, tendremos en cuenta el tiempo de ejecución
de los algoritmos, si bien no será un aspecto relevante a la hora de evaluarlos.

Las ideas detrás de cada algoritmo provienen de una combinación de mis conocimientos
sobre la materia y el estudio que se lleva realizando sobre el problema durante décadas,
reflejado en numerosos artículos, que han servido como base en la mayoría de implementaciones.
Se referenciará en cada caso el material utilizado y cómo se ha adaptado a este trabajo.

\subsection{Implementación}

La implementación del código de este trabajo ha sido realizada en el lenguaje de
programación \textbf{Ruby}, usando su versión 2.3.0. Ruby es un lenguaje
orientado a objetos, interpretado, con una sintaxis muy similar a la de \textbf{Python}.
Está disponible para descarga bajo una licencia de software libre.

He elegido este lenguaje por haber trabajado con él en varias ocasiones, y aporta
simplicidad e intuitividad a la hora de programar los algoritmos, además de salir
un poco de los lenguajes más utilizados en la carrera.

Al utilizar un lenguaje interpretado como es Ruby, cabe esperar unos
tiempos de ejecución mayores que con otros lenguajes, debido a que no está hecho
para ejecutar código lo más rápido posible. En el caso de querer tener en cuenta la
velocidad, lenguajes compilados como \textbf{C++} serían más útiles que Ruby.

El código implementado cumple varias funciones: la primera de ellas es llevar a
memoria todas las instancias disponibles del problema. Para ello, he tenido que
implementar un lector de ficheros que cree una matriz de adyacencia para cada archivo
disponible. Estos archivos siguen una estructura determinada, que permite
automatizar dicho proceso, no teniendo que considerar casos independientes.
Cada archivo dará lugar a un problema, compuesto de una matriz de adyacencias,
el número de nodos y aristas involucrados, y el nombre del grafo en cuestión.

El segundo punto, y más importante, es la implementación de las heurísticas.
Cada una de ellas está contenida en una clase propia, que contendrá uno o varios
métodos para la resolución de los problemas generados a partir de los archivos.
Adicionalmente, podrán contener métodos auxiliares para uso propio o de otras
heurísticas, en caso de que utilicemos técnicas híbridas. Cada clase estará
en un archivo independiente, y todas ellas se recogerán en un mismo módulo.

Finalmente, existe un archivo adicional con métodos auxiliares, que pueden ser
requeridos por varias heurísticas. La idea bajo su creación es encapsular
operaciones importantes, que serán repetidas frecuentemente y en diversos
escenarios, para facilitar el trabajo y la comprensión del código. Dichos
métodos no se encuentran dentro de una clase, pero sí que están en un
módulo aparte, para diferenciarlos de las implementaciones de las heurísticas.


\section{El repositorio DIMACS}
Las siglas DIMACS hacen referencia al centro de matemática discreta y computación
teórica (del inglés, Center for Discrete Mathematics and Theoretical Computer Science).
Es una colaboración entre dos universidades estadounidenses, las de Rutgers y Princeton,
además de las secciónes de investigación de las empresas AT\&T, Bell Labs, Applied
Communication Science y NEC. Fue fundada en el año 1989, y desde 1990 convoca desafíos
de implementación para determinados problemas de interés en el ámbito científico, con
el objetivo de conseguir algoritmos de calidad.

El problema del clique máximo fue, junto con el coloreo de grafos y el problema de satisfabilidad
de cláusulas (SAT), protagonista del segundo desafío DIMACS, que tuvo lugar entre 1990 y 1991.
Para el problema, se proporcionó un conjunto de instancias de grafos, que son de acceso público,
y son las que utiliza la comunidad científica para evaluar nuevos algoritmos, que he utilizado
para los míos. Pueden encontrarse en la URL \url{http://iridia.ulb.ac.be/~fmascia/maximum_clique/DIMACS-benchmark}.

En dicho repositorio podemos encontrar un total de $37$ grafos, de los cuales he tomado $35$.
Ha quedado fuera los grafos \textit{MANN\_a81} y \textit{MANN_a45}, debido a que por su tamaño
los tiempos de ejecución eran muy elevados. Para ejecutar sobre ellos los algoritmos con los
mismos parámetros que sobre el resto de grafos, se hubieran requerido de más de una semana. 

Los grafos están subdivididos en distintos grupos, que hacen referencia al método de
generación de los mismos. Debido a esto, es posible que algunos algoritmos muestren
un buen comportamiento en ciertos tipos de grafos, y den peores resultados en grafos
que correspondan a otra familia. Los métodos de generación de cada familia de grafos
se pueden consultar en profundidad en la web del repositorio.

En ciertas instancias se ha confirmado el valor del óptimo global, mientras que en
otras solo se conoce una cota inferior a la mejor solución posible. Estos valores
serán tomados como referencia a la hora de evaluar cada algoritmo implementado y
determinar la calidad de las soluciones.


%************************************************
\chapter{Algoritmos}\label{ch:algoritmos}
%************************************************

En este capítulo se tratarán los diferentes algoritmos considerados en el trabajo, desde su
fundamento teórico hasta su implementación concreta.

\section{Consideraciones iniciales}

Antes de describir cada algoritmo, vamos a ver consideraciones más generales sobre el problema, que serán
utilizadas en los algoritmos que se implementen. Discutiremos la representación de las soluciones, el
operador de vecindad, la función objetivo y la inicialización del programa.
% Representación de la solución. Inicialización del programa (lectura). Operador vecino. Entorno. Función objetivo.

Cada solución del problema será un clique que se encuentre dentro del grafo. Este clique estará representado
por los nodos que contiene, de forma que las soluciones serán \textit{arrays} de la forma $[x_0, \dots, x_p]$.

La función para medir la calidad de una solución (función objetivo) será el tamaño de cada clique, si bien
se considerarán algunas variantes en ciertos algoritmos. No obstante, como el objetivo final el obtener un clique
del mayor tamaño posible, es lógico que el factor de peso en cualquier función objetivo sea el tamaño del mismo.

Como operador de vecino tenemos tres distintos, los cuales aplicados sobre el clique nos permitirán obtener
su entorno. El primero de ellos, al que llamaremos operador \textit{add}, consiste en añadir al clique un
nuevo nodo que conserve la estructura de clique, esto es, que esté conectado con todos los nodos del clique.
Usualmente tendremos calculado el conjunto de nodos que podemos añadir a cada clique con el que estemos trabajando.
Llamaremos a dicho conjunto $C_0$ o posibles adiciones.

El segundo operador, al que nos referiremos como \textit{swap}, consiste en intercambiar un nodo del clique
con uno de fuera, de forma que se siga manteniendo la estructura. Para ello, el nodo que no pertenece al
clique debe estar conectado con todos salvo uno de sus nodos, que será el nodo que saldrá del clique. Nos
referiremos a este conjunto de nodos como posibles intercambios o $C_1$. Al igual que el anterior, también
tendremos calculado dicho conjunto de nodos para cada clique con el que trabajemos.

Finalmente, el tercer operador de vecino será el operador \textit{drop}, que consiste simplemente en quitar
un nodo del clique. Es claro que se mantiene la estructura de clique al aplicar este operador.

Estos serán los tres operadores que utilicemos, que definirán el entorno de una solución, compuesto por
aquellas a las que se puede llegar mediante el uso de cualquiera de los tres. No obstante, no siempre se
usarán los tres a la vez. El operador \textit{add} es esencial, pues nos permite ampliar el tamaño de un
clique, por lo que siempre lo utilizaremos. Los otros dos operadores se usarán dependiendo del algoritmo,
siendo el operador \textit{swap} bastante frecuente, y el \textit{drop} menos frecuente. Debido a esto,
el entorno variará dependiendo de aquellos consideramos. Detallaremos en cada caso los operadores
considerados y las razones principales.

\subsection{Inicialización del programa}
Antes de aplicar ningún algoritmo al conjunto de instancias, hay que resolver es la lectura de los
archivos y la creación de las estructuras de datos correspondientes para cada grafo. Para ello,
se ha creado el archivo \textit{reader.rb}, que contiene el método de lectura de los archivos.
Cada uno de ellos dará lugar a un objeto de la clase \textbf{Problem}, que almacena el número de
vértices, el número de arcos, la matriz de adyacencia y el nombre del archivo del que proviene el grafo.

Para la lectura, simplemente recorreremos el directorio en el que están almacenados los archivos,
y generaremos un problema por cada uno de ellos, marcando en la matriz de adyacencia los arcos
que están conectados. Almacenaremos el conjunto de problemas en una variable, para tenerlos
disponibles a la hora de aplicarles los algoritmos.

El archivo \textit{reader.rb} también contendrá el método principal, desde donde se ejecutarán
cada uno de los algoritmos. El cambio de algoritmo y el paso de parámetros se harán modificando
el código correspondiente en el método \textit{main}.

Pasamos ya a ver los algoritmos que se han incluido en este trabajo.

\section{Algoritmos \textit{greedy}}

Los algoritmos voraces, o algoritmos \textit{greedy}, son aquellos que resuelven problemas tomando
las decisiones óptimas en cada instante. Esto no lleva necesariamente a un óptimo global, pero puede
aproximarlo de forma razonable en un periodo de tiempo generalmente muy pequeño. Este tipo de algoritmos
basan sus decisiones en las decisiones que ya se han tomado, pero nunca en lo que podría suceder más
adelante.

Por lo general, en los algoritmos \textit{greedy} tendremos un conjunto de candidatos, una función
de selección, que se encarga de elegir al mejor candidato en cada momento, una función objetivo, que
da un valor a cada solución, una función de factibilidad, que elige los candidatos que se pueden incluir
en una solución, y una función que nos dice si tenemos una solución.

Este tipo de algoritmos suelen ser irreversibles, esto es, no se pueden deshacer elecciones ya hechas.
No obstante, existen varios tipos de algoritmos \textit{greedy}, en los que esta condición se relaja. Así,
distinguiremos entre algoritmos \textit{greedy} puros, relajados y ortogonales. En este trabajo se han
implementado dos algoritmos \textit{greedy}, uno puro y otro relajado.

\subsection{\textit{Greedy} en MCP}

En el problema del clique máximo, los algoritmos \textit{greedy} tienen dos enfoque principales:
o bien se parte de un conjunto vacío y se añaden vértices, respetando la estructura de clique,
hasta que no sea posible añadir más, o partimos del conjunto total de vértices y vamos eliminando
nodos hasta que el grafo obtenido sea un clique.

Los dos algoritmos \textit{greedy} que he implementado siguen el primer enfoque de los
mencionados anteriormente. Se diferencian principalmente en la elección del vértice que
se añade al clique, aunque en uno de ellos modificaremos también la estructura de entornos,
considerando intercambios además de adiciones.

A continuación, se detallan en más profundidad los dos algoritmos implementados:

\subsection{Primer algoritmo}

Este algoritmo es el enfoque más básico posible dentro de los algoritmos voraces para MCP.
Comenzando desde un clique vacío, añade nodos al clique, respetando la estructura, hasta
que no se pueda añadir ninguno más. Para elegir el vértice que añadir, toma, de aquellos que
pueden añadirse, el que tiene mayor número de adyacencias, usando el operador \textit{add}.
Otras posibles opciones serían elegir un nodo aleatorio entre los disponibles o tomar aquel que
maximice el tamaño del conjunto $C_0$ al añadirse al clique.

Este algoritmo es irrevesible, pues los nodos que se añadan al clique no saldrán de él
en ningún caso.

Vemos a continuación el pseudocódigo del algoritmo.

\begin{algorithm}[H]
\caption{Greedy}
  \begin{algorithmic}
  \State $Clique = [ ]$
  \State $C_0 = $ Vértices
  \Repeat
    \State{Añadir $v \in C_0$ con más adyacencias}
    \State{Actualizar $C_0$}
  \Until{$C_0 = \emptyset$}
  \State{Devolver clique}
  \end{algorithmic}
\end{algorithm}


\subsection{Segundo algoritmo}

En este algoritmo, se han seguido las ideas de Grosso, Locatelli y Della Croce, usando un algoritmo
que nos permite deshacer, en cierta forma, decisiones tomadas anteriormente. La idea es no solo
considerar aquellos vértices que pueden ser añadidos al clique, sino también tomar aquellos que
pueden ser intercambiados por uno del clique, es decir, usar operadores \textit{add} y \textit{swap}.

El criterio para elegir un nodo será el número de conexiones con el conjunto de posibles adiciones
al clique, buscando maximizar este número. Así, al añadir el nodo al grafo, el tamaño del conjunto
$C_0$ será el máximo de los posibles. Como dicho tamaño es una cota para el número de nodos que se
pueden añadir al clique, nos interesa que sea lo más elevado posible.

Antes de empezar a considerar intercambios, queremos que el clique tenga un tamaño mínimo.
Se ha fijado dicho tamaño en $4$, valor a partir del cual se empiezan a considerar los
intercambios y no solo las adiciones. Detendremos el algoritmo una vez lleguemos a un
clique maximal, esto es, cuando el conjunto $C_0$ sea vacío. No obstante, se ha fijado
un valor máximo de intercambios para evitar que el algoritmo se prolongue durante mucho tiempo.
Una vez satisfecha alguna de estas condiciones, el clique que tengamos será ampliado
hasta su límite usando movimientos \textit{add}, finalizando cuando no sea posible
añadir más nodos. Esta ampliación se realiza por si el algoritmo se ha detenido al
alcanzar el límite de intercambios, lo que significa que es posible añadir más nodos.

Para evitar repetición de intercambios, en el momento en el que se haga uno, guardaremos
el nodo que sale del clique, para no tenerlo en cuenta en la siguiente iteración.
De no hacer esto, cabe la posibilidad de entrar en un bucle, en el que se
intercambien dos nodos de forma indefinida hasta que se alcance el criterio de parada.

Vemos a continuación el pseudocódigo de la parte adaptativa del algoritmo.


\begin{algorithm}[H]
\caption{Greedy adaptativo}
  \begin{algorithmic}
  \State $Clique = [ ]$
  \State $C_0 = $ Vértices
  \State $C_1 = [ ]$
  \State $i = 0$
  \Repeat
    \State Tomar $v \in C_0 \cup C_1 - Swap$ que maximice adyacencias con $C_0$
    \If {$v \in C_0$}
      \State Añadir $v$ a $Clique$
    \Else
      \State Hacer intercambio
      \State Actualizar $Swap$, $i += 1$
    \EndIf
    \State{Actualizar $C_0, C_1$}
  \Until{$C_0 = \emptyset$ o $i = limite$}
  \State{Devolver clique}
  \end{algorithmic}
\end{algorithm}


\section{Búsqueda local}

Los algoritmos de búsqueda local son aquellos que se centran en el entorno de una
solución para buscar nuevas soluciones que sean mejores que la anterior. Dicho entorno
viene dado por aquellas soluciones a las que se puede llegar desde la solución de
partida mediante un \textbf{operador de vecindad}. Llamaremos a estas soluciones
\textbf{soluciones vecinas}.

Las distintas formas de recorrer el entorno y de aceptar soluciones vecinas da
lugar a los distintos algoritmos de búsqueda local. El criterio de aceptación debe
ser monótona, pues queremos que conduzca a soluciones cada vez mejores, finalizando
en un óptimo local. La búsqueda local básica busca mejorar el resultado con un elemento
del entorno, terminando si este no existe y alcanzando un máximo local. Si queremos
permitir que las soluciones puedan ser peores, existen algoritmos como el enfriamiento
simulado o la búsqueda tabú. En cuando a la estructura de entornos, también puede ser
variable, lo que nos llevaría a una búsqueda en entornos variables o, nuevamente,
a una búsqueda tabú. Todo esto se hace con el objetivo de que nuestro algoritmo
sea capaz de escapar de óptimos locales, y alcance el óptimo global.

Las técnicas de búsqueda local se utilizan no solo como heurísticas para resolver
problemas, sino como complementos de otras. Por ejemplo, pueden ser utilizadas junto
a los algoritmos genéticos para mejorar el conjunto de soluciones, o aplicarse después
de técnicas \textit{greedy} para explorar el espacio de soluciones de la solución
generada. También se usan en técnicas multiarranque como ILS o GRASP, que explicaremos
más adelante.

En este trabajo, se han implementado dos técnicas de búsqueda local, que difieren en la
estructura de entornos y en el método de elección de candidatos. Asimismo, se ha implementado
un algoritmo de enfriamiento simulado, y se han incorporado varias de las técnicas mencionadas
anteriormente en otros algoritmos. Estas serán detalladas en cada caso.

\subsection{Búsqueda local en MCP}

En la búsqueda local, las diferencias que hacen a los algoritmos distintos están en la estructura
de entornos y en elección de una solución vecina a la que desplazarse. Estas diferencias son
también las que existen entre los dos algoritmos implementados en este trabajo. Nos valdremos
de los operadores \textit{add, swap} y \textit{drop} como operadores de vecindad, si bien no
siempre los tendremos todos en cuenta.

En el primer algoritmo se consideran los tres movimientos para calcular el entorno.
Si bien esto rompe con la filosofía de que la búsqueda local debe ir a mejores soluciones,
la posibilidad de eliminar elementos del clique se incluye para evitar que estos algoritmos
sean similares a los \textit{greedy}, buscando favorecer una búsqueda en el entorno.

En el segundo algoritmo, solo se consideran movimientos \textit{add} y \textit{swap}.
Si bien puede en principio puede parecer similar al algoritmo \textit{greedy} adaptativo
detallado anteriormente, veremos que tiene varias diferencias fundamentales con respecto
a este.

Veámoslos en más detalle.

\subsection{Primer algoritmo}

Este algoritmo se basa en las ideas de Katayama, Hamamoto y Narihisa. En él, se han considerado
los tres movimientos básicos de la búsqueda local, \textit{add, swap} y \textit{drop}, dando
distintas prioridades a cada uno de ellos. Como nuestro objetivo es construir un clique del
mayor tamaño posible, añadiremos vértices siempre que podamos, tomando aquel que tenga más
conexiones con el conjunto $C_0$ en caso de haber varios. Cuando no sea posible, intentaremos
hacer un swap, siempre que este nos lleve a un escenario en el que podamos añadir más nodos
al clique resultante. La elección de los nodos a intercambiar se hace mediante la técnica del
primer mejor, esto es, tomamos el primero que encontremos. De no poder hacer ningún swap bajo
estas condiciones, eliminaremos del clique el nodo que tenga menos adyacencias con todos los
nodos del grafo.

Para evitar que los nodos que eliminemos sean añadidos inmediatamente al clique, los iremos
acumulando en una lista tabú. Los nodos de esta lista no se tendrán en cuenta a la hora de
hacer un movimiento add. Así, nos aseguramos que el clique cambia, y que no entramos en
un bucle en el que eliminamos un nodo y acto seguido lo volvemos a añadir al clique.
Esta lista tabú será vaciada en el momento en el que se realice un movimiento que no sea
quitar un nodo del clique, pues nos habremos asegurado de que el clique cambia.

Este algoritmo no tiene un criterio de parada inducido por el entorno, pues siempre podremos
hacer alguno de los tres movimientos. Por tanto, es necesario establecer un límite, que este
caso, ha sido establecer un límite de swaps y drops a realizar a lo largo del algoritmo.
Una vez alcanzado este límite, el algoritmo finaliza, devolviendo el mejor clique encontrado.
Para evitar que el límite sea una constante común para todas las instancias, se ha establecido
en proporción al número de vértices del grafo, haciendo que el algoritmo se adapte al tamaño
de cada problema.

El hecho de considerar un criterio de elección sobre los nodos del conjunto $C_0$ y $C_1$ hace que
el orden del algoritmo sea $\mathcal{O}(nl)$, donde $n$ el el tamaño del grafo y $l$ el límite impuesto.

Este algoritmo puede trabajar partiendo desde cualquier clique, lo que nos permite utilizarlo en
técnicas híbridas como algoritmos meméticos, o combinarlo con el uso de un algoritmo \textit{greedy}
para tener una solución de partida prometedora. Para evaluar sus prestaciones por separado,
partirá desde un clique vacío, que es equivalente a pasarle una solución inicial construída
con un algoritmo \textit{greedy} que añada nodos utilizando el mismo criterio que esta búsqueda local.

Vemos el pseudocódigo del algoritmo:

\begin{algorithm}[H]
\caption{Búsqueda local 1}
  \begin{algorithmic}
    \State $i = 0$
    \State $Tabu = \emptyset$
    \State Calcular $C_0$ y $C_1$.
    \Repeat
      \If {$C_0$ no es vacío}
        \State Añadir el nodo de $C_0$ que tenga más conexiones con $C_0$.
        \State $Tabu = \emptyset$
      \Else
        \State Buscar en $C_1$ el primer intercambio que nos da un $C_0$ no vacío.
        \If {El intercambio existe}
          \State Hacer intercambio
          \State $Tabu = \emptyset$
          \State $i = i + 1$
        \Else
          \State Quitar el nodo del clique con más adyacencias en el grafo.
          \State Añadirlo a $Tabu$.
          \State $i = i + 1$
        \EndIf
      \EndIf
      \State{Actualizar $C_0, C_1$, teniendo en cuenta $Tabu$.}
    \Until{$i = limite$}
    \State{Devolver mejor clique encontrado}
  \end{algorithmic}
\end{algorithm}

\subsection{Segundo algoritmo}

Este algoritmo está basado en la búsqueda local dinámica o DLS (del inglés, \textit{Dynamic Local Search}),
propuesta por W. Pullan y H. Hoos en el año 2006, y en concreto, se basa en una adaptación hecha
por A. Grosso, M. Locatelli y W. Pullan para un algoritmo de búsqueda local iterada, que veremos
más adelante. Sigue una filosofía más simple que el anterior pues solo considera movimientos
\textit{add} y \textit{swap}, a los que volveremos a ordenar por prioridad. Además, incluye una
lista tabú con los nodos que han salido del clique mediante un swap, evitando que vuelvan a entrar
al clique y proporcionando diversidad. Limitaremos la longitud de esta lista a un 2\% del número
de vértices en el grafo, para que no aloje un gran número de nodos.

Nuestro entorno consistirá en el conjunto de posibles adiciones ($C_0$) y el conjunto de
posibles intercambios ($C_1$), al que quitaremos la lista tabú. Como teóricamente, el entorno
puede ser vacío, si se diera el caso el algoritmo finalizaría. No obstante, se ha establecido un
límite de intercambios, pues en la práctica el entorno no tiene por qué ser vacío.

La prioridad entre movimientos es la siguiente: primero intentaremos añadir un elemento de
$C_0$ que no esté en la lista tabú, comprobando si existe alguno, y tomando uno aleatorio
en caso de existir varios. De no existir ninguno, tomaremos el conjunto $C_1$ y le quitaremos
la lista tabú. En caso de quedar algún nodo, haremos el intercambio, que volverá a ser aleatorio
si tenemos varios candidatos. Dicho nodo será añadido a la lista tabú para evitar que vuelva
a entrar al clique en las iteraciones más inmediatas. Finalmente, en caso de que este conjunto
también fuera vacío, solo nos queda la posibilidad de añadir un elemento de la lista tabú que
esté en $C_0$. Nuevamente, elegiremos uno aleatorio entre los existentes.

Este proceso se repetirá hasta que se satisfaga uno de los dos criterios de parada. Una vez
acabe el algoritmo, devolveremos el mejor clique que haya encontrado.

Como en este caso la elección de nodos es aleatoria, se considera un único bucle,
luego la complejidad de este algoritmo es lineal sobre el límite establecido,
$\mathcal{O}(l)$. Debido a esto, ofrecerá mejores tiempos de cómputo que el otro
algoritmo de búsqueda local implementado.

Este algoritmo también puede trabajar con cualquier clique de partida, siendo posible
utilizarlo como complemento a otras técnicas. Para evaluar su funcionamiento, tomaremos
como clique inicial un clique vacío, y dado el orden de prioridad de los movimientos,
el algoritmo creará un clique maximal de forma aleatoria y a partir de él explorará
su entorno.

Pasamos a ver el pseudocódigo del algoritmo:

\begin{algorithm}[H]
\caption{Búsqueda local 2 - DLS}
  \begin{algorithmic}
    \State $i = 0$
    \State $Tabu = \emptyset$
    \State Calcular $C_0$ y $C_1$.
    \Repeat
      \If{$C_0 - Tabu$ no es vacío}
        \State Añadir un nodo aleatorio de $C_0 - Tabu$
      \ElsIf{$C_1 - Tabu$ no es vacío}
        \State Hacer un cambio con un nodo aleatorio de $C_1 - Tabu$.
        \State Añadir el nodo a la lista tabú que sale del clique a $Tabu$
        \State $i = i + 1$
      \ElsIf{$C_0$ es no vacío}
        \State Añadir un nodo aleatorio de $C_0$.
      \EndIf
      \State Actualizar $C_0, C_1$
      \State Entorno = $C_0 \cup (C_1 - Tabu)$
    \Until{$i = limite$ o el entorno sea vacío.}
    \State{Devolver mejor clique encontrado}
  \end{algorithmic}
\end{algorithm}

\section{Enfriamiento simulado}

El enfriamiento simulado es un algoritmo de búsqueda en entornos, que incluye un
criterio probabilístico de aceptación de soluciones basado en la termodinámica,
el algoritmo de Metrópolis. Fue introducido formalmente en 1983.

Su principal objetivo es evitar que la búsqueda en el entorno finalice en un óptimo
local, cosa que ocurre con la búsqueda local. Para esto, permitirá que algunos
movimientos vayan a soluciones peores, usando una función de probabilidad que decida
si se acepta un movimiento a una solución de menos calidad. Esta probabilidad vendrá
dada por el algoritmo de Metrópolis, e irá descendiendo a medida que avanza el algoritmo.
Así, conseguimos una diversificación de soluciones al principio, mientras que al final
del algoritmo intensificamos la búsqueda en una región concreta del espacio de soluciones.

El algoritmo de Metrópolis utiliza dos valores para calcular la probabilidad. Uno de ellos
es la diferencia entre los valores de la función objetivo de la solución del entorno y
la solución actual. El otro es un valor llamado temperatura, en referencia a su origen
termodinámico, que desciende en cada iteración de la búsqueda, lo que permitirá que
la probabilidad descienda con el valor.

Su forma de actuar es la siguiente: si nos movemos a una solución que mejore el resultado
anterior, la aceptamos directamente. De no ser así, calculamos la diferencia del valor de
la función objetivo, que notamos $\Delta F$, y tomamos como probabilidad de aceptación
$e^{\nicefrac{\Delta F}{T}}$, donde $T$ representa la temperatura. Como $\Delta F \leq 0$,
este valor estará siempre comprendido entre $0$ y $1$, siendo más cercano a $0$ cuanto
menor sean la temperatura y $\Delta F$.

El decrecimiento de la temperatura nos proporciona también un criterio de parada. Fijando
una temperatura final, usualmente cercana a $0$, el algoritmo se ejecutará hasta que la
temperatura alcance dicho valor. Así, conseguimos que la probabilidad de aceptación de
soluciones peores sea mayor al principio, favoreciendo la exploración, mientras que en
las fases finales la probabilidad de aceptación será baja, y el algoritmo intensificará
la búsqueda en la región del espacio de soluciones que se encuentre.

El descenso de la temperatura puede seguir distintos esquemas, aunque el más usual es el
geométrico, consistente en multiplicar la temperatura por una constante menor que $1$
en cada iteración. No obstante, cualquier función estrictamente decreciente nos resultaría
útil para realizar dicho descenso.

\subsection{Enfriamiento simulado en MCP}

El principal problema a la hora de elaborar un algoritmo de enfriamiento simulado en MCP
es la elección de una función objetivo que permita una exploración amplia del entorno.
Las funciones objetivo usuales, como el tamaño del clique, no aportan diversidad de soluciones,
pues una vez lleguemos a un óptimo local, normalmente repetirá movimientos de intercambio
de forma indefinida, ya que estos no cambian el valor de la función objetivo, y son
aceptados por el criterio de Metrópolis.

Establecer una prioridad entre los tres operadores de vecindad daría como resultado un
algoritmo muy similar a la búsqueda local, que seguiría teniendo los mismos problemas.
Por lo tanto, centraremos nuestro esfuerzo en buscar una función objetivo que nos
permita una mayor exploración. No obstante, hemos de tener en cuenta que el algoritmo
debe seguir convergiendo a un óptimo local, por lo que no nos vale cualquier función.

Es posible encontrar algunos trabajos sobre MCP que involucran el enfriamiento simulado.
Sin embargo, se suelen basar en una representación y una estructura diferentes a las
consideradas en este problema, siendo imposible trasladarlos completamente a este caso,
aunque sí pueden traerse algunas de las ideas principales. En el caso de este trabajo,
se ha implementado un algoritmo que utiliza las nociones básicas de enfriamiento simulado
para ver sus diferencias con la búsqueda local, y un segundo algoritmo que recoge algunas
ideas de un artículo sobre el problema.

\subsection{Primer algoritmo}

Este algoritmo es una adaptación de una búsqueda local a un enfriamiento simulado, utilizando
una función objetivo basada en el tamaño de los cliques. En concreto, el valor de cada clique
es su tamaño más el número de vértices que contiene el conjunto $C_0$ dividido por el número
de vértices totales.
\[ f(C) = |C| + \frac{|C_0|}{N} \]
De esta forma, conseguimos que cliques de mayor tamaño tengan siempre un valor mayor que
aquellos con menor tamaño, lo que garantiza la convergencia a un óptimo local. En caso
de empate, utilizamos el tamaño del conjunto $C_0$ para determinar el mejor, pues, como
hemos discutido con anterioridad, cuanto mayor sea $C_0$ más puede crecer el clique.

Para el entorno, he considerado movimientos \textit{add}, \textit{swap} y \textit{drop},
sin priorizar ninguno de ellos. Se tomarán en orden aleatorio hasta encontrar uno que
mejore la función objetivo o satisfaga el criterio de Metrópolis.

Los valores de temperatura han sido inicializados a $T_{\text{inicial}} = 1,
T_{\text{final}} = 0.001$, con un mecanismo de enfriamiento geométrico, consistente
en multiplicar la temperatura por una constante $\beta$, fijada a $\beta = 0.99$.
Una vez se alcance una temperatura menor que la final, completamos el mejor clique
encontrado añadiendo tantos nodos como sea posible (aunque lo más probable es que
no se añada ninguno), eligiendo dichos nodos por mayor número de adyacencias.
Finalmente, el algoritmo devuelve el clique y finaliza.

El algoritmo trabaja siempre partiendo del clique vacío. No obstante, puede ser
adaptado fácilmente para que comience la exploración desde cualquier clique, lo
que nos permitiría combinarlo con otros algoritmos para obtener técnicas híbridas.

Vemos una descripción del algoritmo en pseudocódigo.

\begin{algorithm}[H]
\caption{Enfriamiento Simulado}
  \begin{algorithmic}
    \State $T = 1$
    \State $T_{\text{final}} = 0.001$
    \State $\beta = 0.99$
    \State Calcular $C_0$ y $C_1$.
    \Repeat
      \State Entorno = $C_0 \cup C_1 \cup$ Drops
      \Repeat
        \State Sacar un clique del entorno, calcular $\Delta F$.
        \If{$\Delta F \geq 0$}
          \State Aceptar solución.
          \State Salir del bucle interno.
        \Else
          \State Aceptar solución con probabilidad $e^{\nicefrac{\Delta F}{T}}$.
          \State Salir del bucle interno si se acepta la solución.
        \EndIf
      \Until{Entorno $= \emptyset$.}
      \State $T = T \times \beta$
    \Until{$T < T_{\text{final}}$}
    \State{Ampliar mejor clique encontrado y devolver el resultado.}
  \end{algorithmic}
\end{algorithm}

\subsection{Segundo algoritmo}

En el segundo algoritmo, se han adaptado las ideas de X. Geng, J. Xu, J. Xiao y L. Pan,
que propusieron un algoritmo de enfriamiento simulado para MCP. Dicho algoritmo sigue
un enfoque distinto al de este problema, pues intenta buscar cliques de un tamaño
predeterminado, explorando el espacio de soluciones hasta que lo encuentra, o hasta
llegar al criterio de parada. Se basa en una representación por permutaciones del
vector de vértices, y va realizando intercambios de dos vértices, de tal forma que
uno sale del subgrafo y otro entra en él, siempre que la función objetivo mejore o
se satisfaga el criterio de Metrópolis. En este caso, en lugar de ir moviéndose entre
cliques, se consideran grafos de tamaño fijo, que se modificarán con el objetivo de
alcanzar un clique. Esto permite definir como función objetivo el número de aristas
que le faltan al grafo para ser un clique.

Para poder adaptarme a este enfoque, he tenido que diseñar un algoritmo que se moviera
entre grafos que no fueran un clique. He utilizado los mismos tres operadores \textit{add, swap}
y \textit{drop}, pero en este caso, he permitido añadir cualquier nodo al clique o
intercambiar dos nodos cualesquiera, sin necesidad de respetar la estructura de clique.
La exploración del nuevo entorno se hace igual que en el caso anterior, de forma aleatoria.

Al principio tomé la mismo función objetivo que la propuesta en el artículo, consistente
en el número de aristas necesarias para que el grafo sea un clique. Sin embargo, esta
función no asegura una convergencia a soluciones de calidad; en particular, es ajena al
tamaño de los cliques. Por esto, he usado la longitud de los grafos para corregir esta,
restándola al número de aristas restante. Así, tenemos una función objetivo que es menor
cuanto mayor es el grafo, y cuanto más se parece a un clique.

Una vez definida la función objetivo, el algoritmo explora el entorno y toma un elemento
aleatorio hasta que la función objetivo es menor, o hasta que se acepta el criterio de
Metrópolis. Como ahora no trabajamos con cliques, el grafo que obtenemos será reducido
a un clique, para lo que eliminaremos nodos uno a uno hasta que el resultado sea un clique.
Eliminaremos el nodo que menos adyacencias tenga. El clique obtenido no sustituye al grafo
que tengamos, sino que se compara con el mejor clique obtenido hasta el momento, reemplazándolo
si su longitud fuera mayor. Después del algoritmo continua explorando el entorno del
grafo actual, repitiendo el método descrito.

He vuelto a tomar los valores $T_{\text{inicial}} = 1, T_{\text{final}} = 0.001,
\beta = 0.99$ para los parámetros del algoritmo. Una vez finaliza la fase de enfriamiento,
vuelvo a añadir tantos nodos como sea posible al mejor clique encontrado, mediante un algoritmo
\textit{greedy} que prioriza los nodos con más adyacencias, devolviendo el clique resultante.



\section{Búsquedas multiarranque}

Las búsquedas multiarranque son un tipo de búsqueda local, que intentar salir de óptimos
locales reiniciando la búsqueda para tomar otra solución de partida, buscando explorar
una región distinta del espacio de soluciones. Dicha reinicialización de la búsqueda será
la clave para distinguir entre búsquedas multiarranque, y podemos distinguir entre las siguientes:

Búsqueda multiarrranque básica, que genera una solución inicial aleatoria y explora su entorno
mediante búsqueda local, repitiendo el proceso hasta que se satisfaga una condición de parada.
Es el método más sencillo, que no utiliza ningún tipo de información del problema, y equivale
a aplicar una búsqueda local múltiples veces sobre distintas soluciones aleatorias de partida.

Métodos constructivos de la solución inicial, donde esta varía en cada iteración. En este caso,
la solución inicial se obtiene mediante una construcción, intentando conseguir soluciones de
partida que sean de calidad. Intentaremos que las soluciones sean distintas en cada iteración,
por lo que el proceso de construcción será usualmente aleatorizado. Una vez se tenga la solución
inicial, se aplica búsqueda local, repitiendo hasta satisfacer una condición de parada,
normalmente un número de iteraciones fijo de antemano.

Métodos basados en la modificación del óptimo encontrado, en los que el óptimo local se somete
a una perturbación, dando lugar a una nueva solución de partida. Con esto utilizamos una
solución ya calculada para proporcionar una nueva solución de partida, con la que intentaremos
mejorar la calidad de la solución anterior. Nuevamente, repetiremos el proceso hasta alcanzar
una condición, que habitualmente será un número de repeticiones dado.

En este trabajo he considerado dos algoritmos de este tipo; uno basado en construcción de
la solución inicial, el algoritmo GRASP (\textit{Greedy Randomized Adaptative Search Procedure}),
y otro basado en modificaciones, el algoritmo ILS (\textit{Iterated Local Search}).
Vamos a detallarlos.

\subsection{GRASP}

El algoritmo GRASP, o \textit{greedy} aleatorizado, es un algoritmo multiarranque que consiste
en repetir una fase de construcción de una solución, seguida de una búsqueda local en el entorno
de dicha solución, hasta que se alcance una condición, que suele ser un límite de iteraciones
o de tiempo de ejecución. El un algoritmo sencillo de programar, y su capacidad para obtener
buenas soluciones está ligada a si los algoritmos \textit{greedy} y la búsqueda local ofrecen
buenos resultados sobre el problema, pues está fundamentado en ellos. Fue introducido formalmente
en 1989 (Feo, Resende), aunque la idea había sido descrita en 1987.

La construcción de la solución inicial sigue un procedimiento \textit{greedy}, en el cual,
en lugar de proporcionar un único candidato como se haría en un algoritmo \textit{greedy}
clásico, se proporcionan varios de los mejores, y de entre estos, elegiremos uno al azar.
El número de candidatos a considerar puede dejarse como parámetro, o fijarse de antemano,
quedando la decisión en manos del programador. Nótese que si el número de candidatos es
$1$, el algoritmo será un \textit{greedy} puro.

Esta técnica se utiliza para proporcionar diversidad, pues esta se pierde al tomar soluciones
iniciales generadas por algoritmos voraces. Así, seguimos obteniendo buenas soluciones de
partida, a la vez que estas difieren lo suficiente como para explorar una región amplia del
espacio de soluciones.

Una vez construída la solución inicial, se explora su entorno siguiendo una búsqueda local,
hasta alcanzar un óptimo local. Una vez alcanzado, se compara la solución con la mejor obtenida
hasta el momento, se actualiza si es necesario, y se comienza de nuevo el algoritmo.
Se repetirá el proceso hasta satisfacer una condición de parada.

\subsubsection{GRASP en MCP}

El aspecto fundamental de este algoritmo es la construcción de la solución inicial.
Podemos valorar los elementos del conjunto de candidatos de distintas formas, lo que
dará lugar a distintas soluciones de partida. Por lo demás, el algoritmo se basa en
una búsqueda local, que hemos discutido con anterioridad.

\subsubsection{Algoritmo implementado}

Como se acaba de mencionar, el aspecto fundamental de este algoritmo es la generación de
la solución inicial. Para ello, he tomado la lista de candidatos, $C_0$, y la he ordenado
por número de adyacencias en el grafo, en orden descendente. Una vez ordenada, tomo la
mitad superior como conjunto de candidatos, y hago una elección aleatoria entre los nodos,
en la que todos tienen la misma probabilidad de ser elegidos. El proceso se repetirá hasta
que $C_0$ sea vacío.

La fase de búsqueda local utiliza los dos métodos ya implementados, lo que da lugar a
dos versiones del algoritmo, en las que solo varía la búsqueda local. Como estos métodos
no se detienen en un óptimo, sino que necesitan un número de iteraciones máximo, he
establecido como límite el número de vértices del grafo, para que se adapte al tamaño
de cada problema.

Este proceso se repetirá un número determinado de iteraciones, que se pasará como parámetro.
Para este caso se ha fijado a $20$ en todas las instancias disponibles. Almacenaremos la
mejor solución obtenida, que será la que el algoritmo devuelva al final.

Vemos el pseudocódigo del algoritmo y de la generación de soluciones iniciales.

\begin{algorithm}[H]
\caption{GRASP}
  \begin{algorithmic}
  \Repeat
    \State Obtener solución inicial.
    \State Aplicar búsqueda local a la solución inicial.
    \State Actualizar mejor solución.
  \Until{Se alcance el número de iteraciones.}
  \State{Devolver mejor solución obtenida.}
  \end{algorithmic}
\end{algorithm}

\begin{algorithm}[H]
\caption{Generación de soluciones aleatorias}
  \begin{algorithmic}
    \State $Clique = \emptyset$
    \State $C_0 =$ vértices
    \Repeat
      \State Ordenar $C_0$ por adyacencias.
      \State Tomar la mitad superior de $C_0$.
      \State Añadir un elemento aleatorio a $Clique$.
      \State Recalcular $C_0$.
    \Until{$C_0 = \emptyset$}
    \State Devolver $Clique$
  \end{algorithmic}
\end{algorithm}


\subsection{ILS}

La búsqueda local iterada, o ILS (del inglés, \textit{Iterated Local Search}), es
un algoritmo de búsqueda local multiarranque basado en la perturbación de soluciones
como método para obtener diversidad. Fue propuesta en 1998 por Thomas Stützle,
y desde entonces se utiliza en la resolución de distintos problemas, entre los
que se encuentra el problema del clique máximo.

El algoritmo ILS requiere de cuatro componentes principales: una solución inicial,
un algoritmo de mutación de soluciones, un método de búsqueda local, y un criterio
de aceptación. Definiendo estos cuatro componentes en un problema, estaremos en
condiciones de aplicar el algoritmo.

El funcionamiento general de este tipo de algoritmos es el siguiente: primero, se
toma la solución inicial y se le aplica búsqueda local, hasta llegar a un óptimo local.
Seguidamente, se le aplica una mutación a la solución obtenida, que nos permite crear
una nueva solución inicial a la que volver a aplicar búsqueda local. Una vez finalizada
esta segunda búsqueda local, se usa el criterio de aceptación para ver si reemplazamos
la solución anterior con la obtenida. Este proceso mutación-búsqueda local-aceptación
se repetirá hasta que se satisfaga un criterio de parada. Finalmente, el algoritmo
devuelve la mejor solución encontrada.

Este modelo de ILS es el más básico que existe, pero no es el único. Podemos utilizar
un modelo basado en poblaciones, donde a cada elemento se le aplica ILS y estos se
combinan de alguna forma para una mayor exploración del espacio de soluciones, o el
modelo $(\nu + \lambda)$, donde, partiendo de $\nu$ soluciones iniciales, generamos
$\lambda$ hijos a partir de ellas mediante mutación, que compiten para pasar a la siguiente fase.

También es posible establecer diferencias en el criterio de aceptación. Si bien el
método más usual es tener una función objetivo, pueden utilizarse otros, como una
selección aleatoria, una selección basada en combinaciones, o una basada en probabilidades,
en la que podría utilizarse el algoritmo de Metrópolis, por ejemplo.

\subsubsection{ILS en MCP}

Es sencillo construir una búsqueda local iterada para el problema del clique máximo,
partiendo de un algoritmo de búsqueda local existente. Para ello, son necesarios un
criterio de aceptación, una solución inicial y una función de mutación.

\subsubsection{Algoritmo implementado}

Se ha implementado un algoritmo que sigue la filosofía de la búsqueda local iterada
básica basado en el propuesto por A. Grosso, M. Locatelli y W. Pullan, que usa el
algoritmo de búsqueda local DLS, si bien se propusieron algunas modificaciones que
aquí no se han tenido en cuenta. En mi caso, además de DLS, he utilizado el otro
algoritmo de búsqueda local implementado, para comparar las diferencias entre ambos
al ser utilizados por ILS. Vemos a continuación la función de mutación y el criterio
de aceptación utilizados.

El algoritmo para mutar un clique consiste en tomar un nodo que no se encuentre en él,
eliminar del clique todos los nodos que no estén conectados a él, y finalmente añadir
dicho nodo al clique. Con este método conseguimos una solución de partida que se encuentra
en una zona distinta del espacio de soluciones, a la vez que usamos información de la
iteración anterior para movernos por soluciones que serán potencialemente de calidad.

La función de aceptación se queda siempre con el último clique obtenido tras la búsqueda local.
Tras valorar si quedarme con la última o con la mejor encontrada, decidí hacer lo primero,
para favorecer la diversidad en las soluciones. No obstante, tendremos almacenada la mejor
solución, al ser el valor de salida del programa.

Dado que las dos búsquedas locales implementadas requieren de un límite como criterio de
parada, he fijado dicho límite al número de vértices del grafo, para que se adapte al
tamaño de cada problema. Además, el criterio de parada del algoritmo también es un número
de iteraciones, que ha sido $20$ para todas las instancias.

Vemos el pseudocódigo del algoritmo:

\begin{algorithm}[H]
\caption{ILS}
  \begin{algorithmic}
    \State Generar solución inicial.
    \State Aplicar búsqueda local.
    \Repeat
      \State Modificar solución.
      \State Aplicar búsqueda local a la solución obtenida.
      \State Actualizar la mejor solución, si es necesario.
    \Until{Se alcance el número de iteraciones.}
    \State{Devolver mejor solución obtenida.}
  \end{algorithmic}
\end{algorithm}

\section{Algoritmos de colonia de hormigas}

Los \textbf{algoritmos de colonia de hormigas} o ACO (del inglés, \textit{Ant Colony Optimization})
son una técnica de adaptación social, que se basa en el comportamiento de las colonias de
hormigas, para tratar de encontrar caminos en grafos. Para ello, utilizará un recurso
que emula a las feromonas que emiten las hormigas cuando se desplazan.

Las hormigas utilizan este rastro de feromonas para guiarse, pues poseen unos receptores
capaces de captar la intensidad de dicho rastro. Dado su pésimo sentido de la vista,
el uso de las feromonas es vital para que las hormigas se desplacen y sean capaces de
localizar fuentes de alimento, volver al hormiguero, y que el resto de la colonia también
sea capaz de llegar a la comida, siguiendo el rastro que las propia hormigas dejan.

Este rastro de feromonas no permanece de forma indefinida en el ambiente, sino que
sufre un proceso de evaporación constante. De esta forma, caminos más cortos tenderán
a una mayor concentración de feromonas, lo que incrementará el tránsito de hormigas.
Por otra parte, los caminos más largos recibiran menos aporte de feromonas, lo que a
la larga hará que dejen de ser visitados.

La elección por parte de una hormiga del camino que recorrer se fundamenta en un
criterio probabilístico. La hormiga tiende a recorrer caminos en los que el rastro
de feromona es mayor, aunque podrá ir por otros donde el rastro no sea tan intenso.

Su fundamentación teórica hace que sea un algoritmo muy eficaz para problemas que
requieren de construcción de caminos, como el problema del viajante de comercio.
En general, son algoritmos que funcionan bien cuando trabajan sobre estructuras
similares a los grafos.

Este comportamiento será trasladado a la optimización por colonia de hormigas,
mediante la construcción de hormigas artificiales. Tendremos un rastro de feromona
artificial, que simula el que dejan las hormigas al recorrer un camino. Es posible
complementar el uso de las feromonas con información heurística del problema,
con el objetivo de crear un mecanismo probabilistico de construcción de soluciones
que no solo se base en feromonas, sino que use más información disponible.

Las hormigas serán las encargadas de ir recorriendo el grafo, utilizando la información
de la que dispone para decidir a qué nodo moverse mediante una función de probabilidad.
Los nodos que compongan una solución serán reforzados con un aumento de feromonas, y
se hará de forma que mejores soluciones reciban más cantidad de feromona.
Se utiliza también un mecanismo de evaporación de feromona simulando al que ocurre en
la naturaleza, que evita un crecimiento ilimitado de la misma. Para esto, se reducirán
todos los rastros de feromona eliminando un parte de su valor actual, usualmente un
porcentaje fijo. En la práctica, esto se traduce en evitar un estancamiento del algoritmo,
debido a que los nodos más visitados tendrían altas probabilidades de repetirse en la
siguiente solución.

Por norma general, se lanzarán un número determinado de hormigas artificiales independientes,
las cuales generan una solución cada una. Entre ellas, se toma la mejor, y esta será la
que cambie los niveles de feromona. Este proceso se repite hasta que se satisfaga una
condición de parada, que suele ser un número de iteraciones o un límite temporal.

Existen numerosas variantes de las colonias de hormigas, dependiendo de la función de
probabilidad, la actualización de feromonas, etcétera. En este trabajo he tomado el
algoritmo más clásico, sin entrar en otros que requerirían de mayor experimentación
para hacer ajustes sobre los parámetros.

También es posible hibridar este método con otros, usualmente un algoritmo de búsqueda
basado en entornos, como una búsqueda local o un enfriamiento simulado. La hibridación
se lleva a cabo aplicando el método a cada hormiga de la colonia una vez ha construído su solución.

\subsection{ACO en MCP}

Al ser el problema del clique máximo un problema de grafos, los algoritmos de colonias
de hormigas son fácilmente implementables, y se pueden esperar buenos resultados en su
ejecución. Si bien en MCP no tratamos de recorrer un grafo, como se haría en el problema
del viajante de comercio, el concepto es adaptable para trabajar con conjuntos de candidatos
sobre los cuales realizamos una acción, como un \textit{add} o un \textit{swap}.

En este trabajo, he considerado dos algoritmos de colonias de hormigas clásicos,
lanzando una colonia de hormigas durante un número determinado de iteraciones y
reforzando el nivel de feromonas en las soluciones de calidad. La evaporación se
realiza reduciendo un porcentaje de feromona, fijo en todas las iteraciones.
La diferencia entre ambos algoritmos reside en la función de probabilidad;
mientras que uno de ellos utiliza solo las feromonas, el segundo complementa
el uso de las feromonas con información del grafo.

\subsection{Primer algoritmo}

El primer algoritmo implementa una colonia de hormigas básica, en la que un número
fijo de hormigas crean una solución cada una, y la mejor es la que se tiene en cuenta
para modificar los niveles de feromona. He considerado únicamente movimientos \textit{add},
por lo que las hormigas irán recorriendo los nodos del conjunto $C_0$ hasta que
este sea vacío.

Cada hormiga comenzará desde un clique vacío y añadirá nodos del conjunto $C_0$ según
una regla probabilística, que consiste en dar a cada nodo del conjunto una probabilidad
proporcional a su feromona, normalizándola para obtener un valor entre $0$ y $1$, y
conseguir que la suma de todas las probabilidades de los nodos de $C_0$ sea $1$.
Para ello,  solo hemos de dividir el valor de la feromona de un nodo por la suma de
los valores de los nodos de $C_0$. Inicialmente, la probabilidad será la misma para
cada nodo,  por lo que no damos ventaja a ninguno y dejamos que sea el proceso el
que determine los mejores nodos.

El proceso se repite hasta que la hormiga finaliza la creación de su clique, dando
paso a la siguiente, que crea uno distinto, y así hasta que todas las hormigas completen
el suyo. Cada vez que una hormiga finaliza, compara su clique con el mejor obtenido
en esa iteración, sustituyéndolo de ser necesario. Así, tendremos siempre el mejor
clique creado por la colonia a lo largo de cada iteración.

Una vez finalizado el trabajo de todas las hormigas, se produce la evaporación de
feromona, consistente en multiplicar los valores en cada nodo por una constante $\beta$.
Seguidamente, se aumentan los niveles de feromona en los nodos de la mejor solución,
multiplicando el valor de esta por un valor $\alpha$. Este valor dependerá del tamaño
del clique encontrado, siendo mayor cuanto más tamaño tiene el clique. Así, conseguimos
dar más peso a soluciones de más calidad.

El proceso se repite un número fijado de iteraciones. Una vez terminadas todas las
iteraciones, el algoritmo finaliza y devuelve la mejor solución encontrada.

Se han fijado los valores $\beta = 0.925$, $\alpha = (5 \times |C| + N) / N$, donde
$|C|$ es el tamaño del clique encontrado, y $N$ es el número de vértices del grafo.
Para fijarlos, me he valido de la experimentación, tomando aquellos que daban mejores
resultados en un conjunto reducido de instancias. No obstante, las diferencias
eran mínimas entre todos los valores probados.

El número de hormigas en cada iteración ha sido fijado a $50$, efectuándose un total
de $40$ iteraciones.

Vemos el pseudocódigo del algoritmo:

\begin{algorithm}[H]
\caption{ACO 1}
  \begin{algorithmic}
    \State $T_{\text{final}} = 0.01$
    \State Inicializar constantes $\alpha, \beta$.
    \State $Mejor Clique = []$.
    \Repeat
    \State $Hormigas = 0$.
      \State Establecer el mejor clique de la colonia al clique vacío.
      \Repeat
        \State $Clique = []$
        \State Añadir vértice aleatorio al clique
        \State Calcular $C_0$
        \Repeat
          \State Calcular probabilidades de los nodos en $C_0$. Normalizar.
          \State Elegir un nodo según la probabilidad y añadirlo a clique.
          \State Recalcular $C_0$.
        \Until{$C_0 = \emptyset$}
        \State Actualizar el mejor clique de la colonia si es necesario.
        \State $Hormigas = Hormigas + 1$.
      \Until{$Hormigas = $ límite hormigas}
      \State Actualizar $Mejor Clique$ si es necesario.
      \State Evaporar feromona.
      \State Aportar feromona a los nodos del mejor clique de la colonia.
      \State $i = i+1$.
    \Until{$i = limite$}
    \State{Devolver mejor clique encontrado.}
  \end{algorithmic}
\end{algorithm}


\subsection{Segundo algoritmo}

Este algoritmo toma algunas de las ideas de X. Xu, J. Ma y J. Lei, que trata de
mejorar el ACO básico considerando información relativa al clique para calcular
las probabilidades de elegir un nodo. La única diferencia con el anterior está
dicho cálculo, para el que utilizaremos una técnica de enfriamiento simulado,
además del uso de los valores de feromona.

La probabilidad se compone de dos partes: la primera de ellas es idéntica a la
anterior, un valor entre $0$ y $1$ obtenido de las feromonas de cada nodo en $C_0$,
normalizando para que la suma total sea $1$. Además, sumaremos un valor que depende
del grado de cada nodo, y estará multiplicado por una temperatura, de forma que la
importancia de este término desciende en el tiempo. La fórmula para el cálculo
de la probabilidad queda de la siguiente forma:
\[ p(v_i) = \frac{\tau(v_i)}{\sum_{v_j \in C_0} \tau(v_j)} + T \frac{D(v_i)}{N} \]
donde $\tau(v_i)$ es el valor de la feromona en el nodo $v_i$, $T$ es la temperatura,
$N$ el número de nodos del grafo, y $D(v_i)$ es el grado de $v_i$. Una vez calculadas
todas las probabilidades de los nodos de $C_0$, las normalizaremos para tener valores
entre $0$ y $1$ y que la suma siga siendo $1$.

La temperatura sigue un esquema de descenso geométrico similar al utilizado en los
algoritmos de enfriamiento simulado, en los que la temperatura se multiplica por
una constante $\gamma$ en cada iteración. Dicho valor se ha fijado a $\gamma = 0.95$.
Los valores de $\alpha$ y $\beta$ siguen siendo los mismos que en el algoritmo anterior,
al igual que el número de hormigas y de iteraciones, que son $50$ y $40$ respectivamente.

Puesto que solo cambia la función del cálculo de la probabilidad, no incluyo pseudocódigo
de este algoritmo, por ser similar al anterior.

\section{Algoritmos genéticos}

Los algoritmos genéticos entran dentro de la llamada \textbf{computación evolutiva},
que está compuesta por modelos de evolución basados en poblaciones. Se basan en
los procesos evolutivos que están presentes en la naturaleza, en los que dos padres
combinan sus genes para producir hijos (usualmente dos), los cuales presentan
características de ambos. Las diferencias entre generaciones pueden llevar a nuevas
configuraciones que devengan en una mejor adaptación al entorno. Esto, en el ámbito
que nos encontramos, se traduce en cambios que nos lleven mejores soluciones, pues
estas se adaptan mejor al problema.

En los algoritmos genéticos, las poblaciones estarán formadas por soluciones al
problema, que combinaremos dos a dos para producir descendientes que intenten mejorar
a sus progenitores. Esta combinación de soluciones se produce mediante el llamado
\textbf{operador de cruce}, que emula cómo los cromosomas se mezclan para producir
una nueva generación. Este operador dependerá del problema y no tiene por qué ser único.
El requisito fundamental es que se tome información de ambos progenitores para crear
a sus descendientes.

En la naturaleza, es posible que al combinar dos individuos se produzca una mutación,
con lo cual el hijo tendría una característica que no ha adquirido de ninguno de sus
padres, y que podría conllevar una ventaja o un perjuicio para el mismo. Esta técnica
se reproducirá en los algoritmos genéticos, incluyendo la posibilidad de que una
solución mute al ser generada. Al igual que el operador de cruce, las mutaciones
pueden seguir distintos enfoques, aunque se suele requerir que introduzcan modificaciones
significativas en los elementos de la población.

Controlaremos el cruce entre individuos y las mutaciones de sus descendientes con
parámetros que establezcan la probabilidad de que ambos sucesos ocurran. Para el
cruce, seleccionaremos dos individuos que o bien generarán descendientes, o bien
pasarán directamente a la siguiente generación. En caso de generar descendientes,
cada uno de ellos podrá ser mutado en el momento de su creación.

Existen distintas formas de seleccionar a los padres. El muestreo puede ser aleatorio,
se puede dar mayor probabilidad de combinarse a los mejores, etc. No hay restricción
en cuanto al número de veces que se puede reproducir un elemento, si bien podría
establecerse de considerarse necesario.

De igual forma, existen diversos métodos para crear una nueva generación. Una de
ellas sería que los hijos sustituyan a los padres, el llamado modelo generacional.
En este caso, se crea una nueva población en cada iteración del ciclo. Es posible
incluir elitismo, de forma que la mejor solución de la población pasa siempre a
la siguiente generación, lo que nos permite conservarla siempre sin necesidad de
almacenarla explícitamente. Otro modelo es el llamado estacionario, en el cual
se eligen dos padres, que generan dos hijos, y estos compiten para pasar a la
siguiente generación, pudiendo hacerlo con los padres o con la población en general.
Este modelo ya es elitista, luego no necesita introducir elitismo explícitamente.

\subsection{Genéticos en MCP}

Implementar un algoritmo genético para el problema del clique máximo se basa en
considerar operadores de cruce y mutación apropiados, esto es, que mantengan la
estructura de clique. A priori, no podemos afirmar que la combinación arbitraria
de dos soluciones o la mutación produzcan cliques, por lo que estos operadores
deben ser muy específicos o ser rectificados de alguna forma.

Una vez tengamos estos operadores, la ejecución del algoritmo consiste simplemente
en ir obteniendo nuevas generaciones hasta un criterio de parada determinado.
En este caso, para el algoritmo implementado, se ha establecido un número fijo de
iteraciones.

\subsection{Algoritmo implementado}

El algoritmo genético implementado se basa en un modelo estacionario, en el que
los dos descendientes compiten con sus padres para entrar en la nueva generación.
Al principio, se creará una población de soluciones de forma aleatoria mediante
un algoritmo \textit{greedy}, y a partir de aquí se irán generando nuevas soluciones.

El operador de cruce opera de la siguiente forma: dados dos padres, se toman aquellos
nodos comunes en ambos, y estos estarán presentes en ambos hijos. Seguidamente,
tomaremos el conjunto de nodos presentes en uno de los dos padres, y los distribuiremos
entre los dos hijos de forma aleatoria y equiprobable.

Una vez tengamos los hijos, procedemos a mutarlos con la probabilidad dada. El
algoritmo de mutación toma dos valores aleatorios entre $0$ y el número de nodos
del grafo, y cambia los nodos del intervalo que forman, quitándolos del hijo si
pertenecen a él, e introduciéndolos si no lo hacen.

Como no tenemos garantizado que los hijos sean cliques antes o después de la mutación,
los reduciremos a un clique eliminando nodos hasta que el conjunto restante sea un
clique. Una vez reducido, lo ampliaremos hasta que el clique sea maximal.
Ambos procesos de reducción y ampliación se hacen de forma aleatoria, si bien
podríamos haber tomado cualquier otro método basado en el información del grafo,
como ordenar por número de adyacencias.

La población inicial se ha tomado con un tamaño de $50$ individuos, con probabilidad
de cruce de $1$ y probabilidad de mutación de $0.1$, un valor alto para lo usual,
fijado así para introducir diversidad en la población. El criterio de parada es un
límite de iteraciones, que se ha fijado a $40$ en las ejecuciones realizadas.
Estos parámetros pueden ser modificados para evaluar la funcionalidad del algoritmo
bajo otras circunstancias.

[Pseudocódigo aquí]

\section{Algoritmos meméticos}

Los algoritmos meméticos son el resultado de la combinación de algoritmos meméticos
con búsqueda local. Tratan de aprovechar la buena capacidad de exploración de los
algoritmos genéticos, y combinarla con la explotación de soluciones de las técnicas
de búsqueda local.

El diseño de estos algoritmos queda muy abierto, y por lo general, suele ser necesario
adaptarlos de una forma u otra dependiendo del problema. Suelen ser algoritmos específicos,
que aprovechan alguna característica en concreto del problema.

Entre las diversas decisiones que hay que tomar, una de las más relevantes es el
proceso de búsqueda local a considerar y cuándo aplicarlo. Podemos elegir cualquier
técnica por trayectorias simples, desde una búsqueda local básica a una busqueda
tabú, y podemos aplicarlas a toda la población o a una parte, en distintas
fases del algoritmo.

\subsection{Meméticos en MCP}

Debido a que disponemos de técnicas de búsqueda local y algoritmos genéticos, la
implementación de un algoritmo memético consiste en tomar las decisiones correspondientes
con respecto a la hibridación de ambas técnicas.

\subsection{Algoritmo implementado}

En este trabajo, he optado por una combinación básica de dos algoritmos ya implementados,
el algoritmo genético y la búsqueda local. De esta forma, se mantiene la estructura
de algoritmo genético estacionario, introduciendo una búsqueda local al final de
la creación de la nueva generación, lo que nos permite explorar su entorno antes de
pasar a la siguiente.

El desarrollo del algoritmo es muy simple, pues utiliza como base el algoritmo
genético introducido anteriormente. Utiliza el mismo operador de cruce y mutación,
y añade una búsqueda local después de que los hijos sean reducidos y ampliados.

He utilizado la segunda búsqueda local introducida, \textit{DLS}, debido a su menor
complejidad computacional, pues hibridar con la otra búsqueda local producía tiempos
de ejecución demasiado elevados.

Como parámetros, he tomado las mismas probabilidades de cruce y mutación, $1$ y $0.1$,
y he establecido el límite de la búsqueda local en un octavo del número de nodos
del grafo, para explorar únicamente la zona del espacio de soluciones más cercana
a cada elemento de la población.


%************************************************
\chapter{Implementación}\label{ch:implementacion}
%************************************************

% Hablar de la implementación de los algoritmos, involucrando a Ruby.

En este capítulo vamos a tratar la implementación del programa encargado de
aplicar los algoritmos al conjunto de instancias. Comenzaremos introduciendo
el lenguaje de programación utilizado, siguiendo con una descripción de la
funcionalidad completa del programa, y concluyendo con un desglose de carácter
más técnico de las partes que lo componen.

\section{Lenguaje utilizado}

La implementación del código de este trabajo ha sido realizada en el lenguaje de
programación \textbf{Ruby}, usando su versión 2.3.0. Ruby es un lenguaje orientado
a objetos, interpretado, con una sintaxis muy similar a la de \textbf{Python}.
Está disponible para descarga bajo una licencia de software libre.

He elegido este lenguaje por ser capaz de aportar simplicidad e intuitividad a
la hora de programar los algoritmos. Como ya había trabajado con este lenguaje en
ocasiones anteriores, conocía sus puntos fuertes, y he tratado de explotarlos
lo máximo posible, para que el resultado sea un código legible y adaptable.

%Además, sale de los lenguajes más utilizados en la carrera, lo que me permite incluir una herramienta menos usual.

Al utilizar un lenguaje interpretado como es Ruby, cabe esperar unos
tiempos de ejecución mayores que con otros lenguajes, debido a que no está hecho
para ejecutar código lo más rápido posible. En el caso de querer realizar un programa
capaz de ejecutar los algoritmos en el menor tiempo posible, lenguajes compilados
como \textbf{C++} serían más útiles que Ruby.

\section{Funcionalidad}

El código implementado cumple varias funciones: la primera de ellas es llevar a
memoria todas las instancias disponibles del problema. Para ello, he tenido que
implementar un lector de ficheros que cree una matriz de adyacencia para cada archivo
disponible. Estos archivos siguen una estructura determinada, que permite
automatizar dicho proceso, no teniendo que considerar casos independientes.
Cada archivo dará lugar a un problema, compuesto de una matriz de adyacencias,
el número de nodos y aristas involucrados, y el nombre del grafo en cuestión.

El segundo punto, y más importante, es la implementación de las heurísticas.
Cada una de ellas está contenida en una clase propia, que contendrá uno o varios
métodos para la resolución de los problemas generados a partir de los archivos.
Adicionalmente, podrán contener métodos auxiliares para uso propio o de otras
heurísticas, en caso de que utilicemos técnicas híbridas.

Finalmente, existe un conjunto de funciones auxiliares, que pueden ser
requeridas por varias heurísticas. La idea bajo su creación es encapsular
operaciones importantes, que serán repetidas frecuentemente y en diversos
escenarios, evitando repetición de código y sirviendo de ayuda para facilitar
el trabajo y para la comprensión del código.


\section{Inicialización y ejecución del programa}

Antes de aplicar ningún algoritmo al conjunto de instancias, hay que resolver la
lectura de los archivos y la creación de las estructuras de datos correspondientes
para cada grafo. Estas operaciones se engloban en un método llamado \textit{readerMethod},
que se encuentra en el archivo \textit{reader.rb}.

Cada instancia del problema dará lugar a un objeto de la clase \textbf{Problem}, que
almacena los datos de cada grafo: su nombre, número de vértices, número de aristas,
su matriz de adyacencias y el orden de cada nodo, esto es, su número de adyacencias.

El método de lectura de los archivos recorre el fichero en el que están almacenadas
las instancias del problema, y por cada uno de ellos construye un objeto de la clase
Problem. Dado que los ficheros siguen todos el mismo formato, es posible extraer
los datos de cada grafo sin más que hacer una distinción de casos.

La parte fundamental de la lectura, que es la creación de la matriz de adyacencias,
se realiza partiendo de una matriz de ceros del tamaño adecuado, e iremos tomando
los pares que se proporcionan en el archivo, que representan los nodos que están
conectados, para marcarlos como unos en la matriz de adyacencias. Dado que los
nodos comienzan su numeración desde el $1$ en los ficheros, y desde el $0$ en el
programa (debido a que las estructuras de datos comienzan en el $0$), hemos de
corregirlo a la hora de marcar las aristas en la matriz de adyacencias. Además,
marcaremos todos los nodos como conectados con ellos mismos, pues será necesario
para el correcto funcionamiento de algunas operaciones. Esto hace que todos los
nodos de un clique pertenezcan a su conjunto $C_0$ asociado, por lo que cuando lo
calculemos tendremos que quitar los nodos del propio clique.

El número de adyacencias de cada nodo no lo proporciona el archivo del grafo,
sino que hemos de calcularlo. Para ello, nos valemos de una función auxiliar
programada con tal fin.

El archivo \textit{reader.rb} también contendrá el método principal. En él, se
ejecutarán el método de lectura y los algoritmos implementados que incluyamos,
para lo que se debe añadir el código necesario. Además, se mostrarán por pantalla
los resultados obtenidos, con la idea de poder redireccionarlos a un archivo de
texto en el que almacenar los resultados.


\section{Algoritmos}

Todos los algoritmos implementados en el programa se encuentran recogidos en
un mismo módulo, el módulo \textbf{Clique}, bajo el que también se encuentra el
archivo \textit{reader.rb} detallado con anterioridad. Cada tipo de algoritmo
estará dentro de una clase, y se separarán los distintos algoritmos del mismo
tipo en métodos diferentes.

\subsection{Consideraciones}

Algunos algoritmos usan aleatoriedad en una o varias fases. Para conseguirla, se
ha utilizado la clase \textbf{Random}, que proporciona un generador de números
aleatorios. El método \textit{rand} nos permite obtener un valor aleatorio
dentro de un rango dado, pudiendo ser un número entero o real, lo que nos permite
cubrir todas las necesidades en este aspecto.

% Hablar de que no se usa información del grafo para construir una clase.

\subsection{Algoritmos \textit{greedy}}

Los dos algoritmos \textit{greedy} implementados se encuentran en el archivo
\textit{greedy.rb}, dentro de la clase \textbf{Greedy}. El algoritmo \textit{greedy}
básico está en el método \textit{solve}, mientras que el \textit{greedy} adaptativo
se encuentra en el método \textit{solve2}. Ambos métodos reciben únicamente una
instancia de la clase Problem como parámetro, el problema a resolver.

Además, en este archivo también se encuentran los distintos algoritmos voraces
encargados de completar cliques o de reparar grafos, utilizados en otras heurísticas.
Son los siguientes:

\begin{itemize}
  \item \textit{complete\_clique}, método que amplía un clique hasta que sea posible,
        tomando de entre los candidatos aquel con más adyacencias.
  \item \textit{complete\_clique\_random}, que, al igual que el anterior, amplía
        un clique hasta que sea maximal, solo que en este caso elige de forma aleatoria.
  \item \textit{solve\_random}, método que crea un clique partiendo desde cero, usando
        \textit{complete\_clique\_random}. Lo usaremos para generar soluciones iniciales
        aleatorias.
  \item \textit{repair}, método que elimina nodos de un grafo hasta que sea un clique.
        Elimina el nodo con menor número de adyacencias.
  \item \textit{repair\_random}, con la misma función que el anterior, pero en este
        caso elimina un nodo aleatorio.
\end{itemize}

\subsection{Búsqueda local}

Los dos algoritmos de búsqueda local están en el archivo \textit{localsearch.rb},
bajo la clase \textbf{LocalSearch}. Para la creación de la clase, solo es necesario
crear un generador de números aleatorio, por lo que no hay que pasarle argumentos
al constructor.

La heurística 1LS está implementada en el método \textit{solve\_with\_solution}, y
recibe como parámetro el elemento de la clase Problem a resolver, una solución
de partida, y un límite de cambios.

Por otra parte, la heurística DLS se encuentra en el método \textit{solve\_dynamic}.
Los parámetros serán los mismos que en el caso anterior: la instancia de la clase
Problem a resolver, la solución de partida y el límite de cambios.

Para la ejecución de cualquiera de los dos métodos partiendo desde el clique
vacío, existe el método \textit{solve}, que será el que llamemos cuando queramos
ejecutar uno de los dos algoritmos. Recibe como parámetros el problema que resolver
y el número de cambios, que los pasará al algoritmo que ejecute, además del clique
vacío como solución de partida.

\subsection{Enfriamiento simulado}

En el archivo \textit{simulatedannealing.rb} se encuentran los dos algoritmo de
enfriamiento simulado considerados, en la clase \textbf{SA}. El constructor de
clase no necesita parámetros, pues crea un generador de números aleatorios y
una instancia de la clase Greedy, ya que el algoritmo de Geng et al. necesita
las funciones para reparar y completar grafos.

El algoritmo de enfriamiento simulado básico está implementado en el método
\textit{solve}, que recibe como único parámetro el problema a resolver. Para el
criterio de aceptación del algoritmo Metrópolis, genera un número aleatorio entre
$0$ y $1$, y si es menor que el valor de la probabilidad, acepta la solución.
Para tomar elementos aleatorios del entorno, se utiliza la función \textit{shuffle},
que baraja los elementos del entorno, y toma un número aleatorio como semilla.

El algoritmo de Geng et al. se encuentra en el método \textit{solve2}, que vuelve
a tener como único parámetro el problema que se quiere resolver. Efectúa las mismas
operaciones para el criterio de Metrópolis y para el baraje del entorno.

Ambas funciones declaran los valores de las temperaturas y del parámetro $\beta$
al inicio. Para cambialos, basta con modificar dichos valores en el código.

\subsection{GRASP}

El algoritmo GRASP está implementado en el archivo \textit{grasp.rb}, dentro de la
clase \textbf{GRASP}. Además de un generador de números aleatorios, el constructor
declara un objeto de la clase LocalSearch, que necesitaremos para aplicar el
algoritmo de búsqueda local que consideremos.

El método \textit{generate\_random\_solution} es el encargado de generar las soluciones
iniciales pseudoaleatorias, tomando un problema como argumento. El tamaño de la
lista restringida de candidatos está fijado por un valor en el código, que debemos
cambiar en caso de querer aumentar o reducir el tamaño de dicha lista.

Como los dos algoritmos implementados solo difieren en la búsqueda local utilizada,
es posible implementar ambos utilizando un único método, y cambiar la búsqueda local
para alternar entre uno u otro. El método \textit{solve} es el encargado de esto,
recibiendo un problema y un número de iteraciones como parámetros. El número de
cambios para el algoritmo de búsqueda local está declarado como una variable, que
deberemos modificar si queremos variar su valor.

\subsection{ILS}

El algoritmo ILS se encuentra en el archivo \textit{ils.rb}, en la clase \textbf{ILS}.
Su construcción no necesita parámetros, y se encarga de crear el generador de números
aleatorios y un objeto de la clase LocalSearch para las búsquedas locales.

La clase tiene un único método, \textit{solve}, que es el encargado de aplicar la
búsqueda local y las permiutaciones a las soluciones obtenidas. Recibe el problema
a resolver y el número de iteraciones. El límite para la búsqueda local es una
variable, que modificaremos de ser necesario. Al igual que en el caso de GRASP,
para cambiar la búsqueda local que aplicamos basta con cambiar la llamada al método
de búsqueda local.

\subsection{Algoritmos de colonia de hormigas}

Los dos algoritmos de colonia de hormigas implementados están en el archivo \textit{aco.rb},
dentro de la clase \textbf{ACO}. El constructor no recibe ningún parámetro, pues solo
crea un generador de números aleatorios.

Mi versión del algoritmo está en el método \textit{solve}, que recibe el problema que
se quiere resolver y el número de iteraciones. El número de hormigas y el parámetro
$\beta$, necesario para la evaporación de feromonas, se encuentran en el código,
pudiendo cambiarse su valor de ser necesario.

Para elegir cada nodo al que va una hormiga, una vez calculadas las probabilidades,
se genera un valor aleatorio entre $0$ y $1$, y vamos tomando nodos, acumulando su
valor de probabilidad, hasta que dicho valor supere el valor aleatorio generado.

El algoritmo de Xu et al. se encuentra en el método \textit{solve2}, y, debido a que
su estructura es igual que el algoritmo anterior, salvo el cálculo de las probabilidades,
se le puede aplicar todo lo mencionado anteriormente.

\subsection{Algoritmos genéticos}

Tanto el algoritmo genético como el memético implementados se encuentran dentro de
la clase \textbf{Genetics}, en el archivo \textit{genetics.rb}. El constructor
de la clase no recibe ningún parámetro, y crea un generador de números aleatorios,
un elemento de la clase LocalSearch y uno de la clase Greedy, que se usarán en el
algoritmo memético y para completar o reparar las soluciones, respectivamente.

El algoritmo genético de Zhang et al. está implementado en el método \textit{solve},
que recibe el problema a resolver y el número de iteraciones. Tanto las probabilidades
de cruce y mutación como el tamaño de la población se encuentran en el código, donde
deben ser modificadas de ser necesario. Como en ocasiones anteriores, para decidir
eventos que dependen de una probabilidad, se genera un número entre $0$ y $1$, y si
este es menor que la probabilidad, el suceso se lleva a cabo. Este proceso se lleva
a cabo en los mecanismo de cruce y de mutación, y además se usa también en el propio
operador de cruce, donde los nodos que pertenecen solo a un padre van a cada uno de
los hijos con probabilidad $0.5$.

El algoritmo memético, implementado en \textit{solve\_memetic}, funciona de la misma
forma, aplicando búsqueda local al generar la nueva población, con un límite que
también está en el código. Por lo demás, opera  igual que el genético, por lo que
son aplicables las consideraciones anteriores.

Además, existe un método privado para las mutaciones, llamado \textit{mutation}.
Recibe como parámetro la solución a mutar y la matriz de adyacencias, y devuelve
la solución mutada.


\section{Métodos auxiliares}

Además de los algoritmos, también se han implementado funciones que realizan
cálculos necesarios en varios de los algoritmos, que han sido incluidas en un
archivo adicional creado para almacenarlas. Estas funciones no se encuentran
dentro de una clase, pero sí que están en un módulo distinto, el módulo \textbf{Algorithm},
para diferenciarlas de las implementaciones de las heurísticas. Todas estas funciones
se encuentran en el archivo \textit{algorithm.rb}, y las detallamos a continuación:

\begin{itemize}
  \item \textit{add}: Operador \textit{add}. Recibe un clique y un vértice, y añade
        el vértice al clique. Supone que el vértice no está en el clique y que
        se mantiene la estructura.

  \item \textit{swap}: Operador \textit{swap}. Recibe un clique, un vértice $v$ del conjunto
        $C_1$, y la matriz de adyacencias. Busca el nodo dentro del clique que no está
        conectado a $v$ y lo intercambia con $v$. De nuevo, suponemos que $v \in C_1$,
        por lo que se puede hacer el intercambio.

  \item \textit{swap\_two}: Operador \textit{swap}. En este caso, se pasan como parámetros
        el clique y los dos nodos a intercambiar, $p$ en $C_1$ y $q$ dentro del clique.
        El método elimina $q$ del clique y mete $p$, suponiendo que se verifica lo
        anterior y que $q$ es el único nodo del clique no conectado a $p$.

  \item \textit{drop}: Operador \textit{drop}. Recibe un clique y un nodo, y elimina
        dicho nodo del clique.

  \item \textit{adjacencies}: Método que recibe un nodo y una matriz de adyacencias,
        y devuelve el número de adyacencias de dicho nodo. Se utiliza para calcular
        las adyacencias de todos los nodos al crear un problema.

  \item \textit{connections}: Dados un nodo $n$, un conjunto de nodos $S$ y una matriz de
        adyacencias, este método devuelve el número de nodos de $S$ conectados con $c$.

  \item \textit{is\_connected}: Recibe un nodo $n$, un conjunto de nodos $S$ y una matriz
        de adyacencias, y devuelve si $n$ está conectado con todos los nodos de $S$.
        Este método nos servirá para calcular el conjunto $C_0$.

  \item \textit{is\_connected\_but\_one}: Funciona igual que el anterior, pero comprueba
        que el nodo $n$ esté conectado con todos los nodos de $S$ salvo uno. Nos servirá
        para el cálculo del conjunto $C_1$.

  \item \textit{connected}: Dados un nodo y una matriz de adyacencias, devuelve todos los
        nodos del grafo conectados a él.

  \item \textit{connected\_with\_all}: Dado un conjunto de nodos $S$ y una matriz de adyacencias,
          usando el método \textit{is\_connected}, devuelve el conjunto $M$, formado por los nodos
        conectados con todos los de $S$, y elimina el conjunto $S$ de $M$, pues si $S$ es un clique,
        $S \subseteq M$. Obtenemos por tanto el conjunto $C_0$, que será para lo que utilizemos
        este método.

  \item \textit{missing\_one\_connection}: Dado un conjunto de nodos $S$ y una matriz de adyacencias,
        devuelve el conjunto de nodos conectados con todos los de $S$ salvo uno. Lo usaremos
        para calcular el conjunto $C_1$ asociado a un clique.

  \item \textit{one\_connected\_with\_all}: Recibe una lista y una matriz de adyacencias, y
        devuelve \textit{true} si existe un nodo fuera de la lista conectado a todos los
        nodos de ella.

  \item \textit{value} y \textit{value2}: Funciones objetivo utilizadas en los algoritmos
        de enfriamiento simulado. Ambas reciben un grafo y una matriz de adyacencias, y
        devuelven el valor calculado.

  \item \textit{is\_clique}: Función que recibe un grafo y una matriz de adyacencias, y
        devuelve si dicho grafo es un clique o no.

  \item \textit{print\_solution}: Método para mostrar por pantalla una solución. Primero,
        comprueba si es un clique con el función anterior y lo muestra. Después, imprime
        su longitud. Además, imprime la solución y su longitud por la salida de error,
        lo que permite redirigirlo a un archivo distinto a la hora de ejecutar el programa.

\end{itemize}


%************************************************
\chapter{Resultados y análisis}\label{ch:resultados}
%************************************************

En este capítulo expondremos los resultados obtenidos por cada algoritmo sobre el
conjunto de cliques que se ha considerado, y evaluaremos el funcionamiento individual
de cada uno de ellos. Además, compararemos los algoritmos del mismo tipo para ver cómo
afectan sus diferencias en los resultados.

\section{Consideraciones}

Para la toma de datos, he utilizado un portátil con el sistema operativo Antergos
Linux 64-bit, con un procesador Intel Core i7-3610QM, que cuenta con $4$ núcleos
y $2$ hebras por núcleo, con una frecuencia de reloj de $2.30$ gigahercios. Los
algoritmos han sido ejecutados individualmente, sin uso de paralelismo, por lo que
cada ejecución ha ocupado una sola hebra.

Algunos de los algoritmos implementados en este trabajo cuentan con una o varias
componentes aleatorias: elección de un elemento aleatorio, cálculo de un valor para
aceptar o no una acción que depende de una probabilidad, etcétera. Debido a esto,
para la obtención de resultados, se han realizado varias ejecuciones de estos algoritmos,
de las cuales se ha tomado el mejor valor obtenido. He tomado el mejor, y no la media,
porque representa el tamaño del clique que dicho algoritmo puede alcanzar en un
número razonable de ejecuciones.

En el caso de los algoritmos que no tienen componentes aleatorias, como su ejecución
lleva siempre a la misma solución, el programa se ha ejecutado una única vez en
cada caso. Estos algoritmos son los dos \textit{greedy} y 1LS. Además, los algoritmos
ILS y GRASP que usan 1LS también se han ejecutado una única vez, aunque el método
de obtención de soluciones iniciales tenga componentes aleatorias, debido a que su
tiempo de ejecución es muy alto.

El resto de algoritmos incluyen componentes aleatorias, y se ha considerado un
número variable de ejecuciones para cada uno de ellos, dependiendo del tiempo que
tardasen. Se detallan a continuación:
\begin{itemize}
  \item DLS: 20 iteraciones.
  \item Enfriamiento simulado: 10 iteraciones para cada algoritmo.
  \item Colonia de hormigas: 10 iteraciones para cada algoritmo.
  \item ILS y GRASP con DLS: 10 iteraciones.
  \item Genético: 10 iteraciones.
  \item Memético: 5 iteraciones.
\end{itemize}

\section{Análisis}

A continuación, vamos a ver los resultados obtenidos por cada algoritmo en forma de tabla,
para hacer un análisis de cada uno de ellos. Debido a la cantidad de algoritmos, estos
han sido divididos según su tipo, en grupos de dos: dos \textit{greedy}, dos búsquedas
locales, dos enfriamientos simulados, dos ILS, dos GRASP, dos algoritmos de colonia de
hormigas y los algoritmos genético y memético. Para cada grupo, mostraremos los resultados
y haremos un análisis individual y comparativo de los dos algoritmos que contiene.

Las tablas tienen la siguiente disposición: la primera columna contiene cada uno de los
grafos considerados. La segunda columna indica el tamaño del mejor clique conocido para
cada grafo, notando con un asterisco los valores que están confirmados como óptimos
globales. El resto de columnas contienen el tamaño del mejor clique encontrado para cada
algoritmo, notado como |C|, y el tiempo de ejecución, que se ha tomado el valor medio en
caso de haberse efectuado más de una ejecución. Como medidas para evaluar a los grafos,
he considerado el porcentaje medio de tamaño de los cliques encontrados, calculado como
la media aritmética de los porcentajes en cada grafo, y el número de óptimos encontrados.

\subsection{Algoritmos \textit{greedy}}

\begin{small}
\begin{longtable}{l l l l l l}
  \label{table:greedy}\\
  & & \multicolumn{2}{c}{Greedy básico} & \multicolumn{2}{c}{Greedy adaptativo} \\ \cline{3-6}
  Grafo              & Clique & |C| & Tiempo (s) & |C| & Tiempo (s) \\ \hline
  \endhead
  \endfoot
  brock200\_2        & 12 & 7 & < 0.01 & 9 & 0.18 \\ \hline
  brock200\_4        & 17 & 13 & 0.02 & 15 & 0.27 \\ \hline
  brock400\_2        & 29 & 20 & 0.08 & 22 & 1.61 \\ \hline
  brock400\_4        & 33 & 18 & 0.07 & 21 & 1.58 \\ \hline
  brock800\_2        & 24 & 14 & 0.18 & 19 & 4.98 \\ \hline
  brock800\_4        & 26 & 14 & 0.18 & 19 & 4.41 \\ \hline
  C125.9             & 34* & 29 & 0.02 & 32 & 0.44 \\ \hline
  C250.9             & 44* & 35 & 0.08 & 40 & 1.49 \\ \hline
  C500.9             & 57 & 43 & 0.31 & 51 & 24.10 \\ \hline
  C1000.9            & 68 & 51 & 1.16 & 56 & 126.73 \\ \hline
  C2000.5            & 16* & 10 & 0.72 & 13 & 24.73 \\ \hline
  C2000.9            & 80 & 56 & 4.39 & 66 & 1058.85 \\ \hline
  C4000.5            & 18* & 12 & 2.86 & 15 & 3.64 \\ \hline
  DSJC500\_5         & 13 & 8 & 0.05 & 11 & 0.07 \\ \hline
  DSJC1000\_5        & 15 & 10 & 0.19 & 12 & 6.18 \\ \hline
  gen200\_p0.9\_44   & 44 & 32 & 0.06 & 36 & 1.81 \\ \hline
  gen200\_p0.9\_55   & 55 & 36 & 0.06 & 39 & 1.56 \\ \hline
  gen400\_p0.9\_55   & 55 & 44 & 0.22 & 43 & 12.67 \\ \hline
  gen400\_p0.9\_65   & 65 & 40 & 0.21 & 44 & 12.82 \\ \hline
  gen400\_p0.9\_75   & 75 & 45 & 0.23 & 46 & 13.36 \\ \hline
  hamming10-4        & 40 & 32 & 0.70 & 36 & 120.37 \\ \hline
  hamming8-4         & 16 & 16 & 0.02 & 16 & 0.57 \\ \hline
  keller4            & 11 & 8 & < 0.01 & 11 & 0.25 \\ \hline
  keller5            & 27 & 17 & 0.29 & 24 & 45.48 \\ \hline
  keller6            & 59 & 37 & 9.39 & 48 & 5575.49 \\ \hline
  MANN\_a27          & 126 & 125 & 1.03 & 125 & 17.20 \\ \hline
  p\_hat300-1        & 8 & 7 & 0.01 & 8 & 0.02 \\ \hline
  p\_hat300-2        & 25 & 23 & 0.05 & 23 & 0.83 \\ \hline
  p\_hat300-3        & 36 & 30 & 0.08 & 34 & 2.68 \\ \hline
  p\_hat700-1        & 11 & 7 & 0.007 & 8 & 0.10 \\ \hline
  p\_hat700-2        & 44 & 38 & 0.32 & 43 & 7.51 \\ \hline
  p\_hat700-3        & 62* & 55 & 0.59 & 58 & 10.20 \\ \hline
  p\_hat1500-1       & 12 & 8 & 0.32 & 11 & 0.44 \\ \hline
  p\_hat1500-2       & 65* & 54 & 1.65 & 62 & 42.14 \\ \hline
  p\_hat1500-3       & 94* & 75 & 2.88 & 87 & 94.61 \\ \hline
  \caption{Resultados en algoritmos greedy.}
\end{longtable}
\end{small}

Los resultados obtenidos por los dos algoritmos \textit{greedy} han sido de un óptimo y una media
de $73.2\%$ de tamaño para el \textit{greedy} básico, y tres óptimos y un $84.9\%$ de media
para el \textit{greedy} adaptativo.

Los resultados del algoritmo \textit{greedy} básico muestran que este falla a la hora de
alcanzar el óptimo en el conjunto de instancias, pues lo logra en solo una de las $35$.
El tamaño medio manifiesta que no se obtienen soluciones de calidad, existiendo
un amplio margen de mejora en la mayoría de instancias del problema. En cuanto a
los tiempos, es un algoritmo rápido, que por lo general requiere menos de $1$
segundo para ejecutarse, y que no supera los $10$ segundos en ningún caso.

El \textit{greedy} adaptativo logra mejorar o igualar los resultados obtenidos por el algoritmo
anterior en todos los grafos, lo que confirma que ampliar el entorno con movimientos
\textit{swap} mejora sensiblemente los resultados, aunque aumenta de forma considerable
los tiempos de ejecución, llegando a superar la hora en \textit{keller6}. Aun así,
en la mayoría de los grafos, la ejecución es rápida, no llegando a superar los $60$
segundos. Si nos fijamos en sus resultados por separado, el porcentaje medio de tamaño
de los cliques hallados y la obtención de únicamente $3$ óptimos, hacen que este
algoritmo tampoco consiga buenos resultados si se aplica sin combinar con ningún otro.

Estos resultados muestran que los algoritmos voraces implementados no
son capaces de alcanzar buenas soluciones por sí solos. Aun así, los cliques
obtenidos pueden servir como soluciones de partida para otros algoritmos que sean
capaces de mejorar las carencias que tienen los \textit{greedy}. En este caso,
si trabajamos con cliques de gran tamaño nos interesa un algoritmo más veloz como
es el \textit{greedy} básico. Si por el contrario los grafos son pequeños, a cambio de algo
más de tiempo de ejecución, el algoritmo \textit{greedy} adaptativo nos da soluciones de
partida de más calidad.

\subsection{Búsqueda local}

\begin{small}
  \begin{longtable}{l l l l l l}
  \label{table:bl}\\
    & & \multicolumn{2}{c}{1LS} & \multicolumn{2}{c}{DLS} \\ \cline{3-6}
    Grafo              & Clique & |C| & Tiempo (s) & |C| & Tiempo (s) \\ \hline
    \endhead
    \endfoot
    brock200\_2        & 12 & 9 & 0.36 & 11 & < 0.01   \\ \hline
    brock200\_4        & 17 & 15 & 0.37 & 15 & 0.01   \\ \hline
    brock400\_2        & 29 & 22 & 2.05 & 23 & 0.06   \\ \hline
    brock400\_4        & 33 & 23 & 1.85 & 23 & 0.07   \\ \hline
    brock800\_2        & 24 & 19 & 6.82 & 19 & 0.10   \\ \hline
    brock800\_4        & 26 & 20 & 6.13 & 18 & 0.09   \\ \hline
    C125.9             & 34* & 33 & 0.36 & 34 & 0.04   \\ \hline
    C250.9             & 44* & 42 & 1.91 & 41 & 0.14   \\ \hline
    C500.9             & 57 & 54 & 7.69 & 54 & 0.39   \\ \hline
    C1000.9            & 68 & 63 & 32.02 & 62 & 1.62   \\ \hline
    C2000.5            & 16* & 13 & 46.60 & 14 & 0.19  \\ \hline
    C2000.9            & 80 & 72 & 179.76 & 68 & 3.63  \\ \hline
    C4000.5            & 18* & 15 & 141.92 & 15 & 0.48   \\ \hline
    DSJC500\_5         & 13 & 11 & 1.89 & 12 & 0.03   \\ \hline
    DSJC1000\_5        & 15 & 13 & 8.28 & 13 & 0.05   \\ \hline
    gen200\_p0.9\_44   & 44 & 39 & 1.01 & 43 & 0.08   \\ \hline
    gen200\_p0.9\_55   & 55 & 39 & 1.01 & 55 & 0.12   \\ \hline
    gen400\_p0.9\_55   & 55 & 50 & 5.03 & 49 & 0.35   \\ \hline
    gen400\_p0.9\_65   & 65 & 45 & 9.12 & 49 & 0.37   \\ \hline
    gen400\_p0.9\_75   & 75 & 49 & 5.83 & 75 & 0.37   \\ \hline
    hamming10-4        & 40 & 36 & 19.48 & 40 & 0.68   \\ \hline
    hamming8-4         & 16 & 16 & 0.50 & 16 & < 0.01   \\ \hline
    keller4            & 11 & 11 & 0.22 & 11 & 0.01  \\ \hline
    keller5            & 27 & 23 & 17.27 & 27 & 0.22   \\ \hline
    keller6            & 59 & 51 & 556.38 & 50 & 4.97   \\ \hline
    MANN\_a27          & 126 & 125 & 62.29 & 123 & 3.84   \\ \hline
    p\_hat300-1        & 8 & 7 & 0.55 & 7 & < 0.01   \\ \hline
    p\_hat300-2        & 25 & 25 & 1.35 & 25 & 0.05   \\ \hline
    p\_hat300-3        & 36 & 34 & 1.67 & 36 & 0.13  \\ \hline
    p\_hat700-1        & 11 & 9 & 2.76 & 9 & 0.02   \\ \hline
    p\_hat700-2        & 44 & 44 & 11.17 & 43 & 0.39   \\ \hline
    p\_hat700-3        & 62* & 60 & 28.40 & 61 & 0.76  \\ \hline
    p\_hat1500-1       & 12 & 11 & 17.38 & 11 & 0.08   \\ \hline
    p\_hat1500-2       & 65* & 64 & 132.85 & 63 & 1.85   \\ \hline
    p\_hat1500-3       & 94* & 91 & 128.08 & 91 & 4.61   \\ \hline
  \caption{Resultados en búsqueda local.}
\end{longtable}
\end{small}

Los dos algoritmos de búsqueda local han dado unos resultados de $4$ óptimos y
un $87.5\%$ de media para 1LS, y $9$ óptimos y un $91.0\%$ de media para DLS.

Vemos como 1LS consigue mejorar los resultados obtenidos por ambos algoritmos
\textit{greedy} tanto en número de óptimos como de tamaño medio de los grafos,
lo que reafirma que la exploración del entorno es importante a la hora de buscar
soluciones. No obstante, en vista de los resultados, estos siguen siendo mejorables
en cuanto a la calidad de las soluciones encontradas. Los tiempos de ejecución no
son excesivos, si bien en los grafos más grandes pueden extenderse un poco,
rozando los $10$ minutos en \textit{keller6}. Aun así, en la mayoría de los grafosgreedy
no se superan los $30$ segundos.

DLS logra mejorar los resultados de 1LS, pues es capaz de alcanzar más óptimos
y de obtener un mayor tamaño medio, que en este caso supera el $90\%$, por lo
que tenemos un algoritmo que empieza a obtener soluciones satisfactorias.
Los tiempos de ejecución son significativamente menores, lo cual era de esperar,
ya que DLS, al contrario que ILS, no hace cálculos para la elección de candidatos,
sino que dicha elección es aleatoria, lo que reduce la complejidad del algoritmo.
Todas las ejecuciones se completan en menos de $5$ segundos, la mayoría de ellas
requiriendo menos de $1$ segundo.

Debido a la mejora en tiempos de ejecución y en resultados, podemos afirmar que
DLS funciona mejor que 1LS al aplicarlo directamente sobre los grafos, y cabe esperar
que también ofrezca mejores resultados cuando lo utilicemos en ILS y GRASP. Aun así,
exiten grafos, como \textit{C2000.9} o \textit{MANN\_a27} en los que 1LS proporciona
resultados ligeramente mejores que DLS. Esto puede indicar que en grafos con una
determinada estructura, 1LS pueda ser de más utilidad que DLS.

Recordando el funcionamiento de los algoritmos, 1LS ordenaba los tres movimientos
para obtener el entorno y trataba de aprovechar información del grafo para elegir
entre los nodos candidatos, mientras que DLS no consideraba movimientos \textit{drop},
y las elecciones eran aleatorias. Los resultados parecen apuntar a que el entorno
considerado en DLS da mejores resultados que el tomado en 1LS. Además, el uso de
información inherente al grafo no parece aportar beneficio, si bien sería necesaria
más experimentación para poder afirmarlo.

\subsection{Enfriamiento simulado}

\begin{small}
\begin{longtable}{l l l l l l}
  \label{table:sa}\\
    & & \multicolumn{2}{c}{SA básico} & \multicolumn{2}{c}{SA adaptado} \\ \cline{3-6}
    Grafo              & Clique & |C| & Tiempo (s) & |C| & Tiempo (s) \\ \hline
    \endhead
    \endfoot
    brock200\_2        & 12 & 10 & 1.25 & 9 & 10.00 \\ \hline
    brock200\_4        & 17 & 16 & 2.15 & 13 & 9.67 \\ \hline
    brock400\_2        & 29 & 24 & 7.73 & 18 & 24.15 \\ \hline
    brock400\_4        & 33 & 24 & 5.55 & 19 & 23.30 \\ \hline
    brock800\_2        & 24 & 19 & 11.19 & 16 & 54.90 \\ \hline
    brock800\_4        & 26 & 20 & 11.13 & 16 & 43.35 \\ \hline
    C125.9             & 34* & 34 & 2.06 & 32 & 6.31 \\ \hline
    C250.9             & 44* & 44 & 8.78 & 36 & 14.42 \\ \hline
    C500.9             & 57 & 53 & 22.86 & 45 & 27.79 \\ \hline
    C1000.9            & 68 & 63 & 44.07 & 52 & 57.76 \\ \hline
    C2000.5            & 16* & 15 & 29.26 & 12 & 136.12 \\ \hline
    C2000.9            & 80 & 71 & 83.85 & 57 & 133.82 \\ \hline
    C4000.5            & 18* & 15 & 44.58 & 13 & 220.79 \\ \hline
    DSJC500\_5         & 13 & 12 & 4.10 & 10 & 28.73 \\ \hline
    DSJC1000\_5        & 15 & 14 & 9.28 & 11 & 60.67 \\ \hline
    gen200\_p0.9\_44   & 44 & 44 & 5.29 & 33 & 11.50 \\ \hline
    gen200\_p0.9\_55   & 55 & 55 & 6.18 & 37 & 10.82 \\ \hline
    gen400\_p0.9\_55   & 55 & 51 & 9.95 & 43 & 19.63 \\ \hline
    gen400\_p0.9\_65   & 65 & 65 & 21.35 & 44 & 24.56 \\ \hline
    gen400\_p0.9\_75   & 75 & 75 & 13.42 & 44 & 23.98 \\ \hline
    hamming10-4        & 40 & 40 & 52.06 & 30 & 61.17 \\ \hline
    hamming8-4         & 16 & 16 & 5.23 & 12 & 13.84 \\ \hline
    keller4            & 11 & 11 & 0.78 & 9 & 8.85 \\ \hline
    keller5            & 27 & 25 & 6.38 & 21 & 48.53 \\ \hline
    keller6            & 59 & 49 & 51.53 & 41 & 244.11 \\ \hline
    MANN\_a27          & 126 & 121 & 12.68 & 122 & 17.65 \\ \hline
    p\_hat300-1        & 8 & 8 & 1.22 & 7 & 18.32 \\ \hline
    p\_hat300-2        & 25 & 25 & 4.16 & 21 & 15.20 \\ \hline
    p\_hat300-3        & 36 & 36 & 7.03 & 32 & 16.92 \\ \hline
    p\_hat700-1        & 11 & 10 & 3.65 & 8 & 39.18 \\ \hline
    p\_hat700-2        & 44 & 44 & 22.02 & 37 & 41.78 \\ \hline
    p\_hat700-3        & 62* & 62 & 28.50 & 52 & 41.75 \\ \hline
    p\_hat1500-1       & 12 & 11 & 9.77 & 8 & 89.71 \\ \hline
    p\_hat1500-2       & 65* & 65 & 62.62 & 51 & 94.75 \\ \hline
    p\_hat1500-3       & 94* & 94 & 77.15 & 75 & 91.43 \\ \hline
  \caption{Resultados en enfriamiento simulado.}
\end{longtable}
\end{small}

El algoritmo de enfriamiento simulado básico ha obtenido unos resultados de $16$
óptimos y un $93.5\%$ de media en el tamaño de los cliques, con respecto al máximo
conocido, mientras que la adaptación del algoritmo de Geng et al. ha obtenido
$0$ óptimos y un $75.7\%$ de tamaño medio.

Los resultados en el algoritmo básico mejoran los obtenidos por ambos algoritmos
de búsqueda local, conviertiéndolo en una buena herramienta para resolver el problema,
pues supera ampliamente el $90\%$ de tamaño medio, y alcanza el óptimo casi en la
mitad de instancias. Recordemos que toma un entorno distinto al considerado en la
búsqueda local, sin priorizar ningún movimiento sobre otro. Esto, y el uso de
técnicas de enfriamiento simulado podrían ser los causantes de que el algoritmo
mejore las búsquedas locales. Además, presenta unos tiempos de ejecución bajos,
sin llegar a los $90$ segundos en ningún caso. Esto se debe a que el número de
iteraciones no depende del grafo, sino que viene dado por la temperatura y es
fijo en todos ellos.

Por otra parte, la adaptación del algoritmo de Geng et al. presenta unos resultados
muy pobres, sin ser capaz de encontrar en óptimo en ningún caso. Podemos afirmar que
falla a la hora de resolver el problema de forma satisfactoria. Dado que la exploración
del entorno es aleatoria, como en el caso anterior, cabe plantearse si considerar grafos
cualesquiera en lugar de cliques como espacio de soluciones es la causa de este mal
rendimiento. La función objetivo, con componentes aparte del tamaño de los grafos,
también podría ser causante de que el algoritmo no se dirija a las zonas del espacio
de soluciones de mayor calidad. Además, la función de reducción de grafos a cliques
es otro factor a tener en cuenta, pues no tenemos garantizado que la reducción sea
óptima, y añade otra etapa de posible pérdida de calidad en las soluciones.

Con estos resultados qued comprobado que el algoritmo básico funciona mucho
mejor que la adaptación del algoritmo de Geng et al.

\subsection{ILS}

\begin{small}
\begin{longtable}{l l l l l l}
  \label{table:ils}\\
    & & \multicolumn{2}{c}{ILS - 1LS} & \multicolumn{2}{c}{ILS - DLS} \\ \cline{3-6}
    Grafo              & Clique & |C| & Tiempo (s) & |C| & Tiempo (s) \\ \hline
    \endhead
    \endfoot
    brock200\_2        & 12 & 10 & 4.36 & 10 & 0.12\\ \hline
    brock200\_4        & 17 & 15 & 6.69 & 16 & 0.21\\ \hline
    brock400\_2        & 29 & 24 & 34.90 & 24 & 0.89\\ \hline
    brock400\_4        & 33 & 24 & 35.40 & 24 & 0.90\\ \hline
    brock800\_2        & 24 & 18 & 116.12 & 19 & 1.72\\ \hline
    brock800\_4        & 26 & 18 & 116.78 & 20 & 1.53\\ \hline
    C125.9             & 34* & 34 & 6.91 & 34 & 0.55\\ \hline
    C250.9             & 44* & 44 & 25.94 & 44 & 1.36\\ \hline
    C500.9             & 57 & 57 & 124.56 & 55 & 3.72\\ \hline
    C1000.9            & 68 & 61 & 514.84 & 63 & 10.28\\ \hline
    C2000.5            & 16* & 12 & 547.22 & 14 & 2.59\\ \hline
    C2000.9            & 80 & 72 & 2512.71 & 72 & 29.10\\ \hline
    C4000.5            & 18* & 16 & 2489.99 & 15 & 6.50\\ \hline
    DSJC500\_5         & 13 & 12 & 31.75 & 12 & 0.43\\ \hline
    DSJC1000\_5        & 15 & 14 & 145.22 & 14 & 1.14\\ \hline
    gen200\_p0.9\_44   & 44 & 39 & 14.99 & 40 & 0.98\\ \hline
    gen200\_p0.9\_55   & 55 & 55 & 17.33 & 55 & 0.99\\ \hline
    gen400\_p0.9\_55   & 55 & 51 & 82.47 & 51 & 2.43 \\ \hline
    gen400\_p0.9\_65   & 65 & 49 & 71.32 & 65 & 2.48 \\ \hline
    gen400\_p0.9\_75   & 75 & 75 & 134.14 & 75 & 2.98 \\ \hline
    hamming10-4        & 40 & 38 & 315.82 & 40 & 3.94\\ \hline
    hamming8-4         & 16 & 16 & 8.27 & 16 & 0.13\\ \hline
    keller4            & 11 & 10 & 4.40 & 11 & 0.17 \\ \hline
    keller5            & 27 & 23 & 180.97 & 25 & 2.91 \\ \hline
    keller6            & 59 & 48 & 10034.52 & 54 & 58.89\\ \hline
    MANN\_a27          & 126 & 126 & 1106.07 & 125 & 55.67 \\ \hline
    p\_hat300-1        & 8 & 7 & 8.36 & 8 & 0.12 \\ \hline
    p\_hat300-2        & 25 & 25 & 21.33 & 25 & 0.83\\ \hline
    p\_hat300-3        & 36 & 36 & 26.64 & 36 & 1.24\\ \hline
    p\_hat700-1        & 11 & 9 & 53.69 & 10 & 0.47\\ \hline
    p\_hat700-2        & 44 & 44 & 179.95 & 44 & 4.36\\ \hline
    p\_hat700-3        & 62* & 61 & 282.32 & 62 & 6.67 \\ \hline
    p\_hat1500-1       & 12 & 10 & 255.15 & 10 & 1.24\\ \hline
    p\_hat1500-2       & 65* & 65 & 1417.50 & 65 & 22.65\\ \hline
    p\_hat1500-3       & 94* & 93 & 2167.37 & 94 & 35.91\\ \hline
  \caption{Resutados en ILS.}
\end{longtable}
\end{small}

El algoritmo ILS que utiliza 1LS, ha obtenido unos resultados de $11$ óptimos y
un $90.3\%$ de tamaño medio, mientras que el que utiliza DLS ha obtenido $15$
óptimos y un $93.3\%$ de tamaño medio.

Los resultados en ambos algoritmos mejoran sensiblemente los alcanzados en las
respectivas búsquedas locales, lo que nos dice que el uso de técnicas de reinicialización
es efectivo a la hora de salir de óptimos globales, obteniendo soluciones de partida
que nos permiten explorar de forma más efectiva el espacio de soluciones y alcanzar
mejores cliques.

ILS + 1LS es capaz de obtener buenas soluciones, estando ligeramente por encima del
$90\%$ de tamaño medio en cliques. Sus tiempos de ejecución empiezan a ser notablemente
altos, con valores que llegan hasta casi $3$ horas en \textit{keller6}, debido a
la repetición de 1LS en $20$ ocasiones. Para tratar de que este algoritmo pueda
competir con el resto, habría que experimentar si se pueden mantener los resultados
reduciendo el número de reinicializaciones, para así disminuir los tiempos de ejecución.

ILS + DLS es mejor que la versión con 1LS tanto en resultados como en tiempos de
ejecución, resultado que era de esperar, vistos los obtenidos en la búsqueda local.
En este caso la ejecución no toma más de $60$ segundos en ninguno de los grafos,
mejorando sensiblemente los tiempos del algoritmo anterior.

Concluímos que ambos algoritmos son capaces de obtener soluciones de calidad,
siendo mejor la versión con DLS por sus mejores resultados y menor tiempo de
ejecución, resultado que es consistente con los obtenidos para DLS e ILS.

\subsection{GRASP}

\begin{small}
\begin{longtable}{l l l l l l}
  \label{table:grasp}\\
    & & \multicolumn{2}{c}{GRASP - 1LS} & \multicolumn{2}{c}{GRASP - DLS} \\ \cline{3-6}
    Grafo              & Clique & |C| & Tiempo (s) & |C| & Tiempo (s) \\ \hline
    \endhead
    \endfoot
    brock200\_2        & 12 & 11 & 4.38 & 11 & 0.28\\ \hline
    brock200\_4        & 17 & 16 & 6.79 & 16 & 0.47\\ \hline
    brock400\_2        & 29 & 24 & 33.76 & 23 & 2.21\\ \hline
    brock400\_4        & 33 & 25 & 32.25 & 33 & 2.23\\ \hline
    brock800\_2        & 24 & 19 & 116.77 & 19 & 4.78\\ \hline
    brock800\_4        & 26 & 19 & 101.59 & 20 & 4.60\\ \hline
    C125.9             & 34* & 34 & 5.90 & 34 & 0.99\\ \hline
    C250.9             & 44* & 44 & 25.92 & 43 & 3.38\\ \hline
    C500.9             & 57 & 53 & 114.25 & 53 & 12.17\\ \hline
    C1000.9            & 68 & 63 & 575.84 & 63 & 36.31\\ \hline
    C2000.5            & 16* & 15 & 598.26 & 15 & 16.16\\ \hline
    C2000.9            & 80 & 71 & 2811.73 & 69 & 134.39\\ \hline
    C4000.5            & 18* & 16 & 2561.37 & 16 & 58.44\\ \hline
    DSJC500\_5         & 13 & 13 & 31.00 & 13 & 1.26\\ \hline
    DSJC1000\_5        & 15 & 14 & 150.85 & 14 & 4.44\\ \hline
    gen200\_p0.9\_44   & 44 & 44 & 17.04 & 44 & 2.24\\ \hline
    gen200\_p0.9\_55   & 55 & 55 & 16.21 & 55 & 2.57\\ \hline
    gen400\_p0.9\_55   & 55 & 51 & 78.42 & 49 & 8.42\\ \hline
    gen400\_p0.9\_65   & 65 & 52 & 76.15 & 65 & 8.72\\ \hline
    gen400\_p0.9\_75   & 75 & 75 & 94.38 & 75 & 9.98\\ \hline
    hamming10-4        & 40 & 40 & 340.81 & 40 & 19.86\\ \hline
    hamming8-4         & 16 & 16 & 8.46 & 16 & 0.47\\ \hline
    keller4            & 11 & 11 & 4.47 & 11 & 0.33\\ \hline
    keller5            & 27 & 26 & 200.16 & 27 & 7.50\\ \hline
    keller6            & 59 & 54 & 8958.52 & 53 & 181.52\\ \hline
    MANN\_a27          & 126 & 125 & 1453.48 & 125 & 79.18\\ \hline
    p\_hat300-1        & 8 & 8 & 9.00 & 8 & 0.32\\ \hline
    p\_hat300-2        & 25 & 25 & 21.48 & 25 & 1.46\\ \hline
    p\_hat300-3        & 36 & 36 & 31.68 & 36 & 2.70\\ \hline
    p\_hat700-1        & 11 & 9 & 54.35 & 11 & 1.51\\ \hline
    p\_hat700-2        & 44 & 44 & 215.50 & 44 & 9.21\\ \hline
    p\_hat700-3        & 62* & 62 & 303.31 & 62 & 18.55\\ \hline
    p\_hat1500-1       & 12 & 10 & 266.69 & 11 & 5.94 \\ \hline
    p\_hat1500-2       & 65* & 65 & 1600.58 & 65 & 48.06\\ \hline
    p\_hat1500-3       & 94* & 94 & 2297.67 & 93 & 103.41\\ \hline
  \caption{Resultados en GRASP.}
\end{longtable}
\end{small}

Los resultados obtenidos por ambos algoritmos GRASP han sido de $16$ óptimos y un
$93.5\%$ de tamaño medio para la versión con 1LS, y $18$ óptimos y un $95.3\%$ para
la versión con DLS.

Estos algoritmos son capaces de alcanzar el óptimo en aproximadamente un $50\%$ de
los casos de estudio, reflejando un buen comportamiento general. Además, la versión
que usa DLS como búsqueda local supera el $95\%$ de tamaño medio en los grafos, por
lo que estamos ante un algoritmo capaz de dar soluciones cercanas al óptimo.
El algoritmo con 1LS también ofrece resultados de calidad, incluso mejorando los
resultados obtenidos por el otro GRASP en grafos como \textit{C2000.9} o \textit{p\_hat1500-3}.

Nuevamente, el uso de DLS como búsqueda local proporciona mejores resultados que
si utilizamos la primera búsqueda local, lo que nos reafirma en su mejor comportamiento,
tanto por separado como si se utiliza en técnicas multiarranque. Los tiempos de
ejecución también favorecen a la versión con DLS, sobre todo en grafos con un mayor
número de nodos, donde la mayor complejidad computacional de 1LS afecta notablemente
a dichos tiempos. Ejemplos de esto son grafos como \textit{keller6, C4000.5} y \textit{C2000.9},
en los que la versión con DLS tiene tiempos del orden de los minutos, y la versión
con 1LS tiene tiempos del orden de horas.

Los dos algoritmos proporcionan mejores resultados que sus correspondientes versiones
de búsqueda local iterada, lo que en un principio nos dice que el uso de soluciones
pseudoaleatorias como solución de partida proporciona mejores resultados que el proceso
de mutación utilizado en ILS. No obstante, dado que estos procesos pueden ser modificados,
queda aun margen de mejora en ambos algoritmos, por lo que sería pronto para decir que
GRASP funciona mejor que ILS en el problema del clique máximo. Sería necesaria más
experimentación para poder sacar conclusiones a la hora de comparar ambas técnicas.

\subsection{Algoritmos de colonia de hormigas}

\begin{small}
\begin{longtable}{l l l l l l}
  \label{table:aco}\\
    & & \multicolumn{2}{c}{ACO básico} & \multicolumn{2}{c}{ACO + SA} \\ \cline{3-6}
    Grafo              & Clique & |C| & Tiempo (s) & |C| & Tiempo (s) \\ \hline
    \endhead
    \endfoot
    brock200\_2        & 12 & 12 & 2.24 & 12 & 2.56\\ \hline
    brock200\_4        & 17 & 16 & 4.60 & 16 & 4.96\\ \hline
    brock400\_2        & 29 & 22 & 17.35 & 22 & 18.92\\ \hline
    brock400\_4        & 33 & 23 & 17.45 & 22 & 18.99\\ \hline
    brock800\_2        & 24 & 19 & 25.39 & 19 & 27.56\\ \hline
    brock800\_4        & 26 & 18 & 25.19 & 19 & 27.41\\ \hline
    C125.9             & 34* & 34 & 17.12 & 34 & 14.18\\ \hline
    C250.9             & 44* & 44 & 45.33 & 42 & 38.78\\ \hline
    C500.9             & 57 & 53 & 111.15 & 50 & 104.84\\ \hline
    C1000.9            & 68 & 55 & 264.30 & 54 & 268.65\\ \hline
    C2000.5            & 16* & 14 & 42.72 & 14 & 46.29\\ \hline
    C2000.9            & 80 & 60 & 724.77 & 60 & 748.91\\ \hline
    C4000.5            & 18* & 15 & 98.38 & 16 & 116.37\\ \hline
    DSJC500\_5         & 13 & 13 & 7.32 & 12 & 8.20\\ \hline
    DSJC1000\_5        & 15 & 13 & 17.80 & 13 & 19.62\\ \hline
    gen200\_p0.9\_44   & 44 & 41 & 32.81 & 39 & 27.14\\ \hline
    gen200\_p0.9\_55   & 55 & 55 & 53.89 & 55 & 34.24\\ \hline
    gen400\_p0.9\_55   & 55 & 49 & 89.14 & 46 & 77.33\\ \hline
    gen400\_p0.9\_65   & 65 & 57 & 93.20 & 48 & 79.71\\ \hline
    gen400\_p0.9\_75   & 75 & 75 & 112.08 & 65 & 85.24 \\ \hline
    hamming10-4        & 40 & 33 & 68.76 & 32 & 72.97\\ \hline
    hamming8-4         & 16 & 16 & 3.52 & 16 & 4.03\\ \hline
    keller4            & 11 & 11 & 2.27 & 11 & 2.54\\ \hline
    keller5            & 27 & 24 & 34.34 & 24 & 37.73\\ \hline
    keller6            & 59 & 48 & 618.39 & 45 & 630.96\\ \hline
    MANN\_a27          & 126 & 125 & 665.68 & 124 & 635.35\\ \hline
    p\_hat300-1        & 8 & 8 & 1.88 & 8 & 2.19\\ \hline
    p\_hat300-2        & 25 & 25 & 14.57 & 25 & 13.66\\ \hline
    p\_hat300-3        & 36 & 36 & 32.46 & 35 & 28.17\\ \hline
    p\_hat700-1        & 11 & 9 & 5.52 & 10 & 6.35\\ \hline
    p\_hat700-2        & 44 & 44 & 70.94 & 43 & 64.89\\ \hline
    p\_hat700-3        & 62* & 61 & 182.29 & 61 & 150.63\\ \hline
    p\_hat1500-1       & 12 & 10 & 15.65 & 10 & 28.12 \\ \hline
    p\_hat1500-2       & 65* & 44 & 195.39 & 54 & 203.79\\ \hline
    p\_hat1500-3       & 94* & 87 & 567.45 & 82 & 524.60 \\ \hline
  \caption{Resultados en algoritmos de colonia de hormigas.}
\end{longtable}
\end{small}

Los algoritmos de colonia de hormigas han obtenido unos resultados de $12$ óptimos
y un $89.7\%$ de tamaño medio la versión básica, y $7$ óptimos, con un $88.5\%$
de tamaño medio la versión con enfriamiento simulado.

Ambos algoritmos ofrecen resultados similares en cuanto al tamaño medio de los cliques
encontrados, si bien la versión que añade información del grafo con enfrimiento simulado
es ligeramente peor, encontrando menos óptimos y con un porcentaje menor de tamaño medio.
Estos resultados indican que son necesarios cambios y más experimentación para ver si
añadir dicho  término a la probabilidad es capaz de aportar mejores nodos a la
construcción de soluciones.

Los resultados del algoritmo básico son buenos, rozando el $90\%$ de tamaño medio,
aunque queda por detrás de otras técnicas, como las búsquedas multiarranque o el
enfriamiento simulado básico. Para tratar de mejorar su rendimiento, podríamos
considerar añadir técnicas para construir un algoritmo híbrido, como una busqueda
basada en entornos, aprovechando mejor el comportamiento de las hormigas e intentando
que estas construyan mejores soluciones.

En el algoritmo con enfriamiento simulado, si comparamos con el anterior, se
alcanzan soluciones de menor o igual tamaño en $31$ de los $35$ casos, por lo que
el término añadido al cálculo de la probabilidad no parece haber tenido repercusión
positiva en la obtención de soluciones. Habría que plantearse otra forma de cambiar
la función de probabilidad para que se use la información del grafo de forma positiva.

Ambos algoritmos son similares en cuanto a tiempos, ya que la única diferencia es
el cálculo añadido en la probabilidad en la versión con enfriamiento simulado,
que no parece repercutir de forma significativa. En general, los tiempos no son
muy extensos en todo el conjunto de grafos, aunque alcanzan valores un tanto altos
en las instancias con más nodos, siendo \textit{C2000.9} la que más tiempo requiere,
unos $12$ minutos para cada uno de los dos algoritmos.

\subsection{Algoritmos genéticos}

\begin{small}
\begin{longtable}{l l l l l l}
  \label{table:geneticos}\\
    && \multicolumn{2}{c}{Genético} & \multicolumn{2}{c}{Memético} \\ \cline{3-6}
    Grafo              & Clique & |C| & Tiempo (s) & |C| & Tiempo (s) \\ \hline
    \endhead
    \endfoot
    brock200\_2        & 12 & 12 & 2.10 & 12 & 6.74\\ \hline
    brock200\_4        & 17 & 17 & 4.02 & 17 & 13.11\\ \hline
    brock400\_2        & 29 & 23 & 16.45 & 25 & 50.66\\ \hline
    brock400\_4        & 33 & 25 & 16.44 & 33 & 51.26\\ \hline
    brock800\_2        & 24 & 20 & 24.77 & 20 & 72.04\\ \hline
    brock800\_4        & 26 & 19 & 24.16 & 20 & 77.78\\ \hline
    C125.9             & 34* & 34 & 12.48 & 34 & 42.86\\ \hline
    C250.9             & 44* & 42 & 33.91 & 44 & 143.46\\ \hline
    C500.9             & 57 & 48 & 94.39 & 56 & 284.240\\ \hline
    C1000.9            & 68 & 54 & 261.16 & 65 & 765.52\\ \hline
    C2000.5            & 16* & 15 & 44.22 & 16 & 123.12\\ \hline
    C2000.9            & 80 & 60 & 737.43 & 73 & 2128.81\\ \hline
    C4000.5            & 18* & 12 & 115.13 & 17 & 297.26\\ \hline
    DSJC500\_5         & 13 & 13 & 7.14 & 13 & 20.25\\ \hline
    DSJC1000\_5        & 15 & 13 & 17.62 & 15 & 51.49\\ \hline
    gen200\_p0.9\_44   & 44 & 40 & 24.00 & 44 & 78.15\\ \hline
    gen200\_p0.9\_55   & 55 & 55 & 24.99 & 55 & 54.66\\ \hline
    gen400\_p0.9\_55   & 55 & 46 & 68.00 & 55 & 233.02 \\ \hline
    gen400\_p0.9\_65   & 65 & 47 & 69.90 & 65 & 216.87\\ \hline
    gen400\_p0.9\_75   & 75 & 69 & 77.20 & 75 & 199.47\\ \hline
    hamming10-4        & 40 & 34 & 75.71 & 40 & 205.56\\ \hline
    hamming8-4         & 16 & 16 & 3.58 & 16 & 6.14\\ \hline
    keller4            & 11 & 11 & 1.95 & 11 & 10.73\\ \hline
    keller5            & 27 & 27 & 33.26 & 27 & 192.61\\ \hline
    keller6            & 59 & 47 & 634.94 & 53 & 4265.78\\ \hline
    MANN\_a27          & 126 & 125 & 646.76 & 125 & 1682.67\\ \hline
    p\_hat300-1        & 8 & 8 & 1.92 & 8 & 6.20\\ \hline
    p\_hat300-2        & 25 & 25 & 11.45 & 25 & 50.88\\ \hline
    p\_hat300-3        & 36 & 36 & 25.37 & 36 & 85.84\\ \hline
    p\_hat700-1        & 11 & 11 & 6.05 & 11 & 18.93\\ \hline
    p\_hat700-2        & 44 & 44 & 81.02 & 44 & 270.53\\ \hline
    p\_hat700-3        & 62* & 60 & 149.41 & 62 & 565.92\\ \hline
    p\_hat1500-1       & 12 & 11 & 18.45 & 11 & 56.23\\ \hline
    p\_hat1500-2       & 65* & 61 & 319.06 & 65 & 1472.25\\ \hline
    p\_hat1500-3       & 94* & 86 & 615.29 & 94 & 2964.62\\ \hline
  \caption{Resultados en algoritmos genéticos.}
\end{longtable}
\end{small}

El algoritmo genético ha dado unos resultados de $13$ óptimos encontrados y un
$90.5\%$ de tamaño medio de los cliques, mientras que su versión híbrida con la
búsqueda local DLS ha encontrado $25$ óptimos, con un tamaño medio de un $97.3\%$,
lo que lo convierte en el mejor de los algoritmos implementados en cuanto a resultados.
Queda claro que usar una búsqueda local para mejorar la población del algoritmo
genético ayuda a obtener mejores soluciones.

Por sí solo, el algoritmo genético presenta un buen comportamiento, siendo capaz de
encontrar soluciones aceptables en el conjunto de instancias considerado. Habría
que experimentar con distintos parámetros tanto en probabilidades de cruce y mutación
como en número de iteraciones y tamaño de población para poder realizar un ajuste más
minucioso, que conllevase una mejora de los resultados. También
debería estudiarse la necesidad de reinicializar algunos individuos de la población,
en caso de que estos fueran similares unos a otros, para introducir diversidad.

El algoritmo memético fruto de la combinación del genético y DLS logra mejorar
considerablemente el funcionamiento del genético, proporcionando muy buenos resultados
en el conjunto de instancias. El $97.3\%$ de tamaño medio nos indica que estamos
ante un algoritmo capaz de encontrar soluciones cercanas al óptimo en todas las
instancias, pues si nos fijamos en los grafos por separado, vemos que no hay ninguno
en el que el rendimiento sea especialmente malo.

Dado que la única diferencia que existe entre las dos versiones es el uso de DLS,
se pone de  manifiesto que es capaz de mejorar buenos algoritmos, lo que nos lleva
a considerar su uso para otros algoritmos, como podría ser un ACO, el cual daba
unos resultados similares al genético.

\subsection{Comportamiento general}

El conjunto de algoritmos implementados muestra, en general, un comportamiento dentro de
lo esperado. Hemos visto como los algoritmos \textit{greedy} no son capaces de encontrar
por sí solos buenas soluciones al problema, aunque aportan cliques que pueden ser utilizados
por otros métodos como soluciones iniciales. La exploración del entorno hecha con búsquedas
locales o enfriamiento simulado mejora estos resultados, si bien no todos los algoritmos de
este tipo proporcionan buenas soluciones, como demuestra la adaptación de la idea de Geng
et al. Estas soluciones pueden ser aun mejores si se utilizan técnicas multiarranque,
hecho que respaldan los resultados de los algoritmos ILS y GRASP.

Dado su amplio uso en problemas de grafos, y sus buenos resultados en general, cabía
esperar que los algoritmos de colonia de hormigas funcionasen bien en este problema, y
así ha sido, aunque sus resultados no han destacado especialmente. Por otra parte,
algoritmos menos usuales en problemas de grafos como los genéticos, han mostrado un
comportamiento bueno, en especial el algoritmo memético, claramente el mejor en
cuanto a resultados de todos los considerados. Ajustar en ambos tipos de algoritmos
los parámetros involucrados podría mejorar su comportamiento, sobre todo en el caso
de los ACO, donde los resultados han sido algo peores.

Para realizar una comparativa entre todos los algoritmos, vamos a clasificarlos en
todos los grafos que se han considerado. Tomaremos los tres algoritmos que den mejores
resultados en cuanto a tamaño del clique encontrado, y usaremos el tiempo para decidir
en caso de empate. De esta forma, otorgaremos primer, segundo y trecer puesto para
cada grafo, y, haciendo un recuento para todos los algoritmos, veremos como funcionan
con respecto a los demas.

A continuación se muestra

\begin{small}
\begin{longtable}{l l l l}
  \label{table:posiciones}\\
  \endhead
  \endfoot
    Grafo              & $1^{er}$ puesto & $2$º puesto & $3^{er}$ puesto \\ \hline
    brock200\_2        & Genético & ACO básico & ACO + SA  \\  \hline
    brock200\_4        & Genético & Memético & ILS + DLS  \\  \hline
    brock400\_2        & Memético & ILS + DLS & SA básico  \\  \hline
    brock400\_4        & GRASP + DLS & Memético & Genético  \\  \hline
    brock800\_2        & Genético & Memético & DLS  \\  \hline
    brock800\_4        & ILS + DLS & GRASP + DLS & 1LS  \\  \hline
    C125.9             & DLS & ILS + DLS & GRASP + DLS  \\  \hline
    C250.9             & ILS + DLS & SA básico & GRASP + 1LS  \\  \hline
    C500.9             & ILS + 1LS & Memético & ILS + DLS \\  \hline
    C1000.9            & Memético & ILS + DLS & 1LS  \\  \hline
    C2000.5            & Memético & GRASP + DLS & SA básico  \\  \hline
    C2000.9            & Memético & ILS + DLS & 1LS  \\  \hline
    C4000.5            & Memético & GRASP + DLS & ACO + SA  \\  \hline
    DSJC500\_5         & GRASP + DLS & Genético & ACO básico  \\  \hline
    DSJC1000\_5        & Memético & ILS + DLS & GRASP + DLS  \\  \hline
    gen200\_p0.9\_44   & GRASP + DLS & SA básico & GRASP + 1LS  \\  \hline
    gen200\_p0.9\_55   & DLS & ILS + DLS & GRASP + DLS  \\  \hline
    gen400\_p0.9\_55   & Memético & ILS + DLS & SA básico  \\  \hline
    gen400\_p0.9\_65   & ILS + DLS & GRASP + DLS & SA básico  \\  \hline
    gen400\_p0.9\_75   & DLS & ILS + DLS & GRASP + DLS  \\  \hline
    hamming10-4        & DLS & ILS + DLS & GRASP + DLS  \\  \hline
    hamming8-4         & DLS & Greedy básico & ILS + DLS  \\  \hline
    keller4            & DLS & ILS + DLS & 1LS  \\  \hline
    keller5            & DLS & GRASP + DLS & Genético  \\  \hline
    keller6            & ILS + DLS & GRASP + 1LS & GRASP + DLS  \\  \hline
    MANN\_a27          & ILS + 1LS & Greedy básico & DLS  \\  \hline
    p\_hat300-1        & DLS & Greedy adaptativo & ILS + DLS  \\  \hline
    p\_hat300-2        & DLS & ILS + DLS & 1LS  \\  \hline
    p\_hat300-3        & DLS & ILS + DLS & GRASP + DLS  \\  \hline
    p\_hat700-1        & GRASP + DLS & Genético & Memético  \\  \hline
    p\_hat700-2        & ILS + DLS & GRASP + DLS & 1LS  \\  \hline
    p\_hat700-3        & ILS + DLS & GRASP + DLS & SA básico  \\  \hline
    p\_hat1500-1       & DLS & Greedy adaptativo & GRASP + DLS  \\  \hline
    p\_hat1500-2       & ILS + DLS & GRASP + DLS & SA básico  \\  \hline
    p\_hat1500-3       & ILS + DLS & SA básico & GRASP + 1LS  \\  \hline
  \caption{Clasificación por grafos.}
\end{longtable}
\end{small}

Si hacemos un recuento de las posiciones obtenidas por los grafos, obtenemos los
siguientes resultados:

\begin{small}
\begin{longtable}{l l l l}
  \label{table:recuento}\\
  \endhead
  \endfoot
    Algoritmo              & $1^{er}$ puesto & $2$º puesto & $3^{er}$ puesto \\ \hline
    Greedy básico          &  0  &  2  &  0  \\ \hline
    Greedy adaptativo      &  0  &  2  &  0  \\ \hline
    1LS                    &  0  &  0  &  6  \\ \hline
    DLS                    &  11  &  0  &  2  \\ \hline
    SA básico              &  0  &  3  &  6  \\ \hline
    SA adaptado            &  0  &  0  &  0  \\ \hline
    ILS + 1LS              &  2  &  0  &  0  \\ \hline
    ILS + DLS              &  8  &  12  &  4  \\ \hline
    GRASP + 1LS            &  0  &  1  &  3  \\ \hline
    GRASP + DLS            &  4  &  8  &  8  \\ \hline
    ACO básico             &  0  &  1  &  1  \\ \hline
    ACO + SA               &  0  &  0  &  2  \\ \hline
    Genético               &  3  &  2  &  2  \\ \hline
    Memético               &  7  &  4  &  1  \\ \hline
  \caption{Recuento.}
\end{longtable}
\end{small}


% Ránking. Medallero.

\section{Conclusiones}

Hemos estudiado de forma satisfactoria el funcionamiento de todos los algoritmos implementados,
que pueden verse como un primer acercamiento al problema que nos permite evaluar las
prestaciones de distintas formas de abordarlo. El trabajo ha sido provechoso en cuanto
a resultados obtenidos, y puede utilizarse como base para futuros trabajos en esta materia.

La línea de trabajo a seguir a partir de aquí podría ser una de las siguientes:
\begin{itemize}
  \item Enfatización en un tipo de algoritmos, experimentando con diversos cambios que
        mejoren su funcionamiento. Dependiendo del tipo de algoritmo, podrían considerarse
        estrategias de reinicialización, énfasis en ciertas fases de búsqueda de soluciones,
        o búsqueda de nuevas medidas que clasifiquen mejor a las soluciones obtenidas, entre
        muchas otras. De esta forma, se intentaría obtener un algoritmo más específico para
        resolver el problema, con un mejor funcionamiento.

  \item Ampliación del trabajo realizado. Existen numerosas técnicas que no se han implementado
        en este trabajo, como la búsqueda tabú, la búsqueda en entornos variables, la búsqueda
        dispersa, etc. Analizar su comportamiento nos daría más información sobre el problema,
        aportando nuevas ideas para tratar de mejorar estos u otros algoritmos.
\end{itemize}




% ------------------------------------------------
% Tablas sin tiempos
\iffalse

\begin{table}[H]
\centering
\caption{Resultados}
\begin{tabular}{|l|l|l|l|l|l|l|l|}
\hline
Grafo              & Mejor conocido & Greedy1 & Greedy2 & BL1 & BL2 & SA1 & SA2 \\ \hline
brock200\_2        & 12 & 7 & 9 & 9 & 11 & 10 & 9 \\ \hline
brock200\_4        & 17 & 13 & 15 & 15 & 15 & 16 & 13 \\ \hline
brock400\_2        & 29 & 20 & 22 & 22 & 23 & 24 & 18 \\ \hline
brock400\_4        & 33 & 18 & 21 & 23 & 23 & 24 & 19 \\ \hline
brock800\_2        & 24 & 14 & 19 & 19 & 19 & 19 & 16 \\ \hline
brock800\_4        & 26 & 14 & 19 & 20 & 18 & 20 & 16 \\ \hline
C125.9             & 34* & 29 & 32 & 33 & 34 & 34 & 32 \\ \hline
C250.9             & 44* & 35 & 40 & 42 & 41 & 44 & 36 \\ \hline
C500.9             & 57 & 43 & 51 & 54 & 54 & 53 & 45 \\ \hline
C1000.9            & 68 & 51 & 56 & 63 & 62 & 63 & 52 \\ \hline
C2000.5            & 16* & 10 & 13 & 13 & 14 & 15 & 12 \\ \hline
C2000.9            & 80 & 56 & 66 & 72 & 68 & 71 & 57 \\ \hline
C4000.5            & 18* & 12 & 15 & 15 & 15 & 15 & 13 \\ \hline
DSJC500\_5         & 13 & 8 & 11 & 11 & 12 & 12 & 10 \\ \hline
DSJC1000\_5        & 15 & 10 & 12 & 13 & 13 & 14 & 11 \\ \hline
gen200\_p0.9\_44   & 44 & 32 & 36 & 39 & 43 & 44 & 33 \\ \hline
gen200\_p0.9\_55   & 55 & 36 & 39 & 39 & 55 & 55 & 37 \\ \hline
gen400\_p0.9\_55   & 55 & 44 & 43 & 50 & 49 & 51 & 43 \\ \hline
gen400\_p0.9\_65   & 65 & 40 & 44 & 45 & 49 & 65 & 44 \\ \hline
gen400\_p0.9\_75   & 75 & 45 & 46 & 49 & 75 & 75 & 44 \\ \hline
hamming10-4        & 40 & 32 & 36 & 36 & 40 & 40 & 30 \\ \hline
hamming8-4         & 16 & 16 & 16 & 16 & 16 & 16 & 12 \\ \hline
keller4            & 11 & 8 & 11 & 11 & 11 & 11 & 9 \\ \hline
keller5            & 27 & 17 & 24 & 23 & 27 & 25 & 21 \\ \hline
keller6            & 59 & 37 & 48 & 51 & 50 & 49 & 41 \\ \hline
MANN\_a27          & 126 & 125 & 125 & 125 & 123 & 121 & 122 \\ \hline
p\_hat300-1        & 8 & 7 & 8 & 7 & 7 & 8 & 7 \\ \hline
p\_hat300-2        & 25 & 23 & 23 & 25 & 25 & 25 & 21 \\ \hline
p\_hat300-3        & 36 & 30 & 34 & 34 & 36 & 36 & 32 \\ \hline
p\_hat700-1        & 11 & 7 & 8 & 9 & 9 & 10 & 8 \\ \hline
p\_hat700-2        & 44 & 38 & 43 & 44 & 43 & 44 & 37 \\ \hline
p\_hat700-3        & 62* & 55 & 58 & 60 & 61 & 62 & 52 \\ \hline
p\_hat1500-1       & 12 & 8 & 11 & 11 & 11 & 11 & 8 \\ \hline
p\_hat1500-2       & 65* & 54 & 62 & 64 & 63 & 65 & 51 \\ \hline
p\_hat1500-3       & 94* & 75 & 87 & 91 & 91 & 94 & 75 \\ \hline
Óptimos            & 35  & 1  & 3  & 4  & 9  & 16 & 0  \\ \hline
Aciertos           & 100\% & 75.5\% & 85.0\% & 88.5\% & 92.2\% & 94.9\% & 76.7\% \\ \hline
Media aciertos     & 100\% & 73.2\% & 84.9\% & 87.5\% & 91.0\% & 93.5\% & 75.7\% \\ \hline
\end{tabular}
\end{table}

\begin{table}[H]
\centering
\caption{Resultados}
\begin{tabular}{|l|l|l|l|l|l|}
\hline
Grafo              & Mejor conocido & ILS1 & ILS2 & GRASP1 & GRASP2 \\ \hline
brock200\_2        & 12 & 10 & 10 & 11 & 11 \\ \hline
brock200\_4        & 17 & 15 & 16 & 16 & 16 \\ \hline
brock400\_2        & 29 & 24 & 24 & 24 & 23 \\ \hline
brock400\_4        & 33 & 24 & 24 & 25 & 33 \\ \hline
brock800\_2        & 24 & 18 & 19 & 19 & 19 \\ \hline
brock800\_4        & 26 & 18 & 20 & 19 & 20 \\ \hline
C125.9             & 34* & 34 & 34 & 34 & 34 \\ \hline
C250.9             & 44* & 44 & 44 & 44 & 43 \\ \hline
C500.9             & 57 & 57 & 55 & 53 & 53 \\ \hline
C1000.9            & 68 & 61 & 63 & 63 & 63 \\ \hline
C2000.5            & 16* & 12 & 14 & 15 & 15 \\ \hline
C2000.9            & 80 & 72 & 72 & 71 & 69 \\ \hline
C4000.5            & 18* & 16 & 15 & 16 & 16 \\ \hline
DSJC500\_5         & 13 & 12 & 12 & 13 & 13 \\ \hline
DSJC1000\_5        & 15 & 14 & 14 & 14 & 14 \\ \hline
gen200\_p0.9\_44   & 44 & 39 & 40 & 44 & 44 \\ \hline
gen200\_p0.9\_55   & 55 & 55 & 55 & 55 & 55 \\ \hline
gen400\_p0.9\_55   & 55 & 51 & 51 & 51 & 49 \\ \hline
gen400\_p0.9\_65   & 65 & 49 & 65 & 52 & 65 \\ \hline
gen400\_p0.9\_75   & 75 & 75 & 75 & 75 & 75 \\ \hline
hamming10-4        & 40 & 38 & 40 & 40 & 40 \\ \hline
hamming8-4         & 16 & 16 & 16 & 16 & 16 \\ \hline
keller4            & 11 & 10 & 11 & 11 & 11 \\ \hline
keller5            & 27 & 23 & 25 & 26 & 27 \\ \hline
keller6            & 59 & 48 & 54 & 54 & 53 \\ \hline
MANN\_a27          & 126 & 126 & 125 & 125 & 125 \\ \hline
p\_hat300-1        & 8 & 7 & 8 & 8 & 8 \\ \hline
p\_hat300-2        & 25 & 25 & 25 & 25 & 25 \\ \hline
p\_hat300-3        & 36 & 36 & 36 & 36 & 36 \\ \hline
p\_hat700-1        & 11 & 9 & 10 & 9 & 11 \\ \hline
p\_hat700-2        & 44 & 44 & 44 & 44 & 44 \\ \hline
p\_hat700-3        & 62* & 61 & 62 & 62 & 62 \\ \hline
p\_hat1500-1       & 12 & 10 & 10 & 10 & 11 \\ \hline
p\_hat1500-2       & 65* & 65 & 65 & 65 & 65 \\ \hline
p\_hat1500-3       & 94* & 93 & 94 & 94 & 93 \\ \hline
Óptimos            & 35  & 11 & 15 & 16 & 18 \\ \hline
Aciertos           & 100\% & 92.6\% & 95.1\% & 94.6\% & 95.8\% \\ \hline
Media aciertos     & 100\% & 90.3\% & 93.3\% & 93.5\% & 95.3\% \\ \hline
\end{tabular}
\end{table}

\begin{table}[H]
\centering
\caption{Resultados}
\begin{tabular}{|l|l|l|l|l|l|}
\hline
Grafo              & Mejor conocido & ACO1 & ACO2 & Genético & Memético \\ \hline
brock200\_2        & 12 & 12 & 12 & 12 & 12 \\ \hline
brock200\_4        & 17 & 16 & 16 & 17 & 17 \\ \hline
brock400\_2        & 29 & 22 & 22 & 23 & 25 \\ \hline
brock400\_4        & 33 & 23 & 22 & 25 & 33 \\ \hline
brock800\_2        & 24 & 19 & 19 & 20 & 20 \\ \hline
brock800\_4        & 26 & 18 & 19 & 19 & 20 \\ \hline
C125.9             & 34* & 34 & 34 & 34 & 34 \\ \hline
C250.9             & 44* & 44 & 42 & 42 & 44 \\ \hline
C500.9             & 57 & 53 & 50 & 48 & 56 \\ \hline
C1000.9            & 68 & 55 & 54 & 54 & 65 \\ \hline
C2000.5            & 16* & 14 & 14 & 15 & 16 \\ \hline
C2000.9            & 80 & 60 & 60 & 60 & 73 \\ \hline
C4000.5            & 18* & 15 & 16 & 12 & 17 \\ \hline
DSJC500\_5         & 13 & 13 & 12 & 13 & 13 \\ \hline
DSJC1000\_5        & 15 & 13 & 13 & 13 & 15 \\ \hline
gen200\_p0.9\_44   & 44 & 41 & 39 & 40 & 44 \\ \hline
gen200\_p0.9\_55   & 55 & 55 & 55 & 55 & 55 \\ \hline
gen400\_p0.9\_55   & 55 & 49 & 46 & 46 & 55 \\ \hline
gen400\_p0.9\_65   & 65 & 57 & 48 & 47 & 65 \\ \hline
gen400\_p0.9\_75   & 75 & 75 & 65 & 69 & 75 \\ \hline
hamming10-4        & 40 & 33 & 32 & 34 & 40 \\ \hline
hamming8-4         & 16 & 16 & 16 & 16 & 16 \\ \hline
keller4            & 11 & 11 & 11 & 11 & 11 \\ \hline
keller5            & 27 & 24 & 24 & 27 & 27 \\ \hline
keller6            & 59 & 48 & 45 & 47 & 53 \\ \hline
MANN\_a27          & 126 & 125 & 124 & 125 & 125 \\ \hline
p\_hat300-1        & 8 & 8 & 8 & 8 & 8 \\ \hline
p\_hat300-2        & 25 & 25 & 25 & 25 & 25 \\ \hline
p\_hat300-3        & 36 & 36 & 35 & 36 & 36 \\ \hline
p\_hat700-1        & 11 & 9 & 10 & 11 & 11 \\ \hline
p\_hat700-2        & 44 & 44 & 43 & 44 & 44 \\ \hline
p\_hat700-3        & 62* & 61 & 61 & 60 & 62 \\ \hline
p\_hat1500-1       & 12 & 10 & 10 & 11 & 11 \\ \hline
p\_hat1500-2       & 65* & 44 & 54 & 61 & 65 \\ \hline
p\_hat1500-3       & 94* & 87 & 82 & 86 & 94 \\ \hline
Óptimos            & 35  & 12 & 7 & 13 & 25 \\ \hline
Aciertos           & 100\% & 89.6 \% & 87.4\% & 89.4\% & 97.6 \\ \hline
Media aciertos     & 100\% & 89.7 \% & 88.5\% & 90.5\% & 97.3 \\ \hline
\end{tabular}
\end{table}

\fi

% ********************************************************************
% Backmatter
%*******************************************************
\appendix
%\renewcommand{\thechapter}{\alph{chapter}}
\cleardoublepage
%\part{Appendix}
%%********************************************************************
% Bibliography
%*******************************************************
% work-around to have small caps also here in the headline
\manualmark
\markboth{\spacedlowsmallcaps{\bibname}}{\spacedlowsmallcaps{\bibname}} % work-around to have small caps also
%\phantomsection
\refstepcounter{dummy}
\addtocontents{toc}{\protect\vspace{\beforebibskip}} % to have the bib a bit from the rest in the toc
\addcontentsline{toc}{chapter}{\tocEntry{\bibname}}
\label{app:bibliography}
\printbibliography

%********************************************************************
% Other Stuff in the Back
%*******************************************************
\cleardoublepage%********************************************************************
% Bibliography
%*******************************************************
% work-around to have small caps also here in the headline
\manualmark
\markboth{\spacedlowsmallcaps{\bibname}}{\spacedlowsmallcaps{\bibname}} % work-around to have small caps also
%\phantomsection
\refstepcounter{dummy}
\addtocontents{toc}{\protect\vspace{\beforebibskip}} % to have the bib a bit from the rest in the toc
\addcontentsline{toc}{chapter}{\tocEntry{\bibname}}
\label{app:bibliography}
\printbibliography

%\cleardoublepage\include{FrontBackmatter/Declaration}
%\cleardoublepage\include{FrontBackmatter/Colophon}
% ********************************************************************
% Game Over: Restore, Restart, or Quit?
%*******************************************************
\end{document}
% ********************************************************************
